\documentclass[11pt,a4paper]{article}
\usepackage[utf8]{inputenc}
\usepackage[T1]{fontenc}
\usepackage[czech]{babel}
\usepackage[margin=2.2cm]{geometry}
\usepackage{amsmath,amssymb,amsthm}
\usepackage{parskip}
\usepackage{enumitem}
\setlist{nosep,leftmargin=1.4em}
\usepackage{booktabs}
\usepackage{array}
\usepackage{xcolor}
\definecolor{linkblue}{RGB}{20,60,150}
\definecolor{defbg}{RGB}{240,244,252}
\definecolor{defframe}{RGB}{100,130,190}
\definecolor{markcol}{RGB}{176,84,0}
\definecolor{markbg}{RGB}{253,246,236}
\usepackage[most]{tcolorbox}
\newtcolorbox{defbox}[1][]{breakable,colback=defbg,colframe=defframe,boxrule=0.6pt,arc=1.5mm,left=2mm,right=2mm,top=1mm,bottom=1mm,title={#1},fonttitle=\bfseries}
\newtcolorbox{poznbox}[1][]{breakable,colback=markbg,colframe=markcol,boxrule=0.6pt,arc=1.5mm,left=2mm,right=2mm,top=1mm,bottom=1mm,title={#1},fonttitle=\bfseries}
\newtheorem{veta}{Věta}[section]
\theoremstyle{definition}
\newtheorem{definice}{Definice}[section]
\usepackage[colorlinks=true,linkcolor=linkblue,urlcolor=linkblue]{hyperref}
\usepackage{algorithm}
\usepackage{algpseudocode}

\floatname{algorithm}{Algoritmus}
\renewcommand{\algorithmicrequire}{\textbf{Vstup:}}
\renewcommand{\algorithmicensure}{\textbf{Výstup:}}

% Značka pro látku, která není ve slidech EVA I / EVA II.
\newcommand{\mimo}{\textcolor{markcol}{\textsf{\footnotesize[$\blacklozenge$~nad rámec slidů]}}}
\newcommand{\mimoq}[1]{\textcolor{markcol}{\textsf{\footnotesize[$\blacklozenge$~nad rámec slidů: #1]}}}

\title{\textbf{Přírodou inspirované počítání}\\[2mm]
\large Učební text k~okruhu státní závěrečné zkoušky}
\author{SZZ Umělá inteligence, MFF UK}
\date{srpen 2026}

\begin{document}
\maketitle

\begin{center}
\begin{minipage}{0.92\textwidth}
\small
Text pokrývá okruh \emph{Přírodou inspirované počítání} magisterské SZZ oboru Umělá
inteligence. Kapitoly odpovídají \textbf{jedna ku jedné} větám oficiálního znění okruhu
(viz mapovací tabulka níže). Vychází ze slidů přednášek NAIL025 (Evoluční algoritmy~I)
a~NAIL086 (Evoluční algoritmy~II) a~je doplněn o~standardní látku z~učebnic
Eiben \& Smith: \emph{Introduction to Evolutionary Computing}, Mitchell:
\emph{An Introduction to Genetic Algorithms}, Michalewicz: \emph{Genetic Algorithms $+$
Data Structures $=$ Evolution Programs} a~Poli, Langdon \& McPhee: \emph{A Field Guide to
Genetic Programming}. Obsahuje definice, pseudokódy, odvození, číselné příklady
a~srovnávací tabulky; každá kapitola končí shrnutím pro ústní zkoušku.
\end{minipage}
\end{center}

\begin{center}
\begin{minipage}{0.92\textwidth}
\begin{poznbox}[Rozsah textu a~značení]
\small
Rozsah kapitol odpovídá váze látky: nejvíce prostoru dostávají genetické algoritmy
s~genetickým programováním, teorie schémat a~aplikace --- tedy to, čemu se přednášky
věnují nejpodrobněji.

Značkou \mimo{} je označena látka, která \textbf{není v~žádných dostupných slidech} ---
tj.\ ani v~EVA~I, ani v~EVA~II, ani ve volitelné \emph{Evoluční robotice} (Mráz) --- ale je
v~textu ponechána, buď proto, že ji \emph{jmenuje oficiální znění okruhu}, nebo proto, že
bez ní nedává okolní výklad smysl. Slouží jako nabídka ke~škrtnutí: co označeno není, je ve
slidech doloženo. U~látky, která je jen letmo zmíněna, značka uvádí i~rozsah zmínky.

Nejnápadnější případ je \textbf{otevřená evoluce}: okruh ji výslovně jmenuje, slidy o~ní
mají jedinou větu. Podobně mravenčí algoritmy (ACO) --- okruh mluví o~rojových algoritmech
v~množném čísle, slidy probírají jen PSO (rojová robotika v~Evoluční robotice je něco
jiného než rojová optimalizace). Naopak \emph{gramatická evoluce} a~\emph{interaktivní
evoluce} označeny nejsou: v~EVA sice nejsou, ale Evoluční robotika je probírá.
\end{poznbox}
\end{minipage}
\end{center}

\vspace{2mm}

\section*{Mapování okruhu na kapitoly}
\addcontentsline{toc}{section}{Mapování okruhu na kapitoly}

\begin{center}
\small
\begin{tabular}{@{}>{\raggedright\arraybackslash}p{8.6cm}>{\raggedright\arraybackslash}p{7.4cm}@{}}
\toprule
\textbf{Věta oficiálního znění okruhu} & \textbf{Kapitola v~tomto textu} \\
\midrule
Genetické algoritmy, genetické a~evoluční programování.
  & \ref{sec:ga}~Genetické algoritmy, genetické a~evoluční programování \\
Teorie schémat, pravděpodobnostní modely jednoduchého genetického algoritmu.
  & \ref{sec:schemata}~Teorie schémat a~pravděpodobnostní modely SGA \\
Evoluční strategie, diferenciální evoluce, koevoluce, otevřená evoluce.
  & \ref{sec:es}~Evoluční strategie, diferenciální evoluce, koevoluce, otevřená evoluce \\
Rojové optimalizační algoritmy.
  & \ref{sec:roje}~Rojové optimalizační algoritmy \\
Memetické algoritmy, hill climbing, simulované žíhání.
  & \ref{sec:lokalni}~Lokální prohledávání a~memetické algoritmy \\
Aplikace evolučních algoritmů (evoluce expertních systémů, neuroevoluce, řešení
kombinatorických úloh, vícekriteriální optimalizace).
  & \ref{sec:aplikace}~Aplikace evolučních algoritmů \\
\bottomrule
\end{tabular}
\end{center}

\vspace{2mm}
\tableofcontents
\newpage

%=====================================================================
\section{Genetické algoritmy, genetické a~evoluční programování}
\label{sec:ga}
%=====================================================================

Nejobsáhlejší věta okruhu; kapitola proto začíná obecným schématem evolučního algoritmu
(bez něj nedává zbytek textu smysl), pokračuje Hollandovým jednoduchým genetickým
algoritmem a~reprezentacemi jedince a~končí genetickým a~evolučním programováním, tedy
evolucí spustitelných struktur.

\subsection{Obecné schéma evolučního algoritmu}
\label{ssec:obecne}
\subsubsection{Zařazení oboru}

\emph{Přírodou inspirované počítání} (\emph{nature-inspired computing}) je souhrnné
označení pro metaheuristické algoritmy, které přebírají principy pozorované
v~živé (a~někdy i~neživé) přírodě. Tradiční členění je následující:

\begin{itemize}
\item \textbf{Výpočetní inteligence} (\emph{computational intelligence}, dříve
  \emph{soft computing}) zahrnuje neuronové sítě, fuzzy systémy a~evoluční
  výpočty. Stojí proti \uv{tvrdé}, logicky a~symbolicky orientované AI.
\item \textbf{Evoluční výpočty} (\emph{evolutionary computation}, EC) jsou
  nadmnožinou zahrnující i~metody, které nejsou evoluční v~úzkém smyslu
  (rojové algoritmy, memetické algoritmy, tabu prohledávání).
\item \textbf{Evoluční algoritmy} (\emph{evolutionary algorithms}, EA) jsou
  podmnožina EC využívající základní principy přírodní evoluce: populaci
  kandidátních řešení, stochastické operátory (křížení, mutace) a~selekci
  řízenou účelovou funkcí.
\end{itemize}

EA jsou \textbf{populační stochastické prohledávací algoritmy}. Používají se
tam, kde nemáme k~dispozici gradient nebo analytický tvar účelové funkce
(tzv.\ \emph{black-box optimalizace}), kde je prostor řešení diskrétní,
strukturovaný (stromy, grafy, permutace) nebo kde hledáme více vzájemně
nedominovaných řešení najednou.

\subsubsection{Biologická inspirace}

\begin{itemize}
\item \textbf{Darwin (1859, O~původu druhů).} Prostředí má omezené zdroje,
  reprodukce je klíčem k~přežití, lépe přizpůsobení jedinci mají větší šanci se
  rozmnožit. Úspěšné znaky fenotypu se reprodukují, modifikují a~rekombinují.
\item \textbf{Mendel (1866).} Gen jako základní jednotka dědičnosti; alely se
  předávají nezávisle. Realita je složitější: \emph{polygenie} (více genů
  ovlivňuje jeden znak), \emph{pleiotropie} (jeden gen ovlivňuje více znaků),
  mitochondriální DNA, epigenetika.
\item \textbf{Watson a~Crick (1953).} Struktura DNA; molekulárně-biologický
  pohled na ukládání a~dědění informace. Kodon (trojice nukleotidů) kóduje jednu
  z~aminokyselin, kód je redundantní.
\item \textbf{Genotyp $\to$ fenotyp.} Zobrazení je jednosměrné a~složité
  (transkripce DNA${\to}$RNA, translace RNA${\to}$protein). \emph{Lamarckismus}
  naopak předpokládá existenci inverzního zobrazení, tedy dědičnost získaných
  vlastností. V~biologii je odmítnut, v~memetických algoritmech se však jako
  technika používá (kap.~\ref{sec:lokalni}).
\end{itemize}

\begin{defbox}[Slovník: příroda vs.\ algoritmus]
\begin{tabular}{@{}ll@{}}
prostředí & optimalizační úloha \\
jedinec / chromozom & kandidátní řešení, bod prohledávaného prostoru \\
gen, alela & složka řešení a~její hodnota \\
genotyp & vnitřní reprezentace (kódování) jedince \\
fenotyp & \uv{vnějšek} jedince, na němž se hodnotí kvalita \\
fitness & míra kvality řešení (účelová funkce) \\
populace & multimnožina kandidátních řešení \\
generace & jedna iterace algoritmu \\
\end{tabular}
\end{defbox}

\subsubsection{Obecné schéma evolučního algoritmu}

\begin{defbox}[Evoluční algoritmus]
Evoluční algoritmus je iterativní stochastický proces, který udržuje populaci
$P(t)$ kandidátních řešení a~z~ní opakovaně vytváří populaci $P(t+1)$ pomocí
\emph{rodičovské selekce}, \emph{rekombinace}, \emph{mutace} a~\emph{environmentální
selekce}, dokud není splněna ukončovací podmínka.
\end{defbox}

\begin{algorithm}[!ht]
\caption{Obecné schéma evolučního algoritmu}
\begin{algorithmic}[1]
\Require reprezentace, fitness $f$, operátory, parametry
\State $t \gets 0$; inicializuj náhodně populaci $P(0)$
\State ohodnoť všechny jedince v~$P(0)$ funkcí $f$
\While{není splněna ukončovací podmínka}
  \State \textbf{rodičovská selekce:} vyber z~$P(t)$ množinu rodičů (pravděpodobnostně, dle fitness)
  \State \textbf{rekombinace:} z~dvojic (n-tic) rodičů vytvoř potomky s~pravděpodobností $p_C$
  \State \textbf{mutace:} každého potomka náhodně poruš s~pravděpodobností $p_M$
  \State ohodnoť potomky $P'(t+1)$
  \State \textbf{environmentální selekce:} z~$P(t) \cup P'(t+1)$ (nebo jen z~$P'(t+1)$) vyber novou populaci $P(t+1)$
  \State $t \gets t+1$
\EndWhile
\Ensure nejlepší nalezený jedinec
\end{algorithmic}
\end{algorithm}

Tento cyklus je společný všem EA; jednotlivé varianty (GA, ES, EP, GP, DE)
se liší volbou reprezentace, důrazem na jednotlivé operátory a~typem selekce.
Klíčovou roli hrají dvě protichůdné síly:

\begin{itemize}
\item \textbf{Variace} (křížení a~mutace) vytváří novou rozmanitost --- zajišťuje
  \emph{explorace} (\emph{exploration}) prostoru.
\item \textbf{Selekce} rozmanitost odebírá a~směřuje hledání ke slibným
  oblastem --- zajišťuje \emph{exploataci} (\emph{exploitation}).
\end{itemize}

Rovnováha mezi nimi je centrální problém návrhu EA. Samotná explorace degeneruje
na náhodnou procházku; samotná exploatace vede k~uvíznutí v~lokálním extrému
a~\emph{předčasné konvergenci} (\emph{premature convergence}).

\subsubsection{Složky evolučního algoritmu}

\paragraph{Reprezentace.} Volba kódování je nejdůležitější návrhové rozhodnutí.
Klasické možnosti: binární řetězec, vektor celých čísel, vektor reálných čísel,
permutace, strom, graf, matice, konečný automat, neuronová síť. Reprezentace
určuje, jaké operátory dávají smysl.

\paragraph{Fitness funkce.} Zobrazení $f:$ prostor fenotypů $\to \mathbb{R}$,
které měří kvalitu. Pro minimalizační úlohy se obvykle převádí na maximalizaci
(např.\ $f' = C_{\max} - f$ nebo $f' = 1/(1+f)$). Fitness může být
\emph{deterministická}, \emph{zašuměná} (simulace), \emph{dynamická}
(mění se v~čase, odst.~\ref{sec:koevoluce}) nebo \emph{kontextová}
(závisí na ostatních jedincích --- koevoluce).

\paragraph{Populace.} Multimnožina jedinců pevné (obvykle) velikosti $n$.
Populace je nositelem \emph{diverzity}; ta se měří např.\ průměrnou Hammingovou
vzdáleností, entropií alel na pozicích nebo rozptylem fitness.

\paragraph{Inicializace.} Obvykle náhodná uniformní. Alternativy: \emph{seeding}
(vložení heuristických nebo expertních řešení --- riziko předčasné konvergence),
konstrukční heuristiky (např.\ nejbližší soused pro TSP), latinské
hyperkrychle pro reálné vektory.

\paragraph{Ukončovací podmínka.} Maximální počet generací či vyhodnocení
fitness, dosažení cílové kvality, stagnace nejlepší fitness po $k$ generací,
ztráta diverzity populace.

\subsubsection{Selekční mechanismy}\label{ssec:selekce}

\paragraph{Ruletová (fitness-proporcionální) selekce}

\begin{defbox}[Ruletová selekce (\emph{roulette wheel selection})]
Jedinec $x_i$ je vybrán s~pravděpodobností
\[
p_i \;=\; \frac{f(x_i)}{\sum_{j=1}^{n} f(x_j)} \;=\; \frac{f(x_i)}{n\,\bar{f}},
\]
kde $\bar f$ je průměrná fitness populace. Očekávaný počet výběrů jedince při
$n$ tazích je $n p_i = f(x_i)/\bar f$. Jde o~výběr \emph{s~vracením} --- jeden
jedinec může být vybrán vícekrát.
\end{defbox}

\textbf{Příklad (klasický Goldbergův).}\label{ex:goldberg} Populace čtyř
5-bitových řetězců, fitness $f(x) = x^2$ (dekódováno jako celé číslo):

\begin{center}
\begin{tabular}{@{}llrrrr@{}}
\toprule
$i$ & řetězec & $x$ & $f=x^2$ & $p_i = f/\sum f$ & oček.\ počet $f/\bar f$ \\
\midrule
1 & 01101 & 13 & 169 & 0{,}144 & 0{,}58 \\
2 & 11000 & 24 & 576 & 0{,}492 & 1{,}97 \\
3 & 01000 &  8 &  64 & 0{,}055 & 0{,}22 \\
4 & 10011 & 19 & 361 & 0{,}309 & 1{,}23 \\
\midrule
  & \multicolumn{2}{l}{součet} & 1170 & 1{,}000 & 4{,}00 \\
  & \multicolumn{2}{l}{průměr $\bar f$} & 292{,}5 & & \\
\bottomrule
\end{tabular}
\end{center}

Ruleta rozdělí obvod kola na výseče o~velikostech $p_i$ a~roztočí se $n$-krát.
Jedinec 2 obsadí zhruba dvě místa v~mezipopulaci, jedinec 3 pravděpodobně
vypadne.

\textbf{Nevýhody ruletové selekce:}
\begin{itemize}
\item \emph{Předčasná dominance} \uv{supermana}: jedinec s~extrémní fitness
  ovládne populaci a~diverzita zmizí.
\item \emph{Ztráta selekčního tlaku} na konci běhu: když jsou všechny fitness
  podobné, selekce se blíží náhodnému výběru.
\item Vyžaduje kladné fitness a~je citlivá na afinní transformace
  ($f$ a~$f + c$ dávají jiné pravděpodobnosti). Řeší se \emph{škálováním}
  (odst.~\ref{ssec:skalovani}).
\item Vysoký rozptyl výběru --- skutečné počty se mohou od očekávaných značně lišit.
\end{itemize}

\paragraph{Stochastické univerzální vzorkování (SUS)}

\emph{Stochastic universal sampling} používá jediné roztočení kola s~$n$
rovnoměrně rozmístěnými ukazateli vzdálenými o~$1/n$. Očekávané počty výběrů
zůstávají stejné jako u~rulety, ale rozptyl je minimální: jedinec s~očekávaným
počtem $e_i$ je vybrán $\lfloor e_i \rfloor$ nebo $\lceil e_i \rceil$-krát.
Je to tedy \uv{spravedlivější ruleta}.

\paragraph{Turnajová selekce}

\begin{defbox}[$k$-turnajová selekce (\emph{tournament selection})]
Vyber náhodně (uniformně, s~vracením) $k$ jedinců a~vrať toho nejlepšího.
Parametr $k$ (typicky 2, 3, 5) přímo řídí selekční tlak: pravděpodobnost, že
bude vybrán $i$-tý nejlepší jedinec z~$n$ (\emph{zde} $i=1$ značí
\textbf{nejlepšího}), je
$\big((n-i+1)^k - (n-i)^k\big)/n^k$ --- je to pravděpodobnost, že všech $k$
losovaných padne mezi $i$-tého a~horší, mínus pravděpodobnost, že všech $k$
padne ostře pod $i$-tého. Deterministický turnaj lze změkčit:
lepší vyhraje s~pravděpodobností $p < 1$.
\end{defbox}

Výhody: nevyžaduje globální informaci (žádné sčítání fitness), je invariantní
vůči monotónním transformacím fitness, snadno paralelizovatelná, funguje i~tam,
kde fitness není explicitně dána a~porovnání se získává \emph{souhrou}
(hraná partie, simulace, koevoluce).

\paragraph{Pořadová selekce}

\emph{Rank-based selection} nahradí fitness pořadím jedince. Lineární pořadová
selekce přiřadí jedinci s~pořadím $i$ (\emph{pozor, zde opačně než u~turnaje:}
nejhorší $i=1$, nejlepší $i=n$)
pravděpodobnost
\[
p_i \;=\; \frac{1}{n}\Big(s^{-} + (s^{+}-s^{-})\,\frac{i-1}{n-1}\Big),
\qquad 1 \le s^{+} \le 2,\quad s^{-} = 2 - s^{+},
\]
kde $s^{+}$ je očekávaný počet kopií nejlepšího jedince. Exponenciální varianta
používá $p_i \propto (1-c)c^{\,n-i}$. Pořadová selekce udržuje konstantní
selekční tlak během celého běhu a~je odolná vůči \uv{supermanům}
i~vůči změně měřítka fitness.

\paragraph{Boltzmannova selekce}

Pravděpodobnost výběru $p_i \propto e^{f(x_i)/T}$ s~\emph{teplotou} $T$, která
klesá v~čase: na začátku (vysoké $T$) je selekce téměř uniformní (explorace),
na konci (nízké $T$) téměř deterministická (exploatace). Analogicky
\emph{Boltzmannův turnaj}: horší jedinec vyhraje s~pravděpodobností závislou
exponenciálně na poměru rozdílu fitness a~teploty. Přímá souvislost se
simulovaným žíháním (kap.~\ref{sec:lokalni}).

\paragraph{Environmentální selekce, elitismus, modely populace}

\begin{itemize}
\item \textbf{Generační model} (\emph{generational}): celá populace je nahrazena
  potomky, $|P'| = |P| = n$.
\item \textbf{Model s~ustáleným stavem} (\emph{steady-state}): v~každém kroku
  vznikne jen $\lambda$ (často 1--2) potomků, kteří nahradí vybrané jedince
  (nejhoršího, náhodného, nejpodobnějšího --- \emph{crowding}). Parametr
  \emph{generation gap} $G = \lambda/n$ udává podíl nahrazované populace.
\item \textbf{Elitismus} (\emph{elitism}): $k$ nejlepších jedinců se
  bezpodmínečně kopíruje do další generace. Zaručuje monotonii nejlepší
  nalezené fitness a~je nutnou podmínkou pro řadu konvergenčních vět;
  příliš silný elitismus však snižuje diverzitu.
\item \textbf{Deterministická selekce ES}: $(\mu+\lambda)$ vybírá $\mu$
  nejlepších z~rodičů i~potomků (implicitní elitismus),
  $(\mu,\lambda)$ jen z~potomků (kap.~\ref{sec:es}).
\end{itemize}

\begin{defbox}[Selekční tlak a~doba převzetí]
\emph{Selekční tlak} (\emph{selection pressure}) je míra zvýhodnění lepších
jedinců. Kvantifikuje se \emph{dobou převzetí} (\emph{takeover time}) $\tau^*$ ---
počtem generací, za který jediná kopie nejlepšího jedince (bez variace) zaplní
celou populaci. Pro turnaj velikosti $k$ platí zhruba
$\tau^* \approx \ln n/\ln k$, pro fitness-proporcionální selekci je doba převzetí
delší a~závisí na rozložení fitness.
\end{defbox}

\subsubsection{Teoretické meze: No Free Lunch}

\begin{veta}[No Free Lunch, Wolpert \& Macready 1995--1997]
Zprůměrováno přes \emph{všechny} možné účelové funkce na konečné doméně mají
libovolné dva deterministické neopakující se prohledávací algoritmy stejný
očekávaný výkon. Neexistuje tedy univerzálně nejlepší optimalizační algoritmus.
\end{veta}

Praktický důsledek: vyplácí se vytvářet \textbf{doménově specifické varianty} EA ---
vhodnou reprezentaci a~operátory zohledňující strukturu úlohy. Věta neříká, že všechny
algoritmy jsou stejně dobré na \emph{reálných} úlohách; reálné funkce tvoří velmi malou
a~strukturovanou podmnožinu všech funkcí.

\subsubsection{Historické větve evolučních algoritmů}

\begin{center}
\small
\begin{tabular}{@{}>{\raggedright\arraybackslash}p{2.3cm}
                  >{\raggedright\arraybackslash}p{3.1cm}
                  >{\raggedright\arraybackslash}p{3.1cm}
                  >{\raggedright\arraybackslash}p{3.1cm}
                  >{\raggedright\arraybackslash}p{3.1cm}@{}}
\toprule
 & \textbf{GA} & \textbf{ES} & \textbf{EP} & \textbf{GP} \\
\midrule
autor, rok & Holland, 1975 & Rechenberg, Schwefel, 1964 &
  L.~Fogel, Owens, Walsh, 1965 & Koza, 1989--92 \\[2pt]
reprezentace & binární řetězec & vektor $\mathbb{R}^n$ + strategické parametry &
  konečný automat (genotyp = fenotyp) & syntaktický strom (LISP) \\[2pt]
hlavní operátor & křížení & mutace & mutace & křížení podstromů \\[2pt]
křížení & 1-bodové & rekombinace více rodičů (lokální/globální) &
  \emph{není} & výměna podstromů \\[2pt]
mutace & bit-flip $p_M$ & gaussovská, samoadaptivní $\sigma$ &
  \uv{chytré} mutace automatu & více typů \\[2pt]
rodičovská selekce & ruleta & náhodná & každý jednou & turnaj \\[2pt]
env.\ selekce & generační náhrada &
  deterministická $(\mu,\lambda)$, $(\mu+\lambda)$ &
  turnaj (rodiče $+$ potomci) & generační \\[2pt]
teorie & schémata & konvergence, 1/5 pravidlo & --- & bloat, schémata \\
\bottomrule
\end{tabular}
\end{center}

Dnes se hranice mezi větvemi stírají (např.\ EP a~ES prakticky splynuly)
a~mluvíme jednotně o~evolučních algoritmech.

\subsection{Genetické algoritmy}
\label{ssec:sga}

\subsubsection{Jednoduchý genetický algoritmus (SGA)}

Původní Hollandův návrh z~roku 1975 se dnes označuje jako \emph{simple genetic
algorithm} (SGA). Používá minimální sadu operátorů, nejjednodušší kódování
a~byl navržen tak, aby umožňoval teoretickou analýzu (teorie schémat,
kap.~\ref{sec:schemata}).

\begin{defbox}[SGA]
Populace $n$ binárních řetězců délky $m$; fitness-proporcionální (ruletová)
selekce; jednobodové křížení s~pravděpodobností $p_C$; bitová mutace
s~pravděpodobností $p_M$ na každý bit; generační náhrada populace
($P(t+1)$ zcela nahradí $P(t)$).
\end{defbox}

\begin{algorithm}[!ht]
\caption{Jednoduchý genetický algoritmus (SGA)}
\begin{algorithmic}[1]
\State $t \gets 0$; vygeneruj náhodně $P(0)$ = $n$ řetězců délky $m$
\While{není splněna ukončovací podmínka}
  \State spočítej $f(x)$ pro každé $x \in P(t)$
  \State $P(t+1) \gets \emptyset$
  \For{$n/2$ krát}
    \State vyber ruletou dvojici rodičů $x, y$ z~$P(t)$
    \State s~pravděpodobností $p_C$ zkřiž $x, y$ jednobodovým křížením
    \State každý bit $x$ i~$y$ převrať s~pravděpodobností $p_M$
    \State vlož $x, y$ do $P(t+1)$
  \EndFor
  \State $t \gets t+1$
\EndWhile
\end{algorithmic}
\end{algorithm}

Typické hodnoty parametrů: $n \in [50, 500]$, $p_C \in [0{,}6; 0{,}9]$,
$p_M \approx 1/m$ (tj.\ v~průměru se změní jeden bit na jedince).

\subsubsection{Binární reprezentace}

\paragraph{Kódování celých a~reálných čísel.} Reálné číslo $x \in [a,b]$
kódované na $k$ bitů jako celé číslo $z \in \{0,\dots,2^k-1\}$ dekódujeme
\[
x \;=\; a \;+\; z \cdot \frac{b-a}{2^k - 1}.
\]
Přesnost je $(b-a)/(2^k-1)$. Výhoda: explicitní kontrola nad přesností
a~komprese prohledávaného prostoru. Nevýhoda: \emph{Hammingovy útesy}
(\emph{Hamming cliffs}) --- sousední hodnoty se mohou lišit ve všech bitech
(0111 $\to$ 1000), takže malá změna fenotypu vyžaduje velkou změnu genotypu.

\paragraph{Grayův kód.} Řeší Hammingovy útesy: sousední celá čísla se
v~Grayově kódu liší právě v~jednom bitu. Převod z~binárního $b$ na Grayův $g$:
$g_1 = b_1$, $g_i = b_{i-1} \oplus b_i$. Zpět: $b_1 = g_1$,
$b_i = b_{i-1} \oplus g_i$.

\paragraph{Hollandův argument pro binární abecedu.} Binární řetězce délky 100
kódují zhruba stejnou informaci jako desítkové délky 30, ale obsahují více
schémat ($3^{100}$ oproti $11^{30}$), takže GA \uv{zpracovává} více informací.
Dnes víme, že schémata nejsou tak zásadní, jak se Holland domníval, a~že
rozhodující je \textbf{přirozenost kódování pro danou úlohu}. Pro reálnou
optimalizaci se používají reálné vektory, pro TSP permutace atd.

\subsubsection{Operátory křížení}

\begin{defbox}[Křížení (\emph{crossover, recombination})]
Operátor, který ze dvou (nebo více) rodičů vytvoří potomka kombinací jejich
genetického materiálu. V~GA je hlavním operátorem; vyjadřuje naději, že
rekombinací dobrých částí (\emph{stavebních bloků}) vzniknou ještě lepší
jedinci. Aplikuje se s~pravděpodobností $p_C$, jinak potomci = kopie rodičů.
\end{defbox}

\paragraph{Jednobodové křížení.} Zvol náhodně bod řezu $k \in \{1,\dots,m-1\}$
a~vyměň koncové části.

\textbf{Příklad.} Rodiče $P_1 = 1011|0110$, $P_2 = 0010|1101$, bod řezu $k=4$:
\[
O_1 = \mathbf{1011}\,1101, \qquad O_2 = \mathbf{0010}\,0110.
\]

\paragraph{Dvoubodové křížení.} Zvol dva body $k_1 < k_2$ a~vyměň prostřední
úsek. Výhoda oproti jednobodovému: řetězec se chová jako kruh, takže nejsou
znevýhodněny geny na okrajích (u~jednobodového křížení je pravděpodobnost
rozbití úseku úměrná jeho \emph{definující délce}, viz kap.~\ref{sec:schemata}).

\textbf{Příklad.} $P_1 = 10|1101|10$, $P_2 = 00|1011|01$ $\Rightarrow$
$O_1 = 10\,1011\,10$, $O_2 = 00\,1101\,01$.

\paragraph{Uniformní křížení.} Pro každou pozici $j$ se hodí mince: s~pravděpodobností
$0{,}5$ pochází gen od prvního rodiče, jinak od druhého (druhý potomek je
komplementární). Nezvýhodňuje sousedící pozice, ale silně narušuje schémata
s~velkou definující délkou --- má vysokou \emph{disruptivitu}.

\textbf{Příklad.} $P_1 = 11111111$, $P_2 = 00000000$, maska $10110010$:
$O_1 = 10110010$, $O_2 = 01001101$.

\begin{center}
\begin{tabular}{@{}lll@{}}
\toprule
\textbf{operátor} & \textbf{p-st rozbití schématu délky $d$} & \textbf{poznámka} \\
\midrule
jednobodové & $d/(m-1)$ & poziční bias, chrání krátká schémata \\
dvoubodové & $\dfrac{2d(m-d)}{m(m-1)}$ & řetězec jako kruh, odstraňuje bias konců \\
$k$-bodové & roste s~$k$ & \\
uniformní & $1 - (1/2)^{o(S)-1}$ & distribuční bias, největší disruptivita \\
\bottomrule
\end{tabular}
\end{center}

\textbf{Poznámka k~dvoubodovému křížení.} Chápeme-li řetězec jako kruh, volíme
dva ze $m$ řezů; schéma je rozbito, padne-li \emph{právě jeden} řez dovnitř
jeho definujícího úseku o~$d$ mezerách, tj.\ s~pravděpodobností
$d(m-d)/\binom{m}{2} = 2d(m-d)/(m(m-1))$. Pro schéma sahající od prvního
k~poslednímu bitu ($d = m-1$) klesne rozbití z~jistoty (jednobodové:
$d/(m-1) = 1$) na pouhé $2/m$ --- to je přesně ono \uv{odstranění okrajového
biasu}. Pro krátká schémata je naopak dvoubodové křížení asi dvakrát
destruktivnější než jednobodové.

\subsubsection{Mutace a~inverze}

\paragraph{Bitová mutace.} Každý bit se s~pravděpodobností $p_M$ nezávisle
překlopí. V~SGA je mutace \emph{sekundárním} operátorem --- působí jako pojistka
proti uvíznutí v~lokálním extrému a~proti trvalé ztrátě alely
z~populace. (Naopak v~EP a~raných ES je mutace jediným zdrojem variability.)
Doporučení: $p_M \approx 1/m$.

\paragraph{Inverze.} Původní Hollandův návrh obsahoval ještě operátor
\emph{inverze}: obrátí pořadí části řetězce, ale \emph{zachová význam bitů}
(gen si nese svůj identifikátor pozice). Cílem bylo umožnit evoluci samotného
uspořádání genů tak, aby spolupracující geny ležely blízko sebe a~nebyly
rozbíjeny křížením. Technicky je to náročné a~přínos se neprokázal;
u~permutační reprezentace se však inverze používá jako velmi účinná mutace
(odst.~\ref{sec:reprez}).

\subsubsection{Škálování fitness}\label{ssec:skalovani}

Ruletová selekce je citlivá na absolutní hodnoty fitness. Na začátku běhu bývá
rozptyl fitness velký (několik \uv{supermanů} ovládne populaci), na konci malý
(selekce ztrácí tlak). Řešením je transformace $f \mapsto f'$:

\begin{itemize}
\item \textbf{Lineární škálování:} $f' = a f + b$, kde $a, b$ se volí tak, aby
  $\bar{f'} = \bar{f}$ a~$f'_{\max} = C_{\mathrm{mult}} \cdot \bar{f}$,
  typicky $C_{\mathrm{mult}} \in [1{,}2; 2]$. Nutná pojistka proti záporným
  hodnotám u~nejhorších jedinců.
\item \textbf{Sigma truncation:} $f' = \max\big(0,\; f - (\bar f - c\,\sigma_f)\big)$,
  kde $\sigma_f$ je směrodatná odchylka fitness a~$c \in [1;3]$. Odstraní
  odlehlé nízké hodnoty a~automaticky se přizpůsobuje rozptylu.
\item \textbf{Mocninné škálování:} $f' = f^{\,k}$, kde $k$ se v~průběhu běhu
  zvyšuje (tlak roste).
\item \textbf{Pořadové škálování:} nahradí fitness pořadím (viz
  odst.~\ref{ssec:selekce}) --- nejrobustnější, dnes nejobvyklejší.
\item \textbf{Boltzmannovo škálování:} $f' = e^{f/T}$ s~klesající teplotou.
\end{itemize}

\textbf{Příklad.} Populace s~fitness $(169, 576, 64, 361)$, $\bar f = 292{,}5$,
$f_{\max} = 576$. Lineární škálování s~$C_{\mathrm{mult}} = 1{,}5$:
chceme $f'_{\max} = 438{,}75$ a~$\bar{f'} = 292{,}5$. Z~podmínek
$a f_{\max} + b = 438{,}75$ a~$a \bar f + b = 292{,}5$ dostáváme
$a = 146{,}25/283{,}5 = 0{,}516$, $b = 292{,}5 - 0{,}516 \cdot 292{,}5 = 141{,}6$.
Nové fitness: $(228{,}8;\, 438{,}8;\, 174{,}6;\, 327{,}9)$; podíl nejlepšího
klesl z~$49{,}2\,\%$ na $37{,}5\,\%$ --- selekční tlak je zmírněn a~nejhorší
jedinec neztrácí šanci.

\subsection{Reprezentace jedince a~genetické operátory}
\label{sec:reprez}

Reprezentace určuje, které operátory mají smysl, a~tím i~to, jak vypadá
\emph{fitness krajina} (\emph{fitness landscape}) --- tedy jak snadné je
prohledávání. Toto je nejdůležitější praktické rozhodnutí při návrhu EA.

\subsubsection{Celočíselná reprezentace}

Jedinec je vektor $(z_1,\dots,z_m) \in \prod_j D_j$, kde $D_j$ jsou konečné
domény. Rozlišujeme:
\begin{itemize}
\item \textbf{Kardinální (ordinální) atributy} --- na hodnotách existuje
  smysluplné uspořádání a~metrika (např.\ počet strojů, věk).
\item \textbf{Nominální atributy} --- hodnoty jsou pouhé kódy (barva, typ
  komponenty). Zde nemá smysl mluvit o~\uv{blízkých} hodnotách.
\end{itemize}

\paragraph{Mutace.}
\begin{itemize}
\item \emph{Nezaujatá} (\emph{unbiased}, \emph{random resetting}): nová hodnota
  je náhodná z~celé domény. Jediná korektní volba pro nominální atributy.
\item \emph{Zaujatá} (\emph{biased}, \emph{creep mutation}): nová hodnota je
  původní posunutá o~malý náhodný krok (např.\ z~diskretizovaného normálního
  rozdělení). Vhodná jen pro ordinální atributy.
\end{itemize}

\paragraph{Křížení.} Jednobodové, vícebodové, uniformní --- stejně jako
u~binární reprezentace, jen se místo bitů kopírují celá čísla.

\subsubsection{Reálná reprezentace}

Jedinec je $\vec{x} = (x_1,\dots,x_n) \in \mathbb{R}^n$. Historicky se reálná
čísla kódovala do bitových řetězců; dnes se pracuje přímo s~reálnými hodnotami
a~operátory se tomu přizpůsobují.

\paragraph{Mutace.}
\begin{itemize}
\item \textbf{Gaussovská (normální) mutace:}
  $x_i' = x_i + \sigma_i \cdot N(0,1)$. Parametr $\sigma$ (šířka kroku) je
  klíčový; může být pevný, deterministicky klesající
  (např.\ $\sigma(t) = 1 - 0{,}9\,t/T$), adaptivní (pravidlo 1/5) nebo
  samoadaptivní (součást genomu, viz kap.~\ref{sec:es}).
\item \textbf{Cauchyho mutace:} těžší chvosty $\Rightarrow$ častější velké skoky
  $\Rightarrow$ lepší únik z~lokálních extrémů (používá se ve \emph{fast EP}).
\item \textbf{Uniformní mutace:} $x_i'$ náhodně z~$[a_i, b_i]$ --- nezaujatá.
\item Ošetření mezí: oříznutí (\emph{clipping}), odraz (\emph{reflection}),
  periodické zabalení (\emph{wrapping}) nebo převzorkování.
\end{itemize}

\paragraph{Křížení.} Rozlišujeme dva typy:
\begin{itemize}
\item \textbf{Strukturální (diskrétní):} 1-bodové, uniformní --- hodnoty se jen
  přeskupují mezi rodiči; potomek je vrcholem hyperkvádru daného rodiči.
\item \textbf{Aritmetické:} hodnoty se kombinují.
\end{itemize}

\begin{defbox}[Aritmetické křížení (\emph{arithmetic crossover})]
Pro rodiče $\vec x, \vec y$ a~$\alpha \in (0,1)$:
\[
\vec z \;=\; \alpha \vec x + (1-\alpha)\vec y \qquad\text{(konvexní kombinace)}.
\]
Varianty: \emph{jednoduché} (kříží se všechny složky --- nejběžnější),
\emph{jednoduché aritmetické} (jen jedna náhodně vybraná složka),
\emph{celé/single} (kombinace s~jednobodovým křížením: za bodem řezu se
aritmeticky kombinuje). Pro $\alpha = 1/2$ jde o~prostý průměr rodičů.
\end{defbox}

\textbf{Příklad.} $\vec x = (2{,}0;\, -1{,}0;\, 4{,}0)$,
$\vec y = (0{,}0;\, 3{,}0;\, 1{,}0)$, $\alpha = 0{,}25$:
\[
\vec z = 0{,}25\,\vec x + 0{,}75\,\vec y = (0{,}5;\; 2{,}0;\; 1{,}75).
\]

Nevýhoda čistě aritmetického křížení: potomci leží uvnitř konvexního obalu
rodičů, takže rozptyl populace nutně klesá (\uv{smršťování}). Proto se používají
rozšířené varianty:
\begin{itemize}
\item \textbf{BLX-$\alpha$} (\emph{blend crossover}): $z_i$ se losuje uniformně
  z~intervalu $[\min - \alpha I,\; \max + \alpha I]$, kde
  $I = |x_i - y_i|$; typicky $\alpha = 0{,}5$. Umožňuje potomkům ležet
  i~vně rodičů.
\item \textbf{SBX} (\emph{simulated binary crossover}): napodobuje chování
  jednobodového binárního křížení v~reálném prostoru. Potomci vzniknou
  \uv{roztažením} dvojice rodičů kolem jejich středu faktorem $\beta$:
  \[
  z_{1,i} = \tfrac12\big[(1+\beta)x_i + (1-\beta)y_i\big],\qquad
  z_{2,i} = \tfrac12\big[(1-\beta)x_i + (1+\beta)y_i\big],
  \]
  takže střed rodičů zůstává zachován a~$|z_1 - z_2| = \beta\,|x - y|$.
  Faktor se losuje z~$u \sim U(0,1)$: $\beta = (2u)^{1/(\eta+1)}$ pro
  $u \le 1/2$, jinak $\beta = \big(2(1-u)\big)^{-1/(\eta+1)}$.
  Intuice: $\beta = 1$ vrátí rodiče, $\beta < 1$ dá potomky mezi ně,
  $\beta > 1$ vně; velké $\eta$ tlačí $\beta$ k~jedničce (potomci blízko
  rodičů), malé $\eta$ dovolí velké skoky. Standardní operátor v~NSGA-II.
\end{itemize}

\subsubsection{Permutační reprezentace}

Jedinec je permutace $\pi$ na $\{1,\dots,n\}$. Typické úlohy: obchodní cestující
(TSP), rozvrhování, přiřazovací úlohy, uspořádání operací. Klasické operátory
zde selhávají --- výsledek jednobodového křížení dvou permutací obecně není
permutace. Je proto třeba operátory, které \emph{zachovávají přípustnost}
a~zároveň dědí něco smysluplného. Které informace mají být zděděny, závisí
na úloze:

\begin{center}
\begin{tabular}{@{}lll@{}}
\toprule
\textbf{operátor} & \textbf{co zachovává} & \textbf{typické použití} \\
\midrule
PMX & absolutní pozice měst & úlohy s~významnou pozicí \\
OX & relativní pořadí & TSP (cyklické cesty) \\
CX & absolutní pozice (cykly) & úlohy s~významnou pozicí \\
ER & hrany (sousednosti) & TSP --- nejlepší kvalita \\
\bottomrule
\end{tabular}
\end{center}

\paragraph{PMX --- částečně mapované křížení}

\begin{defbox}[PMX (\emph{partially mapped crossover}, Goldberg)]
Dvoubodový operátor. (1)~Vyber úsek mezi dvěma body řezu a~vyměň jej mezi
rodiči. (2)~Z~vyměněných úseků odvoď \emph{mapování} odpovídajících si hodnot.
(3)~Zbylé pozice zkopíruj z~původního rodiče; kde by vznikl konflikt (hodnota už
je v~potomkovi), použij mapování (případně opakovaně).
\end{defbox}

\textbf{Příklad.} $P_1 = (123|4567|89)$, $P_2 = (452|1876|93)$.
\begin{enumerate}
\item Výměna úseků: $O_1 = (\ldots|1876|\ldots)$, $O_2 = (\ldots|4567|\ldots)$.
\item Mapování z~úseků: $1{\leftrightarrow}4$, $8{\leftrightarrow}5$,
  $7{\leftrightarrow}6$, $6{\leftrightarrow}7$.
\item Doplnění bezkonfliktních pozic z~$P_1$ do $O_1$:
  $O_1 = (.23|1876|.9)$ (hodnoty 1 a~8 by kolidovaly).
\item Konflikty vyřešíme mapováním: $1 \to 4$, $8 \to 5$, tedy
  $O_1 = (423|1876|59)$. Analogicky $O_2 = (182|4567|93)$.
\end{enumerate}

\paragraph{OX --- křížení zachovávající pořadí}

\begin{defbox}[OX (\emph{order crossover}, Davis)]
Dvoubodový operátor zachovávající \emph{relativní pořadí}. (1)~Zkopíruj úsek
mezi body řezu z~prvního rodiče. (2)~Ze druhého rodiče čti hodnoty cyklicky
počínaje pozicí za druhým bodem řezu a~vynech ty, které už v~potomkovi jsou.
(3)~Zbylé hodnoty vkládej do volných pozic potomka, opět cyklicky od pozice
za druhým bodem řezu.
\end{defbox}

\textbf{Příklad.} $P_1 = (123|4567|89)$, $P_2 = (452|1876|93)$.
\begin{enumerate}
\item $O_1$ převezme úsek z~$P_1$: $O_1 = ({\cdot}{\cdot}{\cdot}|4567|{\cdot}{\cdot})$.
\item Z~$P_2$ čteme cyklicky od pozice 8: $9,3,4,5,2,1,8,7,6$.
\item Vyškrtneme hodnoty už obsažené ($4,5,6,7$): zbývá $9,3,2,1,8$.
\item Doplňujeme od pozice 8 cyklicky (pozice $8,9,1,2,3$):
  $O_1 = (218|4567|93)$.
\item Symetricky (úsek z~$P_2$, pořadí z~$P_1$): $O_2 = (345|1876|92)$.
\end{enumerate}
Princip: vnitřní úsek se převezme beze změny, zbytek se doplní v~relativním
pořadí druhého rodiče s~vynecháním duplicit. Pro TSP je OX vhodnější než PMX,
protože u~okružní jízdy nezáleží na absolutní pozici města, ale na pořadí.

\paragraph{CX --- cyklické křížení}

\begin{defbox}[CX (\emph{cyclic crossover}, Oliver)]
Zachovává \emph{absolutní pozici}: každý gen potomka pochází ze stejné pozice
u~některého z~rodičů. Sestrojí se cykly: začni na pozici 1, vezmi hodnotu
z~$P_1$; podívej se, jaká hodnota je na téže pozici v~$P_2$, najdi ji v~$P_1$
a~pokračuj, dokud se cyklus neuzavře. Sudé cykly ber z~jednoho rodiče, liché
z~druhého.
\end{defbox}

\textbf{Příklad.} $P_1 = (123456789)$, $P_2 = (412876935)$.
Začneme hodnotou 1 z~$P_1$ na pozici 1. Pod ní je v~$P_2$ hodnota 4, tu najdeme
v~$P_1$ na pozici 4 $\Rightarrow$ přidáme; pod ní je 8, ta je v~$P_1$ na pozici
8 $\Rightarrow$ přidáme; pod ní je 3 (pozice 3), pod ní 2 (pozice 2), pod 2 je
1 --- cyklus se uzavřel. Máme
$O_1 = (1\,2\,3\,4\,\_\,\_\,\_\,8\,\_)$; zbytek doplníme z~$P_2$:
$O_1 = (123476985)$, analogicky $O_2 = (412856739)$.

\paragraph{ER --- rekombinace hran}

Pozorování: PMX, OX i~CX zachovávají jen asi $60\,\%$ hran rodičů. Pro TSP
je ale právě \emph{hrana} nositelem kvality. \emph{Edge recombination}
(Whitley) proto:
\begin{enumerate}
\item Pro každé město sestaví \emph{seznam sousedů} ze~\textbf{obou} rodičů.
\item Začne v~náhodném (nebo prvním) městě.
\item Dalším městem je ten dosud nenavštívený soused, který má \emph{nejméně}
  zbývajících sousedů (heuristika: nenechávat města \uv{osamocená}); při
  shodě rozhodne náhoda.
\item Pokud seznam sousedů dojde prázdný, pokračuje se náhodným nenavštíveným
  městem (nastává v~$1$--$1{,}5\,\%$ případů).
\end{enumerate}
\textbf{ER2} navíc značí hrany, které se vyskytují v~\emph{obou} rodičích
(symbolem \#), a~upřednostňuje je při výběru.

\textbf{Příklad.} $P_1 = (123456789)$, $P_2 = (412876935)$ (cyklické cesty).
Seznamy sousedů: $1{:}\,9,2,4$; $2{:}\,1,3,8$; $3{:}\,2,4,9,5$; $4{:}\,3,5,1$;
$5{:}\,4,6,3$; $6{:}\,5,7,9$; $7{:}\,6,8$; $8{:}\,7,9,2$; $9{:}\,8,1,6,3$.
Start v~1: kandidáti 9 (4~sousedi), 2 (3), 4 (3) --- 9 vypadává, mezi 2 a~4
losujeme, vyjde 4. Sousedé 4 jsou 3 a~5; 3 má ještě 3 volné sousedy, 5 jen 2,
bereme tedy 5, atd.\ --- výsledek např.\ $(145678239)$.

\emph{Technická poznámka:} korektně se navštívené město nejprve odstraní ze
\emph{všech} seznamů a~teprve pak se počítají zbývající sousedé (výše jsou
u~startu uvedeny počty ještě \emph{před} odebráním jedničky: po odebrání by
bylo 9$:$3, 2$:$2, 4$:$2). Na pořadí kandidátů to zde nic nemění, ale při
vlastním počítání na papíře je třeba postupovat takto, jinak vycházejí
jiné remízy.

\paragraph{Mutace permutací}

\begin{itemize}
\item \textbf{Swap:} prohození dvou náhodných měst.
\item \textbf{Insert:} vyjmutí města a~vložení jinam (posun zbytku).
\item \textbf{Inverze (2-opt krok):} obrácení souvislého úseku --- pro TSP
  nejúčinnější, protože mění jen dvě hrany.
\item \textbf{Scramble:} náhodné přeházení souvislého úseku.
\item \textbf{Výměna podcest} a~heuristické mutace (2-opt, 3-opt, Lin-Kernighan).
\end{itemize}

\subsubsection{Další reprezentace}

Stromy (odst.~\ref{sec:gp}), grafy a~DAG (Cartesian GP, neuroevoluce),
matice (rozvrhy, TSP jako matice předchůdců), konečné automaty (EP),
seznamy pravidel (LCS, odst.~\ref{sec:lcs}), agenti umělého života.

\subsection{Genetické a~evoluční programování}
\label{sec:gp}

\subsubsection{Historie}

Myšlenku evoluce programů zmínil už Alan Turing v~50.\ letech. Forsyth
(1980, systém BEAGLE) ji aplikoval na rozpoznávání vzorů, Nichael Cramer
(1985) poprvé popsal stromové jedince a~John Koza (1989--1992) formuloval
stromové genetické programování v~dnešní podobě (včetně patentu).

\begin{defbox}[Genetické programování (\emph{genetic programming}, GP)]
Evoluční algoritmus, jehož jedinci jsou \textbf{spustitelné struktury} ---
programy, výrazy, obvody, modely. Fitness se určuje \emph{spuštěním} jedince
na testovacích datech. Velikost i~tvar jedince nejsou pevné, mění se v~průběhu
evoluce.
\end{defbox}

\begin{algorithm}[!ht]
\caption{Obecné schéma GP}
\begin{algorithmic}[1]
\State vygeneruj počáteční populaci náhodných programů
\While{není nalezeno dost dobré řešení}
  \State ohodnoť programy jejich \textbf{spuštěním} na trénovacích datech
  \State selekce podle fitness (typicky turnaj)
  \State křížení dvou programů, mutace programů
  \State vytvoř novou populaci
\EndWhile
\end{algorithmic}
\end{algorithm}

\subsubsection{Stromové GP}

\paragraph{Reprezentace.} Program je \emph{syntaktický strom}. Definujeme dvě
množiny:
\begin{itemize}
\item \textbf{Množina terminálů} $T$ --- listy: proměnné, konstanty, funkce bez
  argumentů (např.\ čtení senzoru).
\item \textbf{Množina neterminálů (funkcí)} $F$ --- vnitřní uzly s~danou
  \emph{aritou}: aritmetické operace, logické spojky, podmínky, cykly.
\end{itemize}
Požaduje se \emph{uzavřenost} (\emph{closure}) --- každá funkce musí umět
zpracovat libovolný výstup libovolné jiné (odtud \uv{chráněné dělení}
$\%$, které pro nulového dělitele vrací 1) --- a~\emph{dostatečnost}
(\emph{sufficiency}) --- z~$T \cup F$ musí být řešení vůbec vyjádřitelné.

Příklad: strom \texttt{(+ (* x x) (sin y))} reprezentuje výraz $x^2 + \sin y$.

\paragraph{Inicializace.}
\begin{itemize}
\item \textbf{Grow:} uzly se losují z~$T \cup F$, dokud se nedosáhne limitu
  hloubky/počtu uzlů $\Rightarrow$ nepravidelné stromy různé velikosti.
\item \textbf{Full:} do dané hloubky se losuje jen z~$F$, na poslední úrovni
  jen z~$T$ $\Rightarrow$ plné, pravidelné stromy.
\item \textbf{Ramped half-and-half} (Koza): polovina populace metodou grow,
  polovina metodou full, pro rovnoměrně rozložené hloubky (např.\ 2--6).
  Standardní volba; vytváří spíše \uv{košaté} pravidelné tvary, což vyhovuje
  úlohám se symetriemi, kde jsou všechny vstupy podobně důležité.
\item \textbf{Uniformní vzorkování} všech možných stromů dává nepravidelnější
  tvary (některé terminály blíže ke kořeni).
\item \textbf{Seeding:} vložení částečných či expertních řešení do počáteční
  populace --- zrychluje, ale hrozí předčasná konvergence.
\end{itemize}

\paragraph{Křížení.}
\begin{itemize}
\item \textbf{Výměna podstromů} (standardní Kozovo křížení): v~každém rodiči se
  zvolí náhodný uzel a~podstromy pod nimi se vymění. Neterminály mívají vyšší
  pravděpodobnost výběru (typicky 90 \% neterminál, 10 \% terminál) ---
  jinak by se často vyměňovaly jen listy.
\item \textbf{Homologní křížení} se snaží zachovat pozici genetického
  materiálu. Nejprve se najde \emph{společná oblast} $C$ (procházením od kořene,
  dokud se shoduje arita uzlů):
  \begin{itemize}
  \item \emph{Jednobodové:} bod křížení se volí ze společné oblasti.
  \item \emph{Uniformní (GPUX, Poli a~Langdon):} každý uzel společné oblasti se
    s~konstantní pravděpodobností vymění; uzly \emph{uvnitř} $C$ se vymění
    bez dopadu na jejich podstromy, uzly na \emph{hranici} $C$ se vymění
    i~s~podstromy. V~raných fázích se tak vyměňují velké podstromy,
    s~konvergencí populace se operátor stává stále lokálnějším.
  \end{itemize}
\item \textbf{Křížení zachovávající kontext} (\emph{context preserving}):
  uzel je identifikován \emph{souřadnicemi} --- posloupností odboček na cestě
  od kořene (např.\ $(2,1,3,1)$). \emph{Silná} varianta povoluje křížení jen
  mezi uzly se \emph{stejnými} souřadnicemi, \emph{slabá} je volnější.
\end{itemize}

\paragraph{Mutace.} V~GP je nutné (nikoli jen užitečné) používat \emph{více
typů} mutací:
\begin{itemize}
\item \textbf{Podstromová mutace} --- nahradí náhodný podstrom novým náhodným.
\item \textbf{Size-fair} podstromová mutace --- velikost nového podstromu se
  losuje tak, aby odpovídala odstraněnému.
\item \textbf{Mutace uzlu} (\emph{node replacement}) --- záměna uzlu za jiný
  \emph{stejné arity}.
\item \textbf{Hoist} --- jedinec se nahradí svým podstromem (zmenšuje).
\item \textbf{Shrink} --- podstrom se nahradí terminálem (zmenšuje).
\item \textbf{Permutační mutace} --- prohození podstromů jednoho neterminálu.
\item \textbf{Mutace konstant} --- GP tradičně špatně dolaďuje číselné hodnoty.
  Pomáhá aritmetická mutace konstant nebo dokonce lokální optimalizace
  \emph{všech} konstant ve stromu (hill climbing, gradientní metoda) ---
  memetický přístup.
\end{itemize}

\paragraph{Fitness a~selekce.} Fitness = chyba na trénovacích datech
(symbolická regrese), počet správných odpovědí, kvalita chování v~simulaci.
Selekce bývá turnajová. Riziko přeučení je stejné jako u~strojového učení
(řeší se validační množinou, penalizací složitosti).

\subsubsection{Bloat}

\begin{defbox}[Bloat]
\emph{Bloat} je tendence programů v~GP růst do velikosti bez odpovídajícího
zlepšení fitness. Typicky jde o~hromadění \emph{intronů} --- částí kódu, které
nemají vliv na výstup (např.\ \texttt{(+ x 0)}, \texttt{(if false ...)}).
\end{defbox}

Proč k~němu dochází: v~prostoru programů je mnohem víc velkých programů se stejným
chováním než malých, takže už samotný neutrální drift tlačí k~růstu; navíc introny chrání
jedince před destruktivním křížením, takže se \uv{svezou} s~dobrým kódem.
V~přírodě růst naráží na fyzikální a~energetické limity, v~GP je nutné bojovat s~bloatem
explicitně:
\begin{itemize}
\item \textbf{Tvrdé limity} na hloubku či počet uzlů (potomek překračující limit se zahodí
  nebo nahradí rodičem).
\item \textbf{Penalizace ve fitness} (\emph{parsimony pressure}):
  $f' = f - \alpha \cdot \text{velikost}$; případně \emph{lexikografická} varianta:
  při shodné fitness vyhrává menší jedinec.
\item \textbf{Anti-bloat operátory} --- hlídat, aby křížení a~mutace jedince příliš
  nezvětšovaly (size-fair varianty), a~speciální zmenšující mutace \emph{hoist}
  a~\emph{shrink}.
\item \textbf{Vícekriteriální přístup} --- velikost jako druhé kritérium
  (odst.~\ref{sec:moea}).
\end{itemize}

\subsubsection{ADF a~typované GP}

\begin{defbox}[ADF (\emph{automatically defined functions})]
Podprogramy evolvované společně s~hlavním programem --- nutná podmínka
\textbf{modularity}, vyšší složitosti a~schopnosti zachytit symetrie úlohy.
Jedinec se skládá z~větve \texttt{main} a~jedné či více větví ADF; každá ADF je
charakterizována svou \emph{aritou} a~může mít omezenou (jinou) množinu terminálů
a~neterminálů. Volání ADF se v~hlavním programu chová jako \textbf{nový neterminál}
dané arity. Genetické operátory pracují na větvích \textbf{odděleně}
(main se kříží s~mainem, ADF$_1$ s~ADF$_1$).

Rozšíření: automaticky definované smyčky, iterace a~rekurze. Alternativou je
\emph{koevoluce} dvou populací --- hlavních programů a~podprogramů
(odst.~\ref{sec:koevoluce}).
\end{defbox}

\textbf{Vynucení struktury.} \emph{Silně typované GP} dává každému uzlu typ vstupů
a~výstupu a~operátory smějí vytvářet jen typově korektní stromy; \emph{gramatické GP}
popíše přípustné tvary stromů bezkontextovou gramatikou. Obojí redukuje prohledávaný
prostor na smysluplné programy.

\subsubsection{Lineární a~grafové GP}

\paragraph{Lineární GP.} Program je posloupnost instrukcí (často strojový nebo
byte kód) pracující nad registry. Výhody: jednoduché operátory (vektorové
křížení a~mutace), rychlá interpretace, přirozená blízkost skutečnému
hardwaru. Nevýhoda: vysoké riziko vzniku nesmyslných programů. Oblíbené
v~umělém životě a~evoluci herních botů. Slavný příklad: \textbf{Tierra}
(T.~Ray) --- simulace umělého ekosystému, kde jedinci jsou sebe-kopírující
programy ve strojovém kódu soutěžící o~paměť a~procesorový čas; emergentně
vznikli paraziti, hyperparaziti a~obrana proti nim. Následovníci: Avida,
Avida-ED.

\paragraph{Grafové GP.} Program je obecný (často acyklický) graf. Původně
rozšíření stromového GP na paralelní programy; ukázalo se, že grafy dobře
popisují elektrické obvody, konečné automaty, neuronové sítě, plány.
Problémem jsou složité operátory --- jak křížit obecné grafy.

\paragraph{Kartézské GP (CGP).} \mimo{} Elegantní obchvat: DAG se reprezentuje svým
rovinným obrazem v~2D mřížce se souřadnicemi uzlů (fenotyp), zatímco genotyp
je \textbf{lineární} posloupnost celých čísel popisujících pro každý uzel
jeho funkci (index do tabulky funkcí), jeho vstupy a~nakonec výstupní uzly
grafu. Vlastnosti:
\begin{itemize}
\item Mutace je bodová mutace lineárního genomu; musí se hlídat, aby vznikaly
  jen platné funkce a~platná propojení (omezení hloubky spojů).
\item Malá změna genotypu může mít velký dopad na fenotyp; naopak řada
  mutací je \emph{neutrální} (uzly nepřipojené k~výstupu tvoří přirozené
  introny, což CGP prospívá).
\item Křížení je málo důležité; naivní jednobodové je příliš destruktivní,
  používá se aritmetická varianta a~řezy na hranicích uzlů.
\item Selekce: obvykle $(1+\lambda)$-ES s~$\lambda = 4$.
\item Genom má konstantní délku (bloat je omezen mřížkou).
\end{itemize}

\paragraph{Vývojové (developmental) GP.} Místo přímé evoluce grafu se
evolvuje \emph{strom-program, který graf postupně staví}. Začíná se
\uv{embryem} a~aplikují se operace pro strukturu (sériové/paralelní zapojení,
třícestné dělení) a~pro elementy (vlož rezistor, kondenzátor). Použití:
Koza --- evoluce elektrických obvodů (znovuobjevení patentovaných zapojení),
Gruau --- \emph{cellular encoding} pro neuronové sítě, AutoML --- evoluce
modelů scikit-learn.

\subsubsection{Gramatická evoluce}

\begin{defbox}[Gramatická evoluce (\emph{grammatical evolution}, GE)]
Genotyp je \textbf{lineární seznam celých čísel} (proměnné délky), fenotyp je program
odvozený z~bezkontextové gramatiky. Číslo z~genomu určuje, které pravidlo se použije pro
expanzi nejlevějšího neterminálu:
\[
\text{pravidlo} \;=\; (\text{gen} \bmod \text{počet pravidel pro tento neterminál}).
\]
\end{defbox}

Použití modula činí kódování \textbf{robustním} (každý genom je interpretovatelný)
a~protože je genom lineární, lze používat \textbf{standardní operátory GA} --- to je hlavní
praktická výhoda. Gramatika navíc zaručuje syntaktickou správnost. Nevýhodou je nelokálnost:
malá změna genotypu (zejména na začátku) může zcela změnit fenotyp (\emph{ripple effect}).
Nedokončí-li se derivace, čte se genom znovu od začátku (\emph{wrapping}) s~limitem počtu
obalení.

\subsubsection{Evoluční programování}

\begin{defbox}[Evoluční programování (\emph{evolutionary programming}, EP)]
Větev EA (L.~Fogel, 1965 --- starší než GA), jejímž cílem bylo evolvovat
\uv{umělou inteligenci}: agenta, který predikuje stav prostředí a~volí vhodnou
akci. Agent je reprezentován \textbf{konečným automatem}; platí
\textbf{genotyp = fenotyp} (žádné zvláštní kódování).
\textbf{Křížení se nepoužívá}, veškerá variabilita pochází z~mutací.
\end{defbox}

\paragraph{EP nad automaty.} Prostředí je posloupnost symbolů konečné abecedy.
Automat dostává symboly na vstup, jeho výstup se porovnává s~následujícím
symbolem posloupnosti; fitness = úspěšnost predikce (absolutní chyba, střední
kvadratická chyba), často s~penalizací za velikost automatu.
Mutace: změna výstupního symbolu, změna cílového stavu přechodu, přidání
stavu, odebrání stavu, změna počátečního stavu. Polovina populace se nahrazuje
novými automaty.

\paragraph{Slavný příklad --- predikce prvočísel.} Cílem je predikovat
posloupnost
\[
111010100010100010100010000010\ldots,
\]
kde 1 znamená \uv{toto číslo je prvočíslo} (tj.\ naučit se Eratosthenovo
síto). Postupuje se iterativně: nejprve se učí krátká posloupnost, pak se
prodlužuje o~jeden symbol. Fitness: bod za každý uhodnutý symbol minus
$0{,}01 \cdot N$ za počet stavů. Výsledek je poučný: automaty se
zjednodušovaly, až nakonec vznikl jednostavový automat generující stále 0 ---
prvočísla jsou totiž řídká, takže konstantní odpověď \uv{není prvočíslo} má
vysokou úspěšnost. Klasická ukázka toho, jak evoluce zneužije špatně
navrženou fitness.

\paragraph{EP dnes.}
\begin{itemize}
\item Kódování má být \emph{přirozené} pro úlohu (obvykle genotyp = fenotyp).
\item Sada \uv{chytrých}, doménově specifických mutací; žádné křížení.
\item Rodičovská selekce: každý jedinec je právě jednou rodičem.
\item Environmentální selekce: turnaj mezi rodiči a~potomky.
\item EP nad reálnými vektory obsahuje --- stejně jako ES --- rozptyly mutací
  v~genomu ($\sigma_j' = \sigma_j(1 + \alpha N(0,1))$, poté
  $x_j' = x_j + \sigma_j' N(0,1)$). EP a~ES tak prakticky splynuly.
\end{itemize}

\subsubsection{Je křížení k~něčemu? Experiment s~bezhlavým kuřetem}

Tradiční obhajoba křížení = výměna stavebních bloků. Kritika (typická pro EP):
křížení je jen podivná velká náhodná mutace, která navíc nerespektuje kódování.

\begin{defbox}[Headless chicken test (Jones, 1995)]
Porovnáme standardní křížení s~\emph{náhodným} křížením, kde se jedinec kříží
s~\textbf{náhodně vygenerovaným} řetězcem. Je-li výsledek stejný nebo lepší,
pak zkoumaný operátor ve skutečnosti nefunguje jako křížení, ale jako
makromutace --- \uv{nenazývejme bezhlavé kuře kuřetem, i~když má mnoho
kuřecích rysů}.
\end{defbox}

Interpretace: existují-li jasné stavební bloky, je pravé křížení lepší.
Vyhraje-li náhodné křížení, znamená to, že v~dané úloze/kódování stavební
bloky nejsou --- špatné kódování, špatný operátor, nebo úloha pro křížení
nevhodná.

\subsection{Shrnutí ke~zkoušce}

\begin{itemize}
\item EA = populační stochastické prohledávání; cyklus \emph{inicializace $\to$
  ohodnocení $\to$ rodičovská selekce $\to$ rekombinace $\to$ mutace $\to$
  ohodnocení $\to$ environmentální selekce}.
\item Variace (křížení, mutace) tvoří diverzitu, selekce ji odebírá a~směřuje
  hledání; návrh EA = hledání rovnováhy explorace/exploatace.
\item Selekce: \textbf{ruleta} ($p_i = f_i/\sum f_j$, vysoký rozptyl, citlivá na
  škálování), \textbf{SUS} (stejné střední hodnoty, nízký rozptyl),
  \textbf{turnaj} (tlak řízen $k$, invariantní k~transformacím fitness, nepotřebuje
  explicitní fitness), \textbf{pořadová} (konstantní tlak),
  \textbf{Boltzmannova} (tlak roste s~časem).
\item Elitismus zaručuje monotonii nejlepší fitness; generační vs.\ steady-state
  model se liší podílem nahrazované populace.
\item No Free Lunch: neexistuje univerzálně nejlepší algoritmus $\Rightarrow$
  doménová specializace reprezentace a~operátorů.
\item Historické větve GA / ES / EP / GP a~jejich charakteristické volby
  (viz srovnávací tabulka) je dobré umět vyjmenovat zpaměti.
\item SGA = binární řetězce + ruleta + 1-bodové křížení + bit-flip mutace
  + generační náhrada. Parametry: $n$, $p_C \approx 0{,}7$, $p_M \approx 1/m$.
\item Binární kódování reálných čísel: $x = a + z(b-a)/(2^k-1)$; problém
  Hammingových útesů řeší Grayův kód.
\item Křížení: 1-bodové (rozbití schématu s~pravděpodobností $d(S)/(m-1)$),
  2-bodové (bez okrajového biasu), uniformní (nejvíce disruptivní, bez
  pozičního biasu).
\item Mutace je v~GA sekundární --- brání nevratné ztrátě alel; inverze je
  historický operátor pro evoluci uspořádání genů.
\item Škálování fitness (lineární, sigma truncation, mocninné, pořadové,
  Boltzmannovo) opravuje slabiny fitness-proporcionální selekce: dominanci
  supermanů na začátku a~ztrátu tlaku na konci.
\item Reprezentace $+$ operátory $=$ fitness krajina. Pro nominální atributy
  jen nezaujaté mutace, pro ordinální i~zaujaté (creep).
\item Reálná reprezentace: gaussovská mutace $x' = x + \sigma N(0,1)$;
  aritmetické křížení $\alpha \vec x + (1-\alpha)\vec y$ (zmenšuje rozptyl),
  proto BLX-$\alpha$ a~SBX.
\item Permutace: PMX (pozice + mapování), OX (relativní pořadí),
  CX (absolutní pozice pomocí cyklů), ER (hrany, nejlepší pro TSP).
  Umět předvést PMX a~OX na konkrétním osmiprvkovém příkladu!
\item Mutace permutací: swap, insert, inverze (2-opt), scramble.
\item Stromové GP (Koza): terminály + neterminály, uzavřenost a~dostatečnost;
  inicializace grow / full / ramped half-and-half; křížení = výměna podstromů
  (neterminály s~vyšší pravděpodobností), více typů mutací včetně mutace
  konstant; fitness = spuštění programu.
\item Bloat = růst velikosti bez zlepšení fitness (introny); tři vysvětlení:
  replication accuracy, removal bias, crossover bias. Obrana: limity
  velikosti, penalizace (parsimony pressure), anti-bloat operátory
  (hoist, shrink, size-fair), vícekriteriální přístup.
\item ADF = evolvované podprogramy (modularita); operátory pracují na větvích
  odděleně. Typované GP a~gramatické GP omezují prostor na smysluplné programy.
\item Varianty: lineární GP (byte kód, Tierra), grafové GP,
  Cartesian GP (lineární genotyp $\to$ 2D DAG, bodová mutace, $(1+4)$-ES),
  vývojové GP (strom staví graf --- obvody, sítě).
\item Gramatická evoluce: lineární genom celých čísel, výběr pravidla
  gramatiky pomocí modula, derivační strom $\to$ fenotyp; robustní kódování,
  standardní operátory GA, ale nelokální mapování.
\item EP: konečné automaty, genotyp = fenotyp, jen mutace, turnajová selekce
  z~rodičů i~potomků; příklad predikce prvočísel (degenerace fitness).
\item Headless chicken experiment: test, zda křížení skutečně kombinuje
  stavební bloky, nebo je jen makromutací.
\end{itemize}

%=====================================================================
\section{Teorie schémat a~pravděpodobnostní modely SGA}
\label{sec:schemata}
%=====================================================================

Kapitola odpovídá na otázku, \emph{proč} genetické algoritmy fungují. Nejprve klasická
Hollandova teorie schémat (a~její kritika, která je u~zkoušky stejně důležitá jako věta
sama), pak exaktní Voseho model, který teorii schémat nahradil.

\subsection{Teorie schémat}
\label{ssec:teorie-schemat}

Teorie schémat je první pokus formálně vysvětlit, \emph{proč} genetické
algoritmy fungují. Pochází od Hollanda (1975) a~popularizoval ji Goldberg
(1989). Dnes je považována za historicky významnou, ale nedostatečnou; přesto
je standardní součástí zkoušky a~její kritika je stejně důležitá jako
věta sama.

\subsubsection{Schéma, jeho řád a~definující délka}

\begin{defbox}[Schéma (\emph{schema})]
Nechť jedinec je slovo délky $m$ nad abecedou $\{0,1\}$. \textbf{Schéma} $S$ je
slovo délky $m$ nad abecedou $\{0,1,*\}$, kde $*$ znamená \uv{na hodnotě
nezáleží} (\emph{don't care}). Schéma reprezentuje množinu všech jedinců, kteří
se s~ním shodují na všech pevných pozicích; schéma s~$r$ hvězdičkami
reprezentuje $2^r$ jedinců.
\end{defbox}

Příklad: schéma $1*0*$ reprezentuje $\{1000, 1001, 1100, 1101\}$; schéma
$**\ldots*$ reprezentuje celý prostor.

\begin{defbox}[Řád, definující délka, fitness schématu]
Pro schéma $S$ délky $m$ definujeme:
\begin{itemize}
\item \textbf{Řád} $o(S)$ = počet pevných pozic (počet nul a~jedniček).
\item \textbf{Definující délka} $d(S)$ = vzdálenost mezi první a~poslední
  pevnou pozicí ($d(S) = 0$ pro schéma s~jedinou pevnou pozicí).
\item \textbf{Fitness schématu} $F(S,t)$ = průměrná fitness jedinců
  v~populaci $P(t)$, kteří schéma $S$ reprezentují. \emph{Pozor:} tato hodnota
  závisí na aktuální populaci, nikoli jen na $S$ a~$f$.
\end{itemize}
\end{defbox}

\textbf{Příklad.} Pro $S = 1{*}{*}0{*}1$ ($m=6$) je $o(S) = 3$
(pozice 1, 4, 6) a~$d(S) = 6 - 1 = 5$.

\paragraph{Kombinatorika schémat.}
\begin{itemize}
\item Schémat délky $m$ je celkem $3^m$.
\item Jeden jedinec délky $m$ je reprezentantem $2^m$ schémat (na každé pozici
  buď jeho bit, nebo $*$).
\item Populace $n$ jedinců reprezentuje mezi $2^m$ a~$n \cdot 2^m$ schémat
  (podle míry jejich podobnosti).
\end{itemize}

\subsubsection{Hollandova věta o~schématech}

\begin{veta}[Věta o~schématech (TST), Holland 1975]
V~jednoduchém genetickém algoritmu s~fitness-proporcionální selekcí,
jednobodovým křížením s~pravděpodobností $p_C$ a~bitovou mutací
s~pravděpodobností $p_M$ platí pro očekávaný počet reprezentantů schématu $S$
v~následující generaci:
\[
\mathbb{E}\big[C(S,t+1)\big] \;\ge\; C(S,t)\,\frac{F(S,t)}{\bar f(t)}
\left[\,1 - p_C \frac{d(S)}{m-1} - p_M\, o(S)\right].
\]
Slovně: \textbf{krátká} (malá definující délka), \textbf{nízkého řádu}
a~\textbf{nadprůměrná} schémata se v~po sobě jdoucích generacích množí
exponenciálně.
\end{veta}

\paragraph{Odvození}

Označme $C(S,t)$ počet jedinců v~$P(t)$ reprezentujících $S$, $n = |P(t)|$,
$\bar f(t)$ průměrnou fitness populace. Postupujeme ve třech krocích podle
operátorů.

\paragraph{Krok 1 --- selekce.} Pravděpodobnost výběru jedince $v$ ruletou je
$p_s(v) = f(v)/F(t)$, kde $F(t) = \sum_{u \in P(t)} f(u)$. Pravděpodobnost, že
vybraný jedinec reprezentuje $S$, je
\[
p_s(S) \;=\; \sum_{v \in S \cap P(t)} \frac{f(v)}{F(t)}
\;=\; \frac{C(S,t)\,F(S,t)}{F(t)}.
\]
Při $n$ nezávislých tazích tedy
\[
\mathbb{E}\big[C(S,t+1)\big] \;=\; n\, p_s(S) \;=\; C(S,t)\,\frac{F(S,t)}{\bar f(t)},
\qquad \bar f(t) = F(t)/n .
\]
Je-li schéma nadprůměrné o~$\varepsilon$, tj.\
$F(S,t) = \bar f(t)\,(1+\varepsilon)$ a~tento poměr se udržuje v~čase, dostáváme
\[
C(S,t+1) = C(S,t)(1+\varepsilon)
\qquad\Longrightarrow\qquad
C(S,t) = C(S,0)\,(1+\varepsilon)^{t},
\]
tedy \textbf{exponenciální růst} (resp.\ pokles pro $\varepsilon<0$).

\paragraph{Krok 2 --- křížení.} Jednobodové křížení volí bod řezu z~$m-1$
možností. Schéma je \emph{jistě} rozbito, padne-li řez mezi jeho první
a~poslední pevnou pozici, tj.\ s~pravděpodobností nejvýše $d(S)/(m-1)$
(nejvýše proto, že i~při řezu uvnitř může být schéma obnoveno, pokud druhý
rodič obsahuje tytéž hodnoty). Křížení se provádí s~pravděpodobností $p_C$,
takže
\[
p_{\mathrm{surv}}^{\mathrm{cross}}(S) \;\ge\; 1 - p_C\,\frac{d(S)}{m-1}.
\]

\paragraph{Krok 3 --- mutace.} Schéma přežije, pokud se nezmění žádná z~jeho
$o(S)$ pevných pozic:
\[
p_{\mathrm{surv}}^{\mathrm{mut}}(S) = (1-p_M)^{o(S)} \;\approx\; 1 - p_M\,o(S)
\quad \text{pro } p_M \ll 1 .
\]

\paragraph{Složení.} Za předpokladu (přibližné) nezávislosti obou destruktivních
jevů a~zanedbáním jejich součinu:
\[
\mathbb{E}\big[C(S,t+1)\big] \;\ge\; C(S,t)\,\frac{F(S,t)}{\bar f(t)}
\left[1 - p_C \frac{d(S)}{m-1} - p_M o(S)\right]. \qquad\blacksquare
\]

\paragraph{Číselný příklad}

Vezměme Goldbergovu populaci ze str.~\pageref{ex:goldberg} ($m=5$, $n=4$):
$01101\,(169)$, $11000\,(576)$, $01000\,(64)$, $10011\,(361)$;
$\bar f = 292{,}5$. Parametry $p_C = 0{,}7$, $p_M = 0{,}001$.

\begin{center}
\begin{tabular}{@{}lccccc@{}}
\toprule
schéma $S$ & reprezentanti & $F(S)$ & $o(S)$, $d(S)$ & přežití & $\mathbb{E}[C(S,t{+}1)]$ \\
\midrule
$1{*}{*}{*}{*}$ & 11000, 10011 & $468{,}5$ & $1;\ 0$ &
  $1 - 0 - 0{,}001 = 0{,}999$ &
  $2 \cdot 1{,}602 \cdot 0{,}999 = 3{,}2$ \\
$1{*}{*}{*}1$ & 10011 & $361$ & $2;\ 4$ &
  $1 - 0{,}7 - 0{,}002 = 0{,}298$ &
  $1 \cdot 1{,}234 \cdot 0{,}298 = 0{,}37$ \\
$0{*}{*}{*}{*}$ & 01101, 01000 & $116{,}5$ & $1;\ 0$ &
  $0{,}999$ &
  $2 \cdot 0{,}398 \cdot 0{,}999 = 0{,}80$ \\
\bottomrule
\end{tabular}
\end{center}

Interpretace: krátké nadprůměrné schéma $1{*}{*}{*}{*}$ se rozšíří
z~2 na cca 3,2 zástupce. Schéma $1{*}{*}{*}1$ je sice nadprůměrné
($F/\bar f = 1{,}23$), ale jeho velká definující délka je fatální ---
křížení ho rozbije a~očekáváme jeho vymizení. Podprůměrné $0{*}{*}{*}{*}$
klesá. To je přesně obsah věty.

\subsubsection{Hypotéza stavebních bloků}

\begin{defbox}[Hypotéza stavebních bloků (\emph{building blocks hypothesis})]
Genetický algoritmus hledá dobré řešení tak, že \textbf{rekombinuje krátká,
nízkého řádu a~nadprůměrná schémata} --- tzv.\ \emph{stavební bloky} ---
do stále lepších řešení.
Goldbergova metafora: \uv{tak jako dítě staví hrad z~jednoduchých dřevěných
kostek, tak GA hledá téměř optimální řešení skládáním stavebních bloků}.
\end{defbox}

Praktické důsledky:
\begin{itemize}
\item \textbf{Na kódování záleží.} Geny, které spolu interagují, mají ležet
  blízko sebe (\emph{linkage}), jinak je křížení rozbíjí. Odtud snahy o~inverzi,
  \emph{messy GA} a~\emph{linkage learning}.
\item \textbf{Na velikosti populace záleží.} Populace musí obsahovat dostatek
  vzorků od každého stavebního bloku (teorie \emph{population sizing}).
\item \textbf{Předčasná konvergence škodí} --- ztráta diverzity znemožní
  objevení chybějících bloků.
\end{itemize}

\subsubsection{Implicitní paralelismus a~víceruký bandita}

\begin{defbox}[Implicitní paralelismus (\emph{implicit parallelism})]
GA pracuje explicitně s~$n$ jedinci, ale implicitně současně \uv{zpracovává}
mnohem více schémat --- mezi $2^m$ a~$n \cdot 2^m$. Holland odhadl, že počet
schémat, která jsou zpracovávána \emph{užitečně} (rostou exponenciálně a~nejsou
příliš rozbíjena), je řádu $n^3$.
\end{defbox}

Bývá to komentováno jako jediný případ, kdy kombinatorická exploze pracuje
v~náš prospěch.

\paragraph{Souvislost s~úlohou o~banditovi.} Hollandova původní motivace byla
\uv{adaptivní plán} hledající rovnováhu mezi explorací a~exploatací. Formálně:
\begin{itemize}
\item \textbf{Dvouruký bandita.} Máme $N$ mincí a~automat se dvěma pákami
  s~neznámými středními výplatami $\mu_1, \mu_2$ a~rozptyly. Alokujeme
  $n$ mincí horší a~$N-n$ lepší páce. Analytické řešení říká, že optimální
  strategie alokuje \emph{exponenciálně} více pokusů momentálně vítězné páce:
  $N - n^* = O(e^{c\, n^*})$.
\item \textbf{GA jako bandita.} Věta o~schématech říká, že GA alokuje
  exponenciálně více \uv{slotů v~populaci} úspěšnějším schématům --- tedy
  že řeší dilema explorace/exploatace (téměř) optimálně.
\item Původně se soudilo, že SGA hraje jednu hru s~$3^m$ pákami. Ve
  skutečnosti hraje mnoho her paralelně: schémata soutěží o~konkrétní pevné
  pozice, a~schémata řádu $k$ hrají hru s~$2^k$ pákami.
\end{itemize}

\subsubsection{Kritika teorie schémat}

\paragraph{Kdy GA selhává --- deceptivní úlohy}

Uvažujme funkci, v~níž jsou nadprůměrná schémata
\[
S_1 = (111{*}\ldots{*}), \qquad S_2 = ({*}\ldots{*}11),
\]
ale zároveň
\[
f(111{*}{*}{*}{*}{*}11) \;\ll\; f(000{*}{*}{*}{*}{*}00).
\]
GA je selekcí veden
k~jedničkovým blokům, přestože optimum leží jinde. Takové úlohy se nazývají
\textbf{deceptivní} (\emph{deceptive}, Goldberg 1987): dílčí stavební bloky
nevedou po složení ke globálnímu optimu. Deception bývá navrhována jako míra
obtížnosti úlohy pro EA (ne všichni s~tím souhlasí).

Opačným testem jsou \textbf{Royal Road funkce} (Mitchell, Forrest, Holland):
fitness je součtem příspěvků za kompletní \uv{bloky} jedniček (např.\ osm bloků po
osmi bitech), krajina je tedy ušitá přesně na míru hypotéze stavebních bloků a~GA
by na ní měl excelovat. Experiment dopadl opačně: jednoduchý GA prohrál s~náhodně
mutujícím horolezeckým algoritmem (RMHC). Příčinou je \emph{hitchhiking}
(\uv{stopování}): spolu s~nalezeným dobrým blokem se selekcí šíří i~nahodilé
\uv{parazitní} bity na ostatních pozicích a~ničí diverzitu potřebnou k~nalezení
zbývajících bloků. Royal Road funkce jsou proto standardním empirickým
protipříkladem k~naivnímu čtení BBH.

\paragraph{Statické průměrné fitness --- statická BBH}

Grefenstette (1991): lidé si BBH vykládají tak, že GA konverguje k~řešením
s~dobrou \emph{skutečnou statickou} průměrnou fitness schématu (přes celý
prostor). GA však odhaduje fitness schématu jen \emph{ze vzorků v~populaci},
a~ten odhad je vychýlený ze dvou důvodů:

\begin{itemize}
\item \textbf{Kolaterální konvergence} (\emph{collateral convergence}).
  Jakmile GA někam konverguje, přestává vzorkovat schémata rovnoměrně.
  Je-li např.\ schéma $111{*}\ldots{*}$ dobré, po několika generacích má
  skoro každý jedinec tento prefix. Pak je ale téměř každý vzorek schématu
  ${*}{*}{*}000\ldots$ zároveň vzorkem $111000\ldots$ a~fitness prvního
  schématu je odhadnuta zcela zkresleně.
\item \textbf{Velký rozptyl fitness schématu.} Uvažujme
  $f(x) = 2$ pro $x \in 111{*}\ldots{*}$, $f(x)=1$ pro $x \in 0{*}\ldots{*}$
  a~$0$ jinak. Statická průměrná fitness schématu $1{*}\ldots{*}$ je $\approx 1/2$,
  zatímco $F(0{*}\ldots{*}) = 1$. GA však odhadne
  $F(1{*}\ldots{*}) \approx 2$, protože v~populaci převáží jedinci
  s~prefixem 111. Vysoký rozptyl vede k~systematickému zkreslení a~k~tomu,
  že GA nevzorkuje schémata nezávisle.
\end{itemize}

\paragraph{Souhrn slabin}

\begin{itemize}
\item Věta je \textbf{nerovnost} a~pouze \textbf{jednokrokový} výsledek. Její
  iterace na více generací vyžaduje předpoklad, že $F(S,t)/\bar f(t)$ zůstává
  konstantní --- což typicky neplatí (obojí se dynamicky mění).
\item Populace jsou \textbf{konečné a~malé} $\Rightarrow$ výběrové chyby;
  exponenciální růst je jen střední hodnota.
\item Věta zohledňuje jen \textbf{destruktivní} účinek operátorů, nikoli
  jejich \emph{konstruktivní} roli (křížení může schéma i~vytvořit).
\item Soutěž schémat je \textbf{silně závislá}, nikoli nezávislá jako páky
  bandity.
\item Analýza je \uv{on-line}: SGA maximalizuje průběžný výkon, hodí se tedy
  spíše pro adaptivní úlohy, zatímco nejběžnější použití je nalezení jednoho
  nejlepšího řešení (off-line kritérium).
\end{itemize}

Přesto zůstává TST užitečnou intuicí: \emph{co je krátké, jednoduché a~dobré,
to se šíří}.

\subsection{Pravděpodobnostní modely jednoduchého GA}
\label{sec:vose}


V~90.\ letech (Vose, Liepins, Nix, Whitley) vznikly \emph{exaktní} modely SGA,
které nahradily přibližnou teorii schémat. Cílem bylo:
popsat přesně, jak vypadá populace, popsat zobrazení přechodu mezi populacemi,
zkoumat jeho vlastnosti a~asymptotické chování SGA. Rozlišují se dva přístupy:
\textbf{nekonečné populace} (deterministický dynamický systém) a~\textbf{konečné
populace} (Markovův řetězec).

\subsubsection{Zjednodušený SGA}

Pro analýzu se SGA ještě zjednoduší: v~každém kroku se vyberou dva rodiče,
s~pravděpodobností $p_C$ se zkříží, \textbf{jeden z~potomků se zahodí},
zbylý zmutuje a~vloží se do nové populace. Tím je tvorba každého jedince
nové populace nezávislá.

\subsubsection{Formalizace populace}

Každý binární řetězec délky $l$ ztotožníme s~číslem
$i \in \{0,1,\dots,2^l-1\}$ (např.\ $00000111 \equiv 7$). Populaci v~čase $t$
popisují dva vektory délky $2^l$:

\begin{defbox}[Vektor populace a~vektor selekce]
\begin{itemize}
\item $p(t) = (p_0(t),\dots,p_{2^l-1}(t))$, kde $p_i(t)$ je \emph{relativní
  četnost} řetězce $i$ v~populaci. Platí $p_i \ge 0$, $\sum_i p_i = 1$, tedy
  $p(t)$ leží v~simplexu $\Lambda$.
\item $s(t) = (s_0(t),\dots,s_{2^l-1}(t))$, kde $s_i(t)$ je pravděpodobnost,
  že selekce vybere řetězec $i$.
\end{itemize}
Vektor $p$ popisuje \emph{složení} populace, vektor $s$ \emph{selekční
pravděpodobnosti}.
\end{defbox}

\subsubsection{Operátor $\mathcal{G}$}

Cílem je najít operátor $\mathcal{G}: \Lambda \to \Lambda$ takový, že
\[
p(t+1) = \mathcal{G}\big(p(t)\big),
\]
a~rozložit ho na dvě části: $\mathcal{G} = \mathcal{M} \circ \mathcal{F}$,
kde $\mathcal{F}$ odpovídá selekci (fitness) a~$\mathcal{M}$ tzv.\ míchání
(\emph{mixing}) --- křížení a~mutaci dohromady. Pak stačí operátor iterovat
z~počáteční populace a~známe úplné chování SGA.

\paragraph{Část $\mathcal{F}$ --- selekce.} Definujme diagonální matici
$\mathbf{F}$ s~$\mathbf{F}_{ii} = f(i)$ a~$\mathbf{F}_{ij} = 0$ pro $i \ne j$.
Fitness-proporcionální selekce je pak přesně
\[
s(t) \;=\; \frac{\mathbf{F}\,p(t)}{\sum_j \mathbf{F}_{jj}\,p_j(t)} .
\]
Čitatel $(\mathbf{F}p)_i = f(i)\,p_i$ je \uv{váha} řetězce $i$ v~populaci,
jmenovatel je normalizační konstanta (průměrná fitness populace).
Vypustíme-li nyní křížení a~mutaci, je nová populace prostě $n$ nezávislých
tahů podle rozdělení $s(t)$, takže $\mathbb{E}[p(t+1)] = s(t)$. Dosazením
do definice selekce dostáváme $\mathbb{E}[s(t+1)] \propto \mathbf{F}\,s(t)$,
tedy iterování $\mathcal{F}$ je (až na normalizaci) opakované násobení
maticí $\mathbf{F}$ --- \emph{mocninná metoda}. Ta konverguje k~vlastnímu
vektoru příslušnému největší diagonální hodnotě, tj.\ k~populaci složené
výhradně z~kopií nejlepšího jedince \emph{přítomného v~počáteční populaci}.
U~konečných populací způsobují statistické chyby odchylky od střední hodnoty;
čím větší populace, tím menší odchylka. Nekonečné populace jsou exaktní, byť
nepraktické.

\textbf{Mini-příklad.} Pro $l = 2$ máme čtyři řetězce $00,01,10,11$, tedy
indexy $0,1,2,3$. Nechť $f = (1, 2, 3, 4)$ a~populace osmi jedinců obsahuje
$2\times 00$, $4\times 01$, $1\times 10$, $1\times 11$, tj.\
$p = (0{,}25;\,0{,}5;\,0{,}125;\,0{,}125)$. Pak
$\mathbf{F}p = (0{,}25;\,1{,}0;\,0{,}375;\,0{,}5)$, součet je $2{,}125$
(= průměrná fitness populace) a
\[
s \;=\; (0{,}118;\; 0{,}471;\; 0{,}176;\; 0{,}235).
\]
Podíl nejlepšího řetězce $11$ vzrostl z~$12{,}5\,\%$ na $23{,}5\,\%$, podíl
nejhoršího $00$ klesl z~$25\,\%$ na $11{,}8\,\%$ --- přesně o~faktor
$f(i)/\bar f$, jak víme z~věty o~schématech.

\paragraph{Část $\mathcal{M}$ --- míchání.} Zavedeme funkci
\[
r(i,j,k) \;=\; \Pr[\text{potomek } k \text{ vznikne z~rodičů } i, j],
\]
která v~sobě zahrnuje jak křížení (přes všechny body řezu), tak mutaci.
Pak
\[
\mathbb{E}\big[p_k(t+1)\big] \;=\; \sum_{i}\sum_{j} s_i(t)\, s_j(t)\, r(i,j,k).
\]
Jde tedy o~\emph{kvadratický} operátor na simplexu. Odvození $r(i,j,k)$ je
technicky náročné, ale možné. Klíčové zjednodušení plyne z~toho, že jak
jednobodové křížení, tak bitová mutace \uv{nevidí} absolutní hodnoty bitů,
jen jejich shody a~neshody. Formálně: pro libovolné $k$ platí
\[
r(i,j,k) \;=\; r(i \oplus k,\; j \oplus k,\; 0),
\]
kde $\oplus$ je bitový XOR. Stačí tedy znát jedinou \emph{míchací matici}
$\mathbf{M}$ s~prvky $\mathbf{M}_{ij} = r(i,j,0)$ (rozměru $2^l \times 2^l$)
a~zbytek se dopočte permutacemi $\sigma_k : x \mapsto x \oplus k$:
$\mathcal{M}(s)_k = (\sigma_k s)^{\!\top}\, \mathbf{M}\, (\sigma_k s)$.
To je hlavní technický přínos Voseho formalismu --- z~$2^{3l}$ čísel se úloha
smrskne na $2^{2l}$.

\subsubsection{Výsledky pro nekonečné populace}

\begin{itemize}
\item SGA je \textbf{diskrétní dynamický systém} na simplexu; posloupnost
  $p(0), p(1), \dots$ je trajektorie tohoto systému.
\item \textbf{Pevné body $\mathcal{F}$}: populace tvořené pouze kopiemi
  nejlepšího jedince z~počáteční populace (selekce \uv{zaostřuje}).
\item \textbf{Pevný bod $\mathcal{M}$}: jediný --- rovnoměrné rozdělení
  (míchání \uv{rozostřuje}, maximalizuje entropii).
\item Pevné body složeného $\mathcal{G}$ nejsou obecně známy; jejich hledání
  je hlavní otevřený problém. Chování je kombinací zaostřování a~rozostřování,
  které se projevuje jako \textbf{přerušované rovnováhy}
  (\emph{punctuated ekvilibria}) --- dlouhá období metastability střídaná
  rychlými přechody. Tento jev je znám i~z~biologie.
\end{itemize}

\subsubsection{Konečné populace: Markovův řetězec}

\begin{defbox}[Markovův řetězec]
Stochastický proces v~diskrétním čase se stavy, kde pravděpodobnost přechodu
do dalšího stavu závisí \emph{jen} na aktuálním stavu (nikoli na historii).
Úplným popisem je matice přechodových pravděpodobností.
\end{defbox}

Modelujeme SGA s~populací velikosti $n$ nad řetězci délky $l$:
\begin{itemize}
\item \textbf{Stavy} = jednotlivé možné populace (multimnožiny $n$ řetězců).
  Jejich počet je
  \[
  N \;=\; \binom{n + 2^l - 1}{\,2^l - 1\,},
  \]
  což je počet multimnožin velikosti $n$ z~$2^l$ typů.
\item \textbf{Matice $\mathbf{Z}$}: sloupce odpovídají populacím,
  $\mathbf{Z}(y,i)$ = počet výskytů jedince $y$ v~populaci $i$. Sloupce
  matice $\mathbf{Z}$ jsou tedy stavy řetězce.
\item \textbf{Matice přechodů $\mathbf{Q}$} rozměru $N \times N$; lze ji
  explicitně odvodit z~$\mathcal{G}$ (přechod je multinomické losování
  $n$ jedinců podle rozdělení $\mathcal{G}(p)$):
  \[
  \mathbf{Q}_{i,j} \;=\; n! \prod_{y=0}^{2^l-1}
  \frac{\big(\mathcal{G}(p_i)_y\big)^{\mathbf{Z}(y,j)}}{\mathbf{Z}(y,j)!}.
  \]
\end{itemize}

\textbf{Problém velikosti.} Už pro $n = l = 8$ má $\mathbf{Q}$ řádově
$10^{29}$ prvků. Model je tedy exaktní, ale prakticky nepoužitelný pro
výpočet; jeho hodnota je teoretická.

\paragraph{Kvalitativní důsledky.}
\begin{itemize}
\item Je-li $p_M > 0$, je řetězec \textbf{ireducibilní a~aperiodický}
  (z~libovolné populace lze mutacemi dosáhnout libovolné jiné) $\Rightarrow$
  existuje jediné stacionární rozdělení a~SGA navštíví každou populaci
  (tedy i~tu s~globálním optimem) nekonečně mnohokrát s~pravděpodobností 1.
\item Bez elitismu ale \textbf{nekonverguje} k~optimu --- optimum se objeví
  a~zase zmizí. \emph{Elitistický} GA (uchovávající nejlepší nalezené řešení)
  konverguje k~globálnímu optimu s~pravděpodobností 1 při $t \to \infty$.
  To je standardní konvergenční výsledek; neříká však nic o~\emph{rychlosti}.
\item Existuje korespondence mezi modelem nekonečné a~konečné populace:
  pro $n \to \infty$ se chování konečného modelu blíží deterministické
  trajektorii $\mathcal{G}$.
\end{itemize}

\subsection{Shrnutí ke~zkoušce}

\begin{itemize}
\item Schéma = slovo nad $\{0,1,*\}$; $o(S)$ = počet pevných pozic;
  $d(S)$ = vzdálenost první a~poslední pevné pozice; $F(S)$ = průměrná fitness
  reprezentantů \emph{v~dané populaci}. Celkem $3^m$ schémat, jedinec je
  reprezentantem $2^m$ z~nich.
\item TST: $\mathbb{E}[C(S,t+1)] \ge C(S,t)\frac{F(S)}{\bar f}
  [1 - p_C \frac{d(S)}{m-1} - p_M o(S)]$; odvození ve třech krocích
  (selekce $\to$ exponenciální růst, křížení $\to$ $d(S)/(m-1)$,
  mutace $\to$ $(1-p_M)^{o(S)}$).
\item BBH: GA skládá krátká, nízkého řádu, nadprůměrná schémata (stavební
  bloky). Důsledky: záleží na kódování, na velikosti populace, škodí předčasná
  konvergence.
\item Implicitní paralelismus ($\sim n^3$ užitečně zpracovaných schémat),
  analogie s~$k$-rukým banditou (exponenciální alokace pokusů vítězi).
\item Kritika: nerovnost pro jeden krok, konečné populace, deceptivní úlohy,
  kolaterální konvergence, velký rozptyl fitness, ignorování konstruktivních
  efektů, závislost schémat.
\item Voseho model: řetězce ztotožníme s~čísly $0..2^l-1$; populaci popisuje
  vektor relativních četností $p(t)$ v~simplexu a~vektor selekčních
  pravděpodobností $s(t)$.
\item Hledáme $\mathcal{G} = \mathcal{M} \circ \mathcal{F}$ s~$p(t+1)=\mathcal{G}(p(t))$.
  $\mathcal{F}$ (selekce) je dána diagonální maticí fitness:
  $s = \mathbf{F}p / \sum_j \mathbf{F}_{jj}p_j$;
  $\mathcal{M}$ (křížení + mutace) je kvadratický operátor
  $\mathbb{E}[p_k'] = \sum_{i,j} s_i s_j\, r(i,j,k)$. Díky symetrii
  $r(i,j,k) = r(i\oplus k, j \oplus k, 0)$ stačí jediná míchací matice.
\item Pevné body: $\mathcal{F}$ --- populace kopií nejlepšího jedince;
  $\mathcal{M}$ --- rovnoměrné rozdělení. Kombinace vede k~přerušovaným
  rovnováhám. Pevné body $\mathcal{G}$ nejsou známy.
\item Konečné populace: Markovův řetězec se stavy = populace,
  $N = \binom{n+2^l-1}{2^l-1}$ stavů, matice $\mathbf{Q}$ je exaktní, ale
  astronomicky velká. Ireducibilita při $p_M>0$; elitistický GA konverguje
  k~optimu s~pravděpodobností 1.
\item Nekonečné populace = deterministický model (exaktní, nepraktický),
  konečné = stochastický (realistický, nezvládnutelný). Limitní přechod je
  spojuje.
\end{itemize}

%=====================================================================
\section{Evoluční strategie, diferenciální evoluce, koevoluce, otevřená evoluce}
\label{sec:es}
%=====================================================================

Kapitola spojuje čtyři témata, která mají společné to, že jdou \emph{za} klasický genetický
algoritmus: dvě metody pro spojitou optimalizaci (ES a~DE) a~dvě témata, kde se mění samotná
fitness --- koevoluce (fitness závisí na ostatních jedincích) a~otevřená evoluce (fitness
není vůbec pevně dána).

\subsection{Evoluční strategie}
\label{ssec:es}
\subsubsection{Motivace a~základní rysy}

Evoluční strategie (\emph{evolution strategies}, ES) navrhli Rechenberg
a~Schwefel v~Berlíně v~60.\ letech pro optimalizaci reálných parametrů
v~obtížných technických úlohách (tvar trysky, ohyb potrubí). Charakteristiky:

\begin{itemize}
\item Jedinec je vektor reálných čísel; \textbf{mutace je hlavní operátor}.
\item Jedinec obsahuje kromě \emph{genetických parametrů} (vlastní řešení)
  i~\emph{strategické, endogenní parametry} (\emph{endogenous parameters}),
  které řídí samotnou evoluci --- typicky směrodatné odchylky mutací.
  Odtud slogan \uv{evoluce evoluce}.
\item Rodičovská selekce je \textbf{náhodná} (nezávislá na fitness);
  environmentální selekce je \textbf{deterministická} --- vybírá se
  $\mu$ nejlepších.
\item Nejmodernějším a~nejúspěšnějším zástupcem je CMA-ES.
\end{itemize}

\begin{defbox}[Notace ES]
$\mu$ --- velikost populace (počet rodičů), $\lambda$ --- počet vytvořených
potomků, $\rho$ --- počet rodičů podílejících se na jednom potomkovi.
\begin{itemize}
\item $(\mu + \lambda)$-ES: nová populace se vybírá z~\emph{$\mu$ rodičů
  i~$\lambda$ potomků} (elitistická).
\item $(\mu, \lambda)$-ES: nová populace se vybírá \emph{jen z~$\lambda$
  potomků}; rodiče vždy vymřou (neelitistická), nutně $\lambda > \mu$.
\item $(\mu/\rho \,\overset{+}{,}\, \lambda)$-ES: explicitní zápis
  s~rekombinací $\rho$ rodičů.
\end{itemize}
\end{defbox}

\begin{center}
\begin{tabular}{@{}p{4cm}p{5.5cm}p{5.5cm}@{}}
\toprule
 & $(\mu,\lambda)$ & $(\mu+\lambda)$ \\
\midrule
pool pro selekci & jen potomci & rodiče + potomci \\
elitismus & ne (fitness může klesnout) & ano \\
únik z~lokálních optim & lepší & horší \\
konvergence & pomalejší & rychlejší \\
doporučení & spojité domény, zašuměná či dynamická fitness & diskrétní domény,
  drahá fitness \\
samoadaptace $\sigma$ & funguje dobře (špatné $\sigma$ vymře) & může uváznout
  (starý rodič s~malým $\sigma$ přežívá) \\
poměr & typicky $\lambda/\mu \approx 7$ & i~$\mu \ge \lambda$; $\lambda = 1$ =
  steady-state ES \\
\bottomrule
\end{tabular}
\end{center}

\subsubsection{$(1+1)$-ES a~pravidlo 1/5}

Nejjednodušší ES: populace je jediný jedinec, vytváří se jediný potomek
gaussovskou mutací a~přijímá se, jen pokud je lepší.

\begin{algorithm}[!ht]
\caption{$(1+1)$-ES s~pravidlem 1/5 (minimalizace $f:\mathbb{R}^n \to \mathbb{R}$)}
\begin{algorithmic}[1]
\State inicializuj $\vec x(0)$ náhodně, $\sigma > 0$, $t \gets 0$
\Repeat
  \State $\vec y(t) \gets \vec x(t) + \sigma \cdot N(\vec 0, \mathbf{I})$
  \If{$f(\vec y(t)) < f(\vec x(t))$} $\vec x(t+1) \gets \vec y(t)$
  \Else\ $\vec x(t+1) \gets \vec x(t)$
  \EndIf
  \State sleduj $p$ = podíl úspěšných mutací za posledních $k$ iterací
  \State každých $k$ iterací uprav $\sigma$:
  \State \quad \textbf{if} $p > 1/5$ \textbf{then} $\sigma \gets \sigma / c$
    \Comment{daří se: dělej větší kroky, více explorace}
  \State \quad \textbf{if} $p < 1/5$ \textbf{then} $\sigma \gets \sigma \cdot c$
    \Comment{nedaří se: zmenši kroky, více exploatace}
  \State \quad \textbf{if} $p = 1/5$ \textbf{then} $\sigma$ beze změny
    \Comment{typicky $0{,}8 \le c \le 1$}
  \State $t \gets t+1$
\Until{$\vec x(t)$ je dost dobré}
\end{algorithmic}
\end{algorithm}

\begin{defbox}[Pravidlo 1/5 úspěchu (\emph{1/5 success rule})]
Heuristika odvozená Rechenbergem z~analýzy dvou modelových funkcí (kulová
a~koridorová): \emph{optimální poměr úspěšných mutací je přibližně $1/5$.}
Je-li podíl úspěchů vyšší, je krok příliš malý (jsme příliš opatrní)
a~$\sigma$ se zvětší; je-li nižší, je krok příliš velký a~$\sigma$ se zmenší.
Jde o~\emph{adaptivní} (nikoli samoadaptivní) řízení parametru --- používá
zpětnou vazbu z~běhu, ale parametr není součástí genomu.
\end{defbox}

\subsubsection{Samoadaptace strategických parametrů}

\begin{defbox}[Samoadaptace (\emph{self-adaptation})]
Strategické parametry (odchylky $\sigma$, případně úhly natočení) jsou
\textbf{součástí genomu} a~podléhají mutaci a~selekci spolu s~řešením.
Selekce je nepřímá: jedinci s~vhodným $\sigma$ produkují lepší potomky,
a~tím se jejich $\sigma$ šíří populací. Pořadí je zásadní: \textbf{nejprve
se mutují strategické parametry, teprve pak se s~nimi mutují genetické
parametry} --- jinak by se hodnotilo staré $\sigma$.
\end{defbox}

Jedinec má tvar $C = [\vec x, \vec \sigma]$, kde $\vec x \in \mathbb{R}^n$
a~$\vec\sigma$ má $k$ složek:

\begin{center}
\begin{tabular}{@{}clp{8.2cm}@{}}
\toprule
$k$ & název & geometrie mutace \\
\midrule
$1$ & jedno globální $\sigma$ & izotropní --- vrstevnice hustoty jsou koule \\
$n$ & nekorelované mutace & elipsoid s~osami rovnoběžnými se souřadnicemi \\
$n(n+1)/2$ & korelované mutace & obecný elipsoid ($n$ odchylek $+$ $n(n-1)/2$
  úhlů natočení $\alpha_{ij}$), odpovídá plné kovarianční matici \\
\bottomrule
\end{tabular}
\end{center}

Kovarianční matice $\mathbf{C}$ je součinem rotační matice: $c_{ii} = \sigma_i^2$,
$c_{ij} = \tfrac12(\sigma_i^2 - \sigma_j^2)\tan(2\alpha_{ij})$.

\paragraph{Standardní pravidla mutace (Schwefel).}
\[
\sigma_i' = \sigma_i \cdot \exp\big(\tau' N(0,1) + \tau N_i(0,1)\big),
\qquad
x_i' = x_i + \sigma_i' \cdot N_i(0,1),
\]
kde $\tau' \propto 1/\sqrt{2n}$ (společná porucha pro celý jedinec)
a~$\tau \propto 1/\sqrt{2\sqrt{n}}$ (individuální porucha složky).
Lognormální tvar zajišťuje kladnost $\sigma$ a~symetrii zvětšení/zmenšení.
Zavádí se dolní mez $\sigma_{\min}$, aby $\sigma$ nezdegenerovalo na nulu.
Úhly se mutují aditivně: $\alpha_{ij}' = \alpha_{ij} + \beta \cdot N(0,1)$.

\subsubsection{Rekombinace v~ES}

Z~$\rho$ rodičů vzniká \emph{jeden} potomek. Podle počtu rodičů rozlišujeme
\textbf{lokální} ($\rho = 2$) a~\textbf{globální} ($\rho = \mu$) rekombinaci;
podle způsobu kombinace:
\begin{itemize}
\item \textbf{Diskrétní (dominantní):} hodnota na dané pozici se náhodně vybere
  od jednoho z~rodičů.
\item \textbf{Intermediární (aritmetická):} hodnota je průměrem hodnot rodičů
  na dané pozici.
\end{itemize}
Obvyklá kombinace: diskrétní rekombinace pro genetické parametry
(zachovává rozmanitost) a~intermediární pro strategické parametry
(stabilizuje adaptaci $\sigma$).

\begin{algorithm}[!ht]
\caption{Obecný cyklus ES}
\begin{algorithmic}[1]
\State $t \gets 0$; inicializuj náhodně populaci $P_t$ o~$\mu$ jedincích $[\vec x, \vec\sigma]$
\State ohodnoť jedince v~$P_t$
\While{není splněna ukončovací podmínka}
  \State $C_{t+1} \gets \emptyset$
  \For{$\lambda$ krát}
    \State vyber náhodně $\rho$ rodičů z~$P_t$
    \State rekombinuj je do jednoho potomka $[\vec x, \vec\sigma]$
    \State mutuj strategické parametry $\vec\sigma$
    \State mutuj genetické parametry $\vec x$ pomocí nových $\vec\sigma$
    \State ohodnoť potomka a~vlož do $C_{t+1}$
  \EndFor
  \State $(\mu,\lambda)$: $P_{t+1} \gets$ $\mu$ nejlepších z~$C_{t+1}$
  \State $(\mu+\lambda)$: $P_{t+1} \gets$ $\mu$ nejlepších z~$P_t \cup C_{t+1}$
  \State $t \gets t+1$
\EndWhile
\end{algorithmic}
\end{algorithm}

\subsubsection{CMA-ES}

\noindent\mimoq{slidy uvádějí jednou větou jako \uv{dnes nejúspěšnější a~nejsložitější} ES}

\begin{defbox}[CMA-ES --- princip]
\emph{Covariance Matrix Adaptation ES} (Hansen, Ostermeier) je dnes referenční algoritmus
pro spojitou black-box optimalizaci. Místo samoadaptace jednotlivých $\sigma$ v~genomu
udržuje \textbf{jediné rozdělení pro celou populaci}
$\mathcal{N}(\vec m, \sigma^2 \mathbf{C})$ a~učí se jeho kovarianční matici z~historie
úspěšných kroků.

Iterace: navzorkuj $\lambda$ bodů, vyber $\mu$ nejlepších \emph{podle pořadí}, posuň střed
$\vec m$ do jejich váženého průměru a~aktualizuj $\mathbf{C}$ a~délku kroku $\sigma$ podle
\emph{evolučních cest} --- exponenciálně vyhlazených záznamů o~tom, kudy se střed pohyboval.
Jdou-li kroky souhlasným směrem, $\sigma$ se zvětší; jsou-li protichůdné, zmenší.
\end{defbox}

Výsledkem je algoritmus, který se automaticky přizpůsobí anizotropii a~korelacím úlohy
(v~podstatě se učí aproximaci inverzní Hessovy matice bez gradientu), je invariantní vůči
monotónním transformacím fitness a~má velmi málo laditelných parametrů.

\subsection{Diferenciální evoluce}
\label{sec:de}
\subsubsection{Motivace}

Diferenciální evoluce (\emph{differential evolution}, DE; Storn a~Price, 1997)
je jednoduchý a~překvapivě silný algoritmus pro black-box optimalizaci funkcí
$f: \mathbb{R}^D \to \mathbb{R}$. Klíčová myšlenka je \textbf{geometrická}:
velikost i~směr mutačního kroku se neodhaduje z~parametru $\sigma$, ale
přebírá se z~\emph{rozdílů dvojic jedinců v~populaci}. Populace tak sama
kóduje měřítko a~orientaci úlohy: na začátku je rozptýlená (velké kroky),
v~konvergujícím stavu jsou rozdíly malé (jemné doladění). Jde o~jakousi
implicitní samoadaptaci zdarma.

Aplikace sahají od strojového učení po plánování optimálních trajektorií
kosmických sond (Evropská kosmická agentura).

\subsubsection{Algoritmus DE/rand/1/bin}

Populace je $N$ vektorů $\vec x_{1},\dots,\vec x_{N} \in \mathbb{R}^D$.
Pro \textbf{každého} jedince $\vec x_i$ (rodičovská selekce je tedy triviální ---
rodičem je postupně každý) se provede:

\begin{enumerate}
\item \textbf{Mutace.} Vyber tři navzájem různé jedince
  $\vec x_{a}, \vec x_{b}, \vec x_{c}$ různé od $\vec x_i$ a~vytvoř
  \emph{dárce} (\emph{donor vector})
  \[
  \vec v_i \;=\; \vec x_{a} \;+\; F\,(\vec x_{b} - \vec x_{c}),
  \qquad F \in [0,2],\ \text{typicky } F \approx 0{,}8 .
  \]
\item \textbf{Křížení (binomické, uniformní \uv{s~pojistkou}).}
  Vytvoř \emph{zkušební vektor} (\emph{trial vector}) $\vec u_i$:
  \[
  u_{j,i} =
  \begin{cases}
  v_{j,i} & \text{pokud } \mathrm{rand}_j \le C \text{ nebo } j = I_{\mathrm{rand}},\\[2pt]
  x_{j,i} & \text{jinak},
  \end{cases}
  \]
  kde $\mathrm{rand}_j \sim U(0,1)$, $C \in [0,1]$ je práh křížení
  a~$I_{\mathrm{rand}}$ je náhodný index z~$\{1,\dots,D\}$. Podmínka
  $j = I_{\mathrm{rand}}$ je \uv{pojistka}: zaručuje, že alespoň jedna složka
  pochází od dárce, takže potomek není nikdy identickou kopií rodiče.
\item \textbf{Environmentální selekce (souboj 1:1).}
  \[
  \vec x_i(t+1) =
  \begin{cases}
  \vec u_i & \text{pokud } f(\vec u_i) \le f(\vec x_i(t)),\\
  \vec x_i(t) & \text{jinak.}
  \end{cases}
  \]
  Selekce je tedy elitistická na úrovni jednotlivce --- populace se nikdy
  nezhoršuje.
\end{enumerate}

\begin{algorithm}[!ht]
\caption{DE/rand/1/bin}
\begin{algorithmic}[1]
\Require velikost populace $N$, měřítko $F$, práh křížení $C$
\State inicializuj populaci $\{\vec x_i\}_{i=1}^{N}$ náhodně v~mezích úlohy
\While{není splněna ukončovací podmínka}
  \For{$i \gets 1$ \textbf{to} $N$}
    \State vyber náhodně různé indexy $a,b,c \ne i$
    \State $\vec v \gets \vec x_a + F(\vec x_b - \vec x_c)$
    \State $I_{\mathrm{rand}} \gets$ náhodný index z~$\{1,\dots,D\}$
    \For{$j \gets 1$ \textbf{to} $D$}
      \State $u_j \gets v_j$ pokud $\mathrm{rand}() \le C$ nebo $j = I_{\mathrm{rand}}$, jinak $u_j \gets x_{j,i}$
    \EndFor
    \If{$f(\vec u) \le f(\vec x_i)$} $\vec x_i' \gets \vec u$ \Else\ $\vec x_i' \gets \vec x_i$ \EndIf
  \EndFor
  \State $\{\vec x_i\} \gets \{\vec x_i'\}$
\EndWhile
\end{algorithmic}
\end{algorithm}

\subsubsection{Číselný příklad jednoho kroku}

Nechť $D = 3$, minimalizujeme $f(\vec x) = \sum_j x_j^2$, parametry
$F = 0{,}8$, $C = 0{,}9$. Aktuální jedinec
$\vec x_i = (4;\, -3;\, 2)$, $f(\vec x_i) = 16+9+4 = 29$.
Náhodně vybraní jedinci:
\[
\vec x_a = (2;\, -1;\, 0{,}5), \quad
\vec x_b = (1;\, 0;\, -1), \quad
\vec x_c = (-1;\, 2;\, 1).
\]

\textbf{Mutace:}
\[
\vec x_b - \vec x_c = (2;\, -2;\, -2), \qquad
\vec v = (2;\,-1;\,0{,}5) + 0{,}8\,(2;\,-2;\,-2) = (3{,}6;\; -2{,}6;\; -1{,}1).
\]

\textbf{Křížení:} nechť $I_{\mathrm{rand}} = 1$ a~náhodná čísla
$\mathrm{rand} = (0{,}40;\ 0{,}95;\ 0{,}20)$. Pak
\begin{itemize}
\item $j=1$: $0{,}40 \le 0{,}9$ (a~navíc $j = I_{\mathrm{rand}}$) $\Rightarrow u_1 = v_1 = 3{,}6$;
\item $j=2$: $0{,}95 > 0{,}9$ a~$j \ne I_{\mathrm{rand}}$ $\Rightarrow u_2 = x_{2,i} = -3$;
\item $j=3$: $0{,}20 \le 0{,}9$ $\Rightarrow u_3 = v_3 = -1{,}1$.
\end{itemize}
Zkušební vektor je $\vec u = (3{,}6;\, -3;\, -1{,}1)$.

\textbf{Selekce:} $f(\vec u) = 12{,}96 + 9 + 1{,}21 = 23{,}17 < 29 = f(\vec x_i)$,
takže $\vec u$ nahrazuje $\vec x_i$ v~nové populaci.

\subsubsection{Varianty mutace}

Zápis \texttt{DE/$x$/$y$/$z$}: $x$ = jak se volí základní vektor,
$y$ = počet rozdílových vektorů, $z$ = typ křížení (\texttt{bin} = binomické
uniformní, \texttt{exp} = exponenciální, souvislý úsek).

\begin{center}
\begin{tabular}{@{}ll@{}}
\toprule
\textbf{varianta} & \textbf{dárce $\vec v$} \\
\midrule
DE/rand/1 & $\vec x_a + F(\vec x_b - \vec x_c)$ \\
DE/rand/2 & $\vec x_a + F(\vec x_b - \vec x_c + \vec x_d - \vec x_e)$ \\
DE/best/1 & $\vec x_{\mathrm{best}} + F(\vec x_b - \vec x_c)$ \\
DE/best/2 & $\vec x_{\mathrm{best}} + F(\vec x_b - \vec x_c + \vec x_d - \vec x_e)$ \\
DE/current-to-best/1 & $\vec x_i + F(\vec x_{\mathrm{best}} - \vec x_i) + F(\vec x_b - \vec x_c)$ \\
\bottomrule
\end{tabular}
\end{center}

Varianty s~\texttt{best} konvergují rychleji, ale snáze uvíznou v~lokálním
optimu; \texttt{rand} je robustnější. Moderní varianty (jDE, SaDE, JADE,
SHADE, L-SHADE) adaptují $F$ a~$C$ za běhu a~patří ke špičce v~soutěžích
CEC.

\subsubsection{Srovnání DE s~ES a~GA}

\begin{center}
\begin{tabular}{@{}p{3.2cm}p{3.7cm}p{3.7cm}p{3.7cm}@{}}
\toprule
 & \textbf{GA (reálný)} & \textbf{ES} & \textbf{DE} \\
\midrule
zdroj kroku mutace & pevné $\sigma$ & evolvované $\sigma$ & rozdíly jedinců \\
rodičovská selekce & podle fitness & náhodná & každý jedinec jednou \\
env.\ selekce & generační / elitismus & $\mu$ nejlepších & souboj rodič vs.\ potomek \\
hlavní operátor & křížení & mutace & diferenční mutace \\
adaptace měřítka & žádná & explicitní (v~genomu) & implicitní (z~populace) \\
\bottomrule
\end{tabular}
\end{center}

\subsection{Koevoluce}
\label{sec:koevoluce}
\subsubsection{Kontextová fitness}

\begin{defbox}[Koevoluce (\emph{coevolution})]
Evoluční algoritmus, v~němž fitness jedince \textbf{není absolutní}, ale závisí
na ostatních jedincích --- na jeho vlastní populaci nebo na jiné, souběžně se
vyvíjející populaci. Fitness je tedy \emph{kontextová} a~fitness krajina je
\textbf{dynamická}: mění se tím, jak se mění protivníci či partneři.
\end{defbox}

Motivace: pro mnoho úloh neumíme napsat absolutní účelovou funkci
(\uv{jak dobrá je strategie pro dámu?}), ale umíme dva jedince \emph{porovnat}
(zahrát partii). Kromě toho koevoluce slibuje \emph{evoluční závody ve
zbrojení} --- protivníci se navzájem tlačí ke stále lepším řešením.
Typy koevoluce:

\begin{center}
\begin{tabular}{@{}p{3cm}p{6cm}p{6cm}@{}}
\toprule
 & \textbf{kompetitivní} & \textbf{kooperativní} \\
\midrule
vztah & fitness jednoho roste na úkor druhého & jedinci se doplňují do
  společného řešení \\
příklad & predátor--kořist, host--parazit, hráč--protihráč, řadicí sítě vs.\
  posloupnosti čísel & dekompozice problému na podproblémy, moduly a~plány
  (CoDeepNEAT), hlavní programy a~ADF \\
fitness & výsledek souboje (turnaje) & kvalita složeného řešení, na němž se
  jedinec podílel \\
riziko & cyklení, odpojení, přespecializace & kredit assignment, líní partneři \\
\bottomrule
\end{tabular}
\end{center}

\subsubsection{Kompetitivní koevoluce}

\paragraph{Klasický příklad: Hillisovy řadicí sítě (1990).} Populace řadicích
sítí koevolvuje proti populaci testovacích posloupností. Fitness sítě = podíl
správně seřazených testů; fitness testu = kolik sítí na něm selhalo. Parazitické
testy se specializují na slabiny sítí, sítě se musí zlepšovat. Hillis takto
nalezl síť pro 16 prvků s~\textbf{61} komparátory. \emph{Pozor na častý omyl:}
nejlepší známý \emph{lidský} návrh (Green, 1969) má 60 komparátorů, koevoluce
ho tedy nepřekonala. Podstatné je srovnání uvnitř experimentu --- týž genetický
algoritmus \emph{bez} koevolvujících parazitů dosáhl jen 65 komparátorů.
Koevoluce zde tedy prokazatelně zlepšila optimalizační schopnost EA a~přiblížila
se lidskému výsledku na jediný komparátor.

\paragraph{Predátor--kořist} v~evoluční robotice: roboti-lovci a~roboti-uprchlíci
se navzájem zdokonalují. \textbf{Host--parazit}: modely v~umělém životě
(Tierra), testování software.

\paragraph{Turnajové hodnocení.} Fitness se získává soubojem: každý s~každým
(drahé), náhodný vzorek protivníků (levné, zašuměné), \emph{hall of fame}
(souboj s~archivem historicky nejlepších --- brání zapomínání), nebo
\emph{turnaj se sdílenou fitness} (\emph{competitive fitness sharing}: za
poražení protivníka, kterého poráží málokdo, je vyšší odměna).

\subsubsection{Patologie koevoluce}

\begin{itemize}
\item \textbf{Cyklení} (\emph{cycling}, intransitivita): populace obíhají
  v~kruhu jako v~kámen--nůžky--papír; zdánlivý pokrok není skutečný.
  Dynamika koevoluce může být až deterministicky chaotická.
\item \textbf{Odpojení} (\emph{disengagement}): jedna populace zcela přeroste
  druhou, všechny souboje končí stejně, fitness ztrácí rozlišovací schopnost
  a~selekce přestane fungovat (gradient zmizí).
\item \textbf{Efekt Červené královny} (\emph{Red Queen effect}): všichni se
  zlepšují, ale relativní fitness zůstává konstantní --- z~naměřených hodnot
  nepoznáme, zda je pokrok skutečný. Řešení: měření proti pevné externí
  sadě protivníků (\emph{master tournament}).
\item \textbf{Průměrné stabilní stavy} (\emph{mediocre stable states}):
  populace se ustálí ve vzájemné rovnováze na nízké úrovni.
\item \textbf{Přespecializace} (\emph{overspecialization}, \emph{focusing}):
  jedinec se optimalizuje na aktuálního protivníka a~ztratí obecnost.
\item \textbf{Zapomínání} (\emph{forgetting}): schopnost porazit staré
  protivníky se ztrácí. Řešení: archivy (hall of fame, Pareto-koevoluční
  archivy, \emph{Nash memory}).
\end{itemize}

\subsubsection{Kooperativní koevoluce}

\begin{defbox}[CCGA (\emph{cooperative coevolutionary GA}, Potter \& De Jong)]
Problém se rozdělí na $k$ komponent (např.\ souřadnice, moduly, pravidla).
Pro každou komponentu se udržuje \textbf{samostatná populace}. Fitness jedince
z~populace $i$ se určí tak, že se \emph{spojí} se zástupci ostatních populací
(typicky s~jejich aktuálně nejlepšími jedinci, případně i~s~náhodnými)
a~ohodnotí se výsledné složené řešení.
\end{defbox}

Výhody: rozklad velkého prostoru na menší, přirozená paralelizace, vhodné pro
modulární úlohy. Problémy: \emph{přiřazení zásluh} (\emph{credit assignment}) ---
jak rozdělit kvalitu složeného řešení mezi komponenty; \emph{provázanost}
(silně interagující komponenty nelze dobře oddělit); \emph{relativní
přespecializace} (populace se přizpůsobí konkrétním partnerům).

Příklady v~praxi: koevoluce hlavních programů a~ADF v~GP,
CoDeepNEAT (populace modulů a~populace \uv{blueprintů}, odst.~\ref{sec:neuro}),
meta-lamarckovské memetické algoritmy (populace memů a~populace řešení,
kap.~\ref{sec:lokalni}).

\subsubsection{Evoluce kooperace a~vězňovo dilema}

Koevoluce úzce souvisí se základní otázkou evoluční biologie: \emph{jak může
vzniknout altruismus}, když z~definice snižuje fitness jedince?
Darwinismus je ve své podstatě soutěživý (\uv{přežití nejzdatnějšího}),
přesto je v~přírodě i~ve společnosti mnoho spolupráce. Teorie:
\begin{itemize}
\item \textbf{Skupinová selekce} --- evoluce působí i~na skupiny (Darwin);
  problém: jak vysvětlit podvodníky uvnitř skupiny.
\item \textbf{Příbuzenská selekce} (\emph{kin selection}) --- zachování téměř
  identických genů u~příbuzných; problém: altruismus vůči cizím a~jiným druhům.
\item \textbf{Sobecký gen} (Dawkins) --- jednotkou evoluce je gen, ne jedinec.
  (Wilson: \uv{organismus je jen způsob, jak DNA vyrábí další DNA}.)
\item \textbf{Reciproční altruismus} (Trivers, 1971) --- vzájemný prospěch,
  \emph{stín budoucnosti}; formální model = iterované vězňovo dilema.
\end{itemize}

\begin{defbox}[Vězňovo dilema (\emph{prisoner's dilemma})]
Hra dvou hráčů s~akcemi $C$ (kooperace) a~$D$ (zrada) a~výplatami
\[
\begin{array}{c|cc}
 & C & D\\\hline
C & R/R & S/T\\
D & T/S & P/P
\end{array}
\qquad
T > R > P > S \quad (\text{a pro iterovanou verzi } 2R > T+S).
\]
Typicky $R=-1$, $T=0$, $P=-2$, $S=-3$. Protože $T > R$ a~$P > S$, je zrada
\textbf{dominantní strategií} obou hráčů; $DD$ je \textbf{Nashovo
ekvilibrium}, ale jako jediný výsledek \textbf{není Pareto-optimální} ---
$CC$ maximalizuje společný zisk.
\end{defbox}

Racionální hráči tedy dospějí k~výsledku, který je pro oba horší než dosažitelná
spolupráce. Tentýž mechanismus v~$n$-hráčské podobě se nazývá \emph{tragédie
obecní pastviny} (\emph{tragedy of the commons}): každému se individuálně
vyplatí čerpat sdílený zdroj o~něco víc, a~zdroj proto zanikne. Otázka
\uv{co je vlastně racionální (a~chovají se tak lidé?)} je jádrem celé
diskuse kolem vězňova dilematu.

\begin{defbox}[Dominantní strategie, Nashovo ekvilibrium, Pareto optimalita]
Strategie $s$ je \emph{dominantní} pro hráče $i$, dává-li stejný nebo lepší
výsledek než kterákoli jiná strategie $i$ proti \emph{všem} strategiím
protihráče. Dvojice $(s_i, s_j)$ je v~\emph{Nashově ekvilibriu}, jsou-li
vzájemně nejlepšími odpověďmi (žádný hráč nemá důvod jednostranně změnit
strategii). Řešení je \emph{Pareto-optimální}, pokud nelze zlepšit výsledek
jednoho hráče, aniž by se zhoršil výsledek jiného.
Nashova věta: každá hra s~konečně mnoha strategiemi má ekvilibrium
ve \emph{smíšených} strategiích (náhodná volba mezi čistými) ---
např.\ kámen--nůžky--papír.
\end{defbox}

\paragraph{Iterovaná verze a~Axelrodovy turnaje.} Hraje-li se opakovaně
a~hráči si pamatují historii, může se vyplatit spolupráce (\emph{stín
budoucnosti}). Zná-li se předem přesný počet kol $N$, indukcí od konce plyne,
že optimální je stále zrazovat --- v~praxi ale hráči konec neznají.
Axelrod (1980) uspořádal turnaje strategií:
\begin{itemize}
\item 1.\ turnaj: 14 strategií + RANDOM, 200 her, každý s~každým (i~sám se
  sebou), 5 opakování. Vyhrála \textbf{TFT} (\emph{tit for tat}): začni
  kooperací, pak opakuj poslední tah soupeře.
\item 2.\ turnaj: 62 strategií, všichni znali výsledky prvního --- TFT
  vyhrála znovu.
\item 3.\ \uv{ekologický} turnaj: simulace jako v~GA --- počet jedinců
  strategie v~další generaci je úměrný počtu vítězství, 1000 generací.
  TFT opět zvítězila.
\end{itemize}

\paragraph{Čtyři vlastnosti úspěšných strategií:}
\emph{slušnost} (niceness --- nezrazuj první),
\emph{dráždivost} (provocability/retaliation --- zradu potrestej okamžitě),
\emph{odpouštění} (forgiveness --- ale rychle se uklidni),
\emph{nezávistivost} (non-enviousness --- neusiluj o~to porazit \emph{konkrétního}
soupeře; TFT nikdy nezískala víc bodů než její protihráč, a~přesto vyhrála turnaj).
Jako pátou vlastnost Axelrod uvádí \emph{srozumitelnost} (clarity --- buď jednoduchý,
aby s~tebou ostatní vůbec dokázali navázat spolupráci).
Žádná strategie nevyhraje proti všem; je nutné být úspěšný proti velmi
rozmanitým soupeřům (ALL-D, TFTT, RANDOM, TRIGGER) a~umět hrát dobře sám
se sebou.

\paragraph{Poučení a~pozdější vývoj.} V~prostředích, kde výplaty přejí
spolupráci a~kde je vysoká pravděpodobnost opakování, se spolupráce obvykle
vyvine --- a~není k~tomu potřeba racionalita ani vědomí, stačí vyšší fitness.
V~zašuměném prostředí (chybné provedení tahu) je úspěšnější strategie
\textbf{Pavlov} (\emph{win-stay, lose-shift}): dostal-li jsem $R$ nebo $P$,
zopakuj tah ($C$); dostal-li jsem $T$ nebo $S$, změň jej ($D$).
Turnaj opakovaný po 20 letech vyhrál \emph{tým} strategií z~University of
Southampton: prvních zhruba 10 tahů sloužilo k~rozpoznání spoluhráče,
poté všichni členové týmu obětavě pomáhali jedné \uv{otcovské} strategii
získat vysoké skóre.

\subsubsection{Diverzita, druhy a~niky}

Malý selekční tlak stačí k~tomu, aby diverzita populace vymizela. (Podobnou
úvahu rozvíjí Flegrova \emph{zamrzlá evoluce}: geneticky proměnlivý druh
se vůči selekci chová pružně jen krátce, poté \uv{zamrzne} a~na změny
prostředí přestane reagovat.) Mechanismy pro udržení diverzity (užitečné
jak pro koevoluci, tak pro dynamickou fitness a~vícekriteriální optimalizaci):
\begin{itemize}
\item \textbf{Fitness sharing:} fitness se dělí počtem podobných jedinců
  $f'(i) = f(i)/\sum_j sh(d(i,j))$, kde
  $sh(d) = 1-(d/\sigma_{\mathrm{share}})^{\alpha}$ pro $d < \sigma_{\mathrm{share}}$,
  jinak 0. Vytváří tlak na obsazení různých nik.
\item \textbf{Crowding:} nový jedinec nahrazuje \emph{nejpodobnějšího}
  jedince (deterministic crowding, RTS).
\item \textbf{Nenáhodné páření} (\emph{non-random mating}): křížení jen mezi
  podobnými jedinci (\emph{speciace}), nebo podle topologie --- jedinci žijí
  na 2D/3D mřížce a~kříží se jen se sousedy.
\item \textbf{Ostrovní model} (\emph{island model}): několik subpopulací
  s~občasnou migrací (podrobněji níže).
\item \textbf{Explicitní druhy} pomocí značkovacích bitů (\emph{tag bits}),
  nebo podle strukturální podobnosti genomů (NEAT, odst.~\ref{sec:neuro});
  páření je pak povoleno jen uvnitř druhu.
\end{itemize}
Problémem je vždy definice podobnosti (Hammingova vzdálenost, vzdálenost
v~prostoru fenotypů, počet shodných inovačních čísel) a~volba prahu niky.

\paragraph{Paralelní modely EA.} Struktura populace úzce souvisí
s~paralelizací; rozlišují se tři modely:
\begin{itemize}
\item \textbf{Master--slave} (globální paralelizace): jedna populace, mezi
  uzly se rozdělí jen \emph{vyhodnocení fitness}. Algoritmus je matematicky
  totožný se sekvenčním; vyplatí se při drahé fitness.
\item \textbf{Ostrovní / hrubozrnný model} (\emph{coarse-grained}): několik
  nezávislých subpopulací, mezi nimiž občas \emph{migruje} několik jedinců.
  Parametry: topologie ostrovů, \emph{migrační interval} a~\emph{migrační
  míra}, výběr migrantů a~nahrazovaných. Vzácná migrace udržuje diverzitu
  (ostrovy konvergují jinam), příliš častá model degeneruje na jedinou populaci.
\item \textbf{Buněčný / jemnozrnný model} (\emph{fine-grained}, difuzní):
  jedinci sedí na 2D/3D mřížce a~kříží se jen se sousedy. Dobré řešení se šíří
  populací jako vlna, což silně zpomaluje předčasnou konvergenci; jde
  o~implicitní formu nenáhodného páření.
\end{itemize}
Ostrovní i~buněčný model tedy nejsou jen implementační trik --- mění samotnou
dynamiku algoritmu a~patří mezi standardní nástroje udržení diverzity.

\subsubsection{Dynamické fitness krajiny}

\begin{defbox}[Fitness krajina (\emph{fitness landscape})]
Zavedl Sewall Wright (1932). Grafické znázornění závislosti fitness na genotypu
(nebo na frekvenci alel či na rysu fenotypu). Teoretický nástroj pro
uvažování o~EA; ve vyšších dimenzích ovšem není intuitivní.
\end{defbox}

Fitness se často mění v~čase: robot v~místnosti, kde někdo rozsvítí; začne
pršet; krize na burze; turnaje agentů (koevoluce). Klasické GA si vedou špatně,
protože jsou \uv{přeučené} na statické úlohy --- rychle konvergují a~ztratí
diverzitu, kterou by pro reakci na změnu potřebovaly.
De Jongovo srovnání: $(1+10)$-ES reaguje na náhlou změnu pomalu, protože
samoadaptací zmenšila krok; generační GA s~ruletovou selekcí a~bez elitismu
si udrží více diverzity a~zotaví se rychleji.

Úpravy EA pro dynamické krajiny:
\begin{itemize}
\item \textbf{Diploidní reprezentace} (Goldberg, Smith) --- dvě sady
  chromozomů s~dominancí fungují jako dlouhodobá paměť při oscilující fitness.
\item \textbf{Hypermutace} (Cobb, Grefenstette) --- klesne-li průměrná fitness
  populace (pravděpodobně se změnilo prostředí), výrazně se zvýší
  pravděpodobnost mutace.
\item \textbf{Náhodná migrace} --- do populace se vloží dávka nových
  náhodných jedinců.
\item \textbf{Udržování diverzity} --- nižší selekční tlak, crowding, niching,
  subpopulace, čárková ES $(\mu,\lambda)$ (rodiče nepřežívají, populace
  se snáz odpoutá od zastaralého optima).
\end{itemize}
Typy změn: malé plynulé (opotřebení hardwaru), morfologické (vznikají
a~zanikají kopce), cyklické (roční období), katastrofické nespojité.
Mění-li se fitness příliš rychle, žádné řešení není trvale dobré
(srov.\ No Free Lunch).

\subsection{Otevřená evoluce}
\label{ssec:otevrena}

\noindent\mimoq{celý oddíl kromě interaktivní evoluce} Okruh otevřenou evoluci výslovně
jmenuje, slidy EVA~I ani EVA~II o~ní ale nemají nic; \emph{Evoluční robotika} ji zmiňuje
jedinou větou (\uv{evolúcia robotov je kontinuálny proces s~otvoreným koncom}).
Zpracováno proto stručně, na úrovni, která na zkoušce obstojí. Výjimkou je
\emph{interaktivní evoluce} na konci oddílu --- ta ve slidech Evoluční robotiky je
(subjektivní fitness, Biomorphs).

\begin{defbox}[Otevřená evoluce (\emph{open-ended evolution})]
Evoluční proces, který \textbf{nemá pevný cíl} a~který neustále produkuje novou složitost
a~nové \uv{druhy} chování, aniž by konvergoval. Vzorem je pozemská biosféra: nemá účelovou
funkci, přesto trvale generuje novinky.
\end{defbox}

\paragraph{Znaky otevřenosti.} Aby bylo možné o~systému vůbec rozhodnout, formulují se
\emph{znaky} (\emph{hallmarks}): \textbf{trvalá produkce novosti} (nové chování se objevuje
i~po libovolně dlouhém běhu); \textbf{neomezený růst složitosti} jedinců i~jejich vztahů;
\textbf{vznik nových tříd entit} a~úrovní organizace (v~biologii \emph{major transitions}:
replikátor $\to$ buňka $\to$ mnohobuněčnost $\to$ společenství); \textbf{evoluce
evolvability} --- mění se i~samotný způsob, jak systém generuje novinky.

\paragraph{Systémy a~proč se zasekávají.} \emph{Tierra} (Ray) a~\emph{Avida} ---
sebekopírující programy ve strojovém kódu, kde emergentně vznikli paraziti
a~hyperparaziti; \emph{Polyworld} --- agenti s~neuronovými sítěmi ve 2D světě;
\emph{ECHO} (Holland) --- agenti s~\uv{tagy} a~zdroji. Prakticky každý z~nich po čase
dospěje do stavu, kdy už se nic kvalitativně nového neobjevuje: prostředí je konečné,
prostor možných \uv{vynálezů} se vyčerpá a~dynamika zamrzne. Hypotézy proč: chybí
dostatečně bohaté a~samo se rozšiřující prostředí, chybí skutečně otevřená reprezentace
a~chybí mechanismus průběžně zvyšující nároky na jedince. Dosažení skutečně otevřené
evoluce zůstává otevřeným výzkumným problémem.

\paragraph{Generativní (nepřímé) reprezentace.} S~otevřenou evolucí souvisí otázka, jak
vůbec zakódovat neomezeně rostoucí složitost. Odpovědí jsou kódování, kde genotyp není
popis výsledku, ale \textbf{návod, jak výsledek vypěstovat}: \emph{L-systémy}
(přepisovací gramatiky generující větvící se struktury), \emph{celulární automaty},
\emph{buněčné kódování} (Gruau) a~\emph{CPPN/HyperNEAT} pro neuronové sítě
(odst.~\ref{sec:neuro}). Společný přínos: malý genom, velký fenotyp, přirozené vyjádření
symetrií a~opakujících se motivů.

\begin{defbox}[Novelty search a~řídkost]
Radikální myšlenka (Lehman \& Stanley, 2011: \emph{Abandoning Objectives}): u~klamavých
(deceptivních) úloh je cílová funkce zavádějící --- vede do lokálních optim. Proto
\textbf{zahoďme cíl a~odměňujme pouze novost}.

Definuje se \emph{prostor chování} (např.\ koncová pozice robota v~bludišti) a~novost
jedince $x$ se měří řídkostí okolí
\[
\rho(x) \;=\; \frac{1}{k}\sum_{i=1}^{k} \mathrm{dist}\big(x, \mu_i\big),
\]
kde $\mu_i$ je $k$ nejbližších sousedů $x$ v~aktuální populaci \textbf{a~v~permanentním
archivu}. Fitness $=\rho$; původní účelová funkce se \emph{nepoužívá vůbec}. Na úlohách
typu \uv{robot v~bludišti} bývá úspěšnější než cílená optimalizace. Navazuje
\textbf{quality-diversity} (MAP-Elites): mřížka buněk indexovaná deskriptory chování,
v~každé buňce nejlepší nalezené řešení.
\end{defbox}

\paragraph{Interaktivní evoluce.} Fitness funkci nahradí \textbf{člověk}: uživateli se
předloží populace kandidátů a~on vybere ty, které se mu líbí. Používá se tam, kde kritérium
kvality neumíme zapsat (estetika, design, hudba). Ukázky: Dawkinsovy \emph{Biomorphs},
\textbf{Picbreeder} a~\textbf{EndlessForms} (evoluce obrázků a~3D tvarů pomocí CPPN).
Omezením je pomalost a~únava uživatele --- nutné malé populace (9--20 jedinců na obrazovku)
a~málo generací. Právě zkušenost z~Picbreederu (nejzajímavější objekty nikdy nevznikly
cíleným hledáním) byla přímou motivací pro novelty search.

\subsection{Shrnutí ke~zkoušce}

\begin{itemize}
\item ES: reálné vektory, mutace hlavní operátor, náhodná rodičovská selekce,
  deterministická environmentální selekce; jedinec $= [\vec x, \vec\sigma]$.
\item $(\mu,\lambda)$ vs.\ $(\mu+\lambda)$: umět vyjmenovat rozdíly
  (elitismus, konvergence, lokální optima, samoadaptace).
\item $(1+1)$-ES a~pravidlo 1/5: $p > 1/5 \Rightarrow$ zvětši $\sigma$,
  $p < 1/5 \Rightarrow$ zmenši; jde o~adaptivní, nikoli samoadaptivní řízení.
\item Samoadaptace: $\sigma' = \sigma \exp(\tau'N(0,1) + \tau N_i(0,1))$,
  poté $x' = x + \sigma' N(0,1)$; \textbf{nejdřív $\sigma$, pak $x$}.
  Počet strategických parametrů: 1 (izotropní), $n$ (nekorelované),
  $n(n+1)/2$ (korelované, s~úhly natočení).
\item Rekombinace: lokální ($\rho=2$) / globální ($\rho=\mu$),
  diskrétní (dominantní) / intermediární (aritmetická).
\item CMA-ES: $\theta = (\vec m,\sigma,\mathbf C,\vec p_\sigma,\vec p_c)$,
  vzorkování $\mathcal N(\vec m, \sigma^2\mathbf C)$, evoluční cesty,
  rank-1 a~rank-$\mu$ update, řízení kroku podle délky $\vec p_\sigma$;
  invariantní vůči monotónním transformacím fitness a~lineárním transformacím
  prostoru.
\item DE/rand/1/bin: dárce $\vec v = \vec x_a + F(\vec x_b - \vec x_c)$;
  binomické křížení s~prahem $C$ a~pojistkou $j = I_{\mathrm{rand}}$;
  souboj potomka s~vlastním rodičem (nahrazení jen při zlepšení).
\item Parametry: $F \in [0,2]$ (typ.\ $0{,}5$--$0{,}9$), $C \in [0,1]$,
  $N \approx 10D$.
\item Geometrická podstata: měřítko a~orientace kroků se přebírají z~aktuálního
  rozptylu populace $\Rightarrow$ automatická adaptace bez strategických
  parametrů.
\item Varianty: rand/1, rand/2, best/1, best/2, current-to-best;
  \texttt{bin} vs.\ \texttt{exp} křížení.
\item Umět předvést jeden krok s~konkrétními čísly (mutace $\to$ křížení
  $\to$ selekce).
\item Koevoluce = kontextová fitness + dynamická krajina.
  \emph{Kompetitivní}: predátor--kořist, host--parazit, Hillisovy řadicí
  sítě. \emph{Kooperativní}: dekompozice na komponenty, CCGA,
  moduly $+$ blueprinty.
\item Patologie: cyklení (in\-tran\-zi\-ti\-vi\-ta), odpojení, efekt Červené
  královny, průměrné stabilní stavy, přespecializace, zapomínání.
  Léky: archivy (hall of fame), sdílená fitness, měření proti pevné sadě.
\item Vězňovo dilema: $T>R>P>S$, $2R > T+S$; $D$ dominantní, $DD$ Nashovo
  ekvilibrium a~jediné Pareto-neoptimální; iterovaná verze $\Rightarrow$ TFT;
  čtyři vlastnosti úspěšné strategie (slušnost, dráždivost, odpouštění,
  nezávistivost) $+$ srozumitelnost jako pátá; Pavlov v~zašuměném prostředí.
\item Diverzita: fitness sharing, crowding, speciace, ostrovní model,
  nenáhodné páření.
\item Dynamická fitness: diploidie, hypermutace, náhodná migrace, udržování
  diverzity, $(\mu,\lambda)$-ES.
\item Otevřená evoluce: bez pevného cíle; znaky = trvalá novost, neomezený růst
  složitosti, nové třídy entit, evoluce evolvability. Systémy Tierra, Avida,
  Polyworld, ECHO --- všechny se dříve či později zaseknou.
  Generativní (nepřímé) reprezentace: L-systémy, celulární automaty,
  buněčné kódování, CPPN.
\item Novelty search měří řídkost $\rho(x)$ v~prostoru chování vůči populaci
  \emph{a} permanentnímu archivu, účelovou funkci vůbec nepoužívá;
  navazuje quality-diversity a~MAP-Elites (mřížka deskriptorů chování).
\item Interaktivní evoluce: fitness = člověk (Biomorphs, Picbreeder,
  EndlessForms); omezením je malá populace a~únava uživatele.
\end{itemize}

%=====================================================================
\section{Rojové optimalizační algoritmy}
\label{sec:roje}
%=====================================================================

\begin{defbox}[Rojová inteligence (\emph{swarm intelligence})]
Přístup, kde globálně inteligentní chování \textbf{emerguje} z~jednoduchých
lokálních pravidel mnoha jedinců bez centrálního řízení. Rysy: decentralizace,
jednodušší jedinci, nepřímá komunikace přes prostředí (\emph{stigmergie}),
robustnost vůči výpadku jedince. Na rozdíl od EA zde obvykle není křížení,
mutace ani \uv{smrt} jedinců --- jedinci se pohybují a~pamatují si.
\end{defbox}

\subsection{Optimalizace rojem částic (PSO)}

\emph{Particle swarm optimization} (Eberhart a~Kennedy, 1995) je inspirována
hejny ptáků a~ryb. Jedinec je \emph{částice} --- bod v~$\mathbb{R}^n$
s~rychlostí. Nepoužívá se křížení ani mutace; místo toho má algoritmus
\textbf{paměť}:
\begin{itemize}
\item $\vec p_i$ (\emph{pBest}) --- nejlepší pozice, kterou částice $i$
  kdy navštívila (kognitivní, individuální paměť),
\item $\vec g$ (\emph{gBest}) --- nejlepší pozice nalezená celým rojem
  (sociální, kolektivní paměť).
\end{itemize}

\begin{defbox}[Pohybové rovnice PSO]
\[
\vec v_i \;\gets\; \underbrace{w\,\vec v_i}_{\text{setrvačnost}}
\;+\; \underbrace{c_1 r_1 (\vec p_i - \vec x_i)}_{\text{kognitivní složka}}
\;+\; \underbrace{c_2 r_2 (\vec g - \vec x_i)}_{\text{sociální složka}},
\qquad
\vec x_i \;\gets\; \vec x_i + \vec v_i ,
\]
kde $r_1, r_2 \sim U(0,1)$ (losují se pro každou složku zvlášť),
$c_1, c_2$ jsou koeficienty učení (v~původní verzi $c_1 = c_2 = 2$)
a~$w$ je \emph{setrvačná váha} (\emph{inertia weight}).
V~původní verzi z~roku 1995 je $w = 1$ (chybí); Shi a~Eberhart ji doplnili
a~doporučují lineární pokles např.\ z~$0{,}9$ na $0{,}4$: velké $w$ podporuje
exploraci, malé exploataci.
\end{defbox}

\begin{algorithm}[!ht]
\caption{PSO}
\begin{algorithmic}[1]
\State inicializuj pozice $\vec x_i$ a~rychlosti $\vec v_i$ náhodně; $\vec p_i \gets \vec x_i$
\Repeat
  \For{každou částici $i$}
    \State vyhodnoť $f(\vec x_i)$
    \If{$f(\vec x_i)$ je lepší než $f(\vec p_i)$} $\vec p_i \gets \vec x_i$ \EndIf
  \EndFor
  \State $\vec g \gets$ nejlepší z~$\{\vec p_i\}$
  \For{každou částici $i$}
    \State aktualizuj rychlost a~pozici dle pohybových rovnic
  \EndFor
\Until{je dosaženo maxima iterací nebo požadované přesnosti}
\end{algorithmic}
\end{algorithm}

\paragraph{Číselný příklad.} Částice v~$\mathbb{R}^2$:
$\vec x = (1;\,2)$, $\vec v = (0{,}5;\,-0{,}5)$, $\vec p = (0;\,3)$,
$\vec g = (-1;\,1)$, $w = 0{,}7$, $c_1 = c_2 = 1{,}5$, $r_1 = 0{,}4$,
$r_2 = 0{,}8$.
\begin{align*}
\vec v' &= 0{,}7\,(0{,}5;\,-0{,}5) + 1{,}5\cdot 0{,}4\,\big((0;3)-(1;2)\big)
  + 1{,}5\cdot 0{,}8\,\big((-1;1)-(1;2)\big)\\
 &= (0{,}35;\,-0{,}35) + 0{,}6\,(-1;\,1) + 1{,}2\,(-2;\,-1)
  = (-2{,}65;\; -0{,}95),\\
\vec x' &= (1;\,2) + (-2{,}65;\,-0{,}95) = (-1{,}65;\; 1{,}05).
\end{align*}
Částice tedy přeletěla přes $\vec g$ --- typické chování PSO, které je zdrojem
explorace. Proto se rychlost omezuje ($|v_j| \le v_{\max}$) nebo se používá
\emph{constriction faktor} (Clerc a~Kennedy):
$\vec v \gets \chi[\vec v + c_1 r_1(\vec p - \vec x) + c_2 r_2 (\vec g - \vec x)]$,
kde $\chi = 2/|2 - \varphi - \sqrt{\varphi^2 - 4\varphi}|$ pro
$\varphi = c_1 + c_2 > 4$; pro $c_1=c_2=2{,}05$ vychází $\chi \approx 0{,}729$.
Tato volba zaručuje konvergenci roje.

\paragraph{Topologie sousedství.} Místo globálního $\vec g$ (topologie
\emph{gbest}, tj.\ úplný graf) lze použít \emph{lbest} --- nejlepší jedinec
v~lokálním okolí:
\begin{itemize}
\item \textbf{Kruh} (každá částice má $k$ sousedů): informace se šíří pomalu
  $\Rightarrow$ pomalejší konvergence, ale lepší explorace a~menší riziko
  uvíznutí.
\item \textbf{Von Neumannova mřížka}, \textbf{hvězda}, \textbf{náhodné grafy}.
\item Empiricky: gbest je rychlejší na jednoduchých úlohách, lbest robustnější
  na multimodálních.
\end{itemize}

\paragraph{Varianty a~vztah k~EA.}
Diskrétní/binární PSO (rychlost interpretována jako pravděpodobnost překlopení
bitu přes sigmoidu), PSO s~více roji, hybridizace s~lokálním prohledáváním,
\emph{grammatical swarm}.

\begin{center}
\begin{tabular}{@{}p{5cm}p{5cm}p{5.2cm}@{}}
\toprule
 & \textbf{genetický algoritmus} & \textbf{PSO} \\
\midrule
společné & náhodná inicializace, fitness, stochastické aktualizace &
  totéž \\
jedinci & vznikají a~zanikají & trvale existují, pohybují se \\
paměť & jen v~populaci & explicitní ($\vec p_i$, $\vec g$) \\
operátory & křížení, mutace & žádné, jen pohyb \\
tok informace & obousměrný mezi rodiči & jednosměrný: od lepších k~ostatním \\
\bottomrule
\end{tabular}
\end{center}

\subsection{Optimalizace mravenčí kolonií (ACO)}

\noindent\mimoq{ve slidech EVA je z~rojových algoritmů jen PSO; okruh ale mluví
o~rojových algoritmech v~množném čísle a~ACO je jejich kanonický druhý zástupce}

\emph{Ant colony optimization} (M.~Dorigo, 1991) napodobuje hledání cesty
u~mravenců: mravenec pokládá \emph{feromonovou stopu}, po kratší cestě se
vrátí dřív, feromon se tam hromadí rychleji, což přitahuje další mravence ---
kladná zpětná vazba. Feromon se navíc \textbf{odpařuje}, takže systém
zapomíná zastaralou informaci a~neuvízne.

\begin{defbox}[Stigmergie]
Nepřímá komunikace prostřednictvím změn v~prostředí. Mravenci spolu
nekomunikují přímo; jediným kanálem je feromon na hranách. Obecný princip
reaktivních multiagentních architektur.
\end{defbox}

\paragraph{Ilustrace --- Steelsův průzkumník Marsu.} Roj robotů má sbírat
vzorky hornin; základna vysílá navigační signál (stačí sledovat gradient).
Robot má \uv{drobky} (radioaktivní značky) pro nepřímou komunikaci. Pravidla
s~prioritou:
R1: vidíš-li překážku, změň směr; R2: neseš-li vzorek a~jsi na základně,
polož ho; R3: neseš-li vzorek, jdi po gradientu nahoru; R4: vidíš-li vzorek,
seber ho; R5: jinak se pohybuj náhodně.
Druhá iterace přidá stigmergii: R7: neseš-li vzorek a~nejsi na základně,
polož 2 drobky a~jdi po gradientu nahoru; R8: cítíš-li drobek, seber 1 drobek
a~jdi po gradientu dolů. Vznikne tak efektivní kolektivní sběr
z~bohatých nalezišť, aniž by roboti měli mapu či komunikaci.

\paragraph{Obecné schéma ACO.} Úlohu převedeme na hledání cesty v~ohodnoceném
grafu; mravenci jsou agenti konstruující řešení po hranách.
\begin{enumerate}
\item Každý mravenec stochasticky zkonstruuje řešení (posloupnost hran).
\item Volitelně \emph{akce démona} (\emph{daemon actions}) --- centralizovaný
  krok, typicky lokální prohledávání zlepšující nalezené cesty.
\item Porovnají se kvality řešení.
\item Aktualizují se feromony na hranách (odpaření + uložení).
\end{enumerate}

\begin{defbox}[Ant System (AS) --- pravidlo přechodu]
Mravenec $k$ v~městě $i$ zvolí následující město $j$ z~množiny dosud
nenavštívených $\mathcal{N}_i^k$ \textbf{ruletou}: každému kandidátovi přiřadí
\emph{skóre} a~to znormalizuje na pravděpodobnost.
\[
w_{ij} \;=\; \big[\tau_{ij}\big]^{\alpha}\,\big[\eta_{ij}\big]^{\beta},
\qquad
p_{ij}^{k} \;=\; \frac{w_{ij}}{\displaystyle\sum_{l \in \mathcal{N}_i^k} w_{il}},
\qquad \eta_{ij} = \frac{1}{d_{ij}},
\]
kde $\tau_{ij}$ je množství feromonu na hraně, $\eta_{ij}$ heuristická
\uv{viditelnost} (převrácená vzdálenost) a~$\alpha, \beta$ váhy obou vlivů
(typicky $\alpha = 1$, $\beta \in [2,5]$).

\textbf{Zapamatovatelný tvar.} Indexy $i$ a~$k$ jsou během jednoho rozhodnutí
\emph{konstanty} ($i$ = kde stojím, $k$ = čí je seznam nenavštívených měst)
a~jmenovatel je jen normalizace. Dosadíme-li navíc $\eta = 1/d$, zbude
\[
p_j \;\propto\; \frac{\tau_j^{\,\alpha}}{d_j^{\,\beta}}
\qquad\text{---\ \uv{feromon nahoře, vzdálenost dole}.}
\]
Z~toho se plný zápis kdykoli zrekonstruuje: dopočítat skóre všem kandidátům
a~vydělit jejich součtem.
\end{defbox}

\begin{defbox}[Ant System --- aktualizace feromonu]
\[
\tau_{ij} \;\gets\; (1-\rho)\,\tau_{ij} \;+\; \sum_{k=1}^{m} \Delta\tau_{ij}^{k},
\qquad
\Delta\tau_{ij}^{k} =
\begin{cases}
Q/L_k & \text{prošel-li mravenec } k \text{ hranou } (i,j),\\
0 & \text{jinak,}
\end{cases}
\]
kde $\rho \in (0,1)$ je míra odpařování a~$L_k$ délka trasy mravence $k$.
Kratší trasa $\Rightarrow$ více feromonu.
\end{defbox}

\paragraph{Číselný příklad (pravidlo přechodu).} Mravenec stojí ve městě $i$
a~může jít do měst $A, B, C$ se vzdálenostmi $d = (5;\,10;\,2)$
a~feromony $\tau = (2;\,4;\,1)$; $\alpha = 1$, $\beta = 2$.
Viditelnosti: $\eta = (0{,}2;\,0{,}1;\,0{,}5)$. Váhy:
\[
w_A = 2 \cdot 0{,}2^2 = 0{,}08,\quad
w_B = 4 \cdot 0{,}1^2 = 0{,}04,\quad
w_C = 1 \cdot 0{,}5^2 = 0{,}25,\quad \textstyle\sum = 0{,}37 .
\]
Pravděpodobnosti: $p_A = 0{,}216$, $p_B = 0{,}108$, $p_C = 0{,}676$.
Přestože hrana do $B$ má nejvíce feromonu, krátká vzdálenost do $C$ (umocněná
$\beta = 2$) rozhoduje.

\paragraph{Číselný příklad (aktualizace).} Nechť $\tau_{ij} = 2$, $\rho = 0{,}5$,
$Q = 1$ a~hranou prošli dva mravenci s~délkami tras $L_1 = 20$, $L_2 = 25$:
\[
\tau_{ij} \gets 0{,}5 \cdot 2 + \left(\tfrac{1}{20} + \tfrac{1}{25}\right)
= 1 + 0{,}09 = 1{,}09 .
\]

\paragraph{ACS (Ant Colony System).} Vylepšení konkurenceschopné se
specializovanými řešiči TSP; inspirováno Q-learningem:
\begin{itemize}
\item \textbf{Globální aktualizace jen nejlepším řešením} (elitismus): feromon
  přidává pouze globálně (nebo iteračně) nejlepší mravenec.
\item \textbf{Lokální aktualizace}: kdykoli mravenec projde hranou, feromon
  na ní se mírně sníží ($\tau \gets (1-\xi)\tau + \xi\tau_0$) --- tím se hrana
  stává méně atraktivní pro ostatní a~roste rozmanitost tras.
\item \textbf{Pseudonáhodné pravidlo výběru}: s~pravděpodobností $q_0$ se
  zvolí deterministicky nejlepší hrana podle
  $\tau_{il}\eta_{il}^{\beta}$ (exploatace), jinak se použije standardní
  pravidlo AS (explorace).
\end{itemize}

\paragraph{Další varianty.}
\emph{Elitist AS} (nejlepší trasa přispívá navíc),
\emph{Rank-based AS} (přispívá $N$ nejlepších, váha podle pořadí),
\textbf{MAX--MIN AS} (feromon je držen v~intervalu
$[\tau_{\min}, \tau_{\max}]$, aktualizuje jen nejlepší, inicializace na
$\tau_{\max}$ --- brání předčasné stagnaci),
ACO pro spojitou optimalizaci (diskretizace prostoru na podintervaly,
feromon vede hledání do lepších podprostorů).

\paragraph{Vlastnosti a~aplikace.} Decentralizace, emergence, nepřímá
komunikace; blízko reaktivním architekturám agentů (Brooksova subsumpční
architektura). Velkou předností je schopnost \textbf{optimalizovat
v~dynamickém prostředí} --- při změně grafu se feromony přizpůsobí.
Použití: TSP, směrování vozidel, směrování v~sítích, sekvenování DNA,
skládání proteinů, detekce hran v~obraze, rozvrhování.

\subsection{Další rojové algoritmy}

\noindent\mimo{}

\begin{itemize}
\item \textbf{Umělá včelí kolonie} (\emph{artificial bee colony}, ABC;
  Karaboga 2005). Populace zdrojů potravy = řešení. Tři role:
  \emph{dělnice} (prohledávají okolí svého zdroje),
  \emph{pozorovatelky} (vybírají zdroj ruletou podle kvality a~prohledávají
  jeho okolí),
  \emph{průzkumnice} (zdroj, který se $limit$ pokusů nezlepšil, se opustí
  a~nahradí náhodným --- vestavěný restart).
\item \textbf{Firefly algorithm} (Yang, 2008). Jedinci = světlušky; jasnost je
  úměrná fitness a~s~vzdáleností klesá jako
  $I(r) = I_0 e^{-\gamma r^2}$. Méně jasná světluška se posouvá směrem
  k~jasnější:
  $\vec x_i \gets \vec x_i + \beta_0 e^{-\gamma r_{ij}^2}(\vec x_j - \vec x_i)
  + \alpha \vec\varepsilon$. Díky omezenému dosahu přitažlivosti se roj
  přirozeně dělí na podskupiny $\Rightarrow$ vhodné pro multimodální úlohy.
\item \textbf{Bat algorithm}, \textbf{cuckoo search}, \textbf{grey wolf
  optimizer} a~desítky dalších \uv{metaforických} algoritmů. Kritika:
  mnohé jsou jen přejmenované varianty PSO či ES a~přinášejí spíše
  terminologický zmatek než nové principy.
\end{itemize}

\subsection{Shrnutí ke~zkoušce}

\begin{itemize}
\item PSO: rovnice
  $\vec v \gets w\vec v + c_1r_1(\vec p - \vec x) + c_2r_2(\vec g - \vec x)$,
  $\vec x \gets \vec x + \vec v$; setrvačnost $w$ (klesá 0{,}9 $\to$ 0{,}4),
  kognitivní a~sociální složka, $c_1=c_2=2$; omezení $v_{\max}$ nebo
  constriction $\chi \approx 0{,}729$; topologie gbest vs.\ lbest (kruh,
  mřížka).
\item Rozdíl PSO vs.\ GA: žádné operátory, částice mají paměť
  ($\vec p_i$, $\vec g$), informace teče jednosměrně od lepších.
\item ACO: pravděpodobnost přechodu
  $p_{ij} \propto \tau_{ij}^{\alpha}\eta_{ij}^{\beta}$ s~$\eta = 1/d$;
  aktualizace $\tau \gets (1-\rho)\tau + \sum_k Q/L_k$ (odpařování + uložení
  úměrné kvalitě).
\item AS (původní, všichni aktualizují) vs.\ ACS (globální update jen
  nejlepším, lokální update při průchodu, pseudonáhodné pravidlo $q_0$)
  vs.\ MAX--MIN (meze feromonu). Akce démona = lokální prohledávání.
\item ABC (dělnice / pozorovatelky / průzkumnice s~limitem = restart),
  firefly (přitažlivost klesá exponenciálně se vzdáleností).
\item Klíčová slova: stigmergie, emergence, decentralizace, kladná zpětná
  vazba + odpařování, vhodnost pro dynamické prostředí.
\end{itemize}

%=====================================================================
\section{Lokální prohledávání a~memetické algoritmy}
\label{sec:lokalni}
%=====================================================================

\subsection{Lokální prohledávání a~horolezecký algoritmus}

\begin{defbox}[Okolí a~lokální prohledávání]
Pro prostor řešení $X$ definujeme \emph{okolí} $N(x) \subseteq X$ --- množinu
řešení dosažitelných z~$x$ jednou elementární změnou (překlopení bitu, prohození
dvou měst, 2-opt krok). \emph{Lokální prohledávání} iterativně přechází
do souseda podle daného pravidla. \emph{Lokální optimum} vzhledem k~$N$ je
$x$ takové, že žádný soused není lepší --- závisí tedy na definici okolí!
\end{defbox}

\begin{defbox}[Horolezecký algoritmus (\emph{hill climbing})]
\textbf{1. Steepest ascent (nejstrmější stoupání):} projdi \emph{celé} okolí
a~přejdi do nejlepšího souseda, je-li lepší než $x$; jinak skonči.\\
\textbf{2. First-improvement (první zlepšení):} procházej okolí a~přejdi do
prvního nalezeného lepšího souseda (levnější, často stejně dobré).\\
\textbf{3. Stochastic hill climbing:} vyber náhodného souseda; přijmi jej,
je-li lepší (případně s~pravděpodobností úměrnou zlepšení).\\
\textbf{4. Random-restart hill climbing:} po uvíznutí začni z~nové náhodné
pozice; pamatuj si nejlepší nalezené řešení. S~rostoucím počtem restartů
konverguje pravděpodobnost nalezení globálního optima k~1.
\end{defbox}

Slabiny: uvíznutí v~lokálním optimu, na \emph{plošinách} (plateau, kde jsou
všichni sousedé stejní --- pomáhají \emph{sideways moves}) a~na hřebenech.
Výhody: jednoduchost, malá paměť, rychlost --- proto je ideální jako
vylepšovací procedura uvnitř EA.

\subsection{Simulované žíhání}

\begin{defbox}[Simulované žíhání (\emph{simulated annealing}, SA)]
Kirkpatrick, Gelatt, Vecchi (1983); fyzikální analogie s~pomalým chladnutím
kovu, při němž atomy zaujmou uspořádání s~minimální energií. Algoritmus je
horolezec, který \textbf{připouští i~zhoršující kroky} s~pravděpodobností
klesající s~velikostí zhoršení a~s~teplotou.
\end{defbox}

\begin{defbox}[Metropolisovo kritérium]
Pro minimalizaci, kandidáta $y \in N(x)$ a~$\Delta = f(y) - f(x)$:
\[
P(\text{přijmout } y) =
\begin{cases}
1, & \Delta \le 0 \quad (\text{zlepšení}),\\[2pt]
e^{-\Delta / T}, & \Delta > 0 \quad (\text{zhoršení}),
\end{cases}
\]
kde $T > 0$ je teplota.
\end{defbox}

\textbf{Číselný příklad.} Zhoršení $\Delta = 2$. Při $T = 1$ je
$P = e^{-2} = 0{,}135$; při $T = 0{,}5$ je $P = e^{-4} = 0{,}018$;
při $T = 5$ je $P = e^{-0{,}4} = 0{,}67$. Na začátku (horko) se tedy algoritmus
chová téměř jako náhodná procházka, na konci (chladno) jako čistý horolezec.

\begin{algorithm}[!ht]
\caption{Simulované žíhání}
\begin{algorithmic}[1]
\State $x \gets$ počáteční řešení; $T \gets T_0$; $x^* \gets x$
\While{$T > T_{\min}$ a~není splněna ukončovací podmínka}
  \For{$k$ krát (délka Markovova řetězce při dané teplotě)}
    \State vyber náhodně $y \in N(x)$; $\Delta \gets f(y) - f(x)$
    \If{$\Delta \le 0$ nebo $\mathrm{rand}() < e^{-\Delta/T}$} $x \gets y$ \EndIf
    \If{$f(x) < f(x^*)$} $x^* \gets x$ \EndIf
  \EndFor
  \State $T \gets \mathrm{cool}(T)$
\EndWhile
\Ensure $x^*$
\end{algorithmic}
\end{algorithm}

\paragraph{Rozvrh chlazení (\emph{cooling schedule}).}
\begin{itemize}
\item \emph{Geometrický:} $T_{k+1} = \alpha T_k$, $\alpha \in [0{,}8; 0{,}99]$ ---
  v~praxi nejběžnější.
\item \emph{Lineární:} $T_k = T_0 - \beta k$.
\item \emph{Logaritmický:} $T_k = c/\log(k+1)$ --- jediný, pro který je
  dokázána konvergence.
\item Adaptivní (reheating: při stagnaci se teplota zvýší).
\end{itemize}

\begin{veta}[Konvergence SA]
Při logaritmickém rozvrhu $T_k \ge c/\log(k+1)$ s~dostatečně velkou konstantou
$c$ (souvisí s~výškou nejhlubší \uv{bariéry} v~krajině) konverguje simulované
žíhání k~množině globálních optim s~pravděpodobností 1.
\end{veta}
Prakticky je tento rozvrh nepoužitelně pomalý; SA je tedy teoreticky
uspokojivá, ale v~praxi se používají rychlejší heuristické rozvrhy. Formálně
je SA nehomogenní Markovův řetězec; při pevné $T$ konverguje k~Boltzmannovu
rozdělení $\pi(x) \propto e^{-f(x)/T}$, které se pro $T \to 0$ koncentruje
na globálních optimech. Souvislost s~EA: \emph{Boltzmannovský turnaj}
(odst.~\ref{ssec:selekce}) přenáší tutéž myšlenku do selekce.

\subsection{Tabu prohledávání}

\begin{defbox}[Tabu search (Glover, 1986)]
Lokální prohledávání s~\emph{krátkodobou pamětí}: uchovává se \emph{tabu list}
nedávno provedených tahů (efektivnější než ukládat celá řešení), které jsou
po dobu \emph{tabu tenure} zakázány. V~každém kroku se přejde do
\textbf{nejlepšího nezakázaného souseda}, a~to i~tehdy, je-li horší než
aktuální řešení --- tím se algoritmus dostane z~lokálního optima a~necyklí.
\end{defbox}

\begin{itemize}
\item \textbf{Tabu tenure} $T$: statická (pevný počet iterací) nebo dynamická
  (náhodná z~rozsahu).
\item \textbf{Aspirační kritérium}: tabu se poruší, vede-li tah k~řešení
  lepšímu než dosud nejlepší nalezené (jinak by tabu mohlo blokovat postup).
\item \textbf{Střednědobá paměť}: pravidla směrující hledání do slibných
  oblastí (intenzifikace).
\item \textbf{Dlouhodobá paměť / frekvenční paměť}: počítadlo, jak často byl
  který prvek řešení použit; často navštěvované konfigurace se penalizují
  (diverzifikace), případně se restartuje z~málo navštívené oblasti.
\end{itemize}
Aplikace: TSP (výměny hran, \emph{ejection chains}), rozvozní úlohy,
sekvenování DNA, rozvrhování. Pro: únik z~lokálních optim, žádné cyklení,
funguje v~diskrétních i~spojitých doménách. Proti: mnoho parametrů, vyšší
výpočetní náročnost (procházení celého okolí).

\subsection{Scatter search}

Populační metaheuristika zaměřená na \textbf{diverzitu}. Udržuje malou
\emph{referenční množinu} $R$ (řádově desítky jedinců), do níž se vybírají
jedinci podle kombinace kvality a~vzájemné vzdálenosti (maximalizuje se
minimální vzdálenost nově přidávaného jedince od zbytku $R$).
Cyklus: (1)~\emph{diverzifikace} --- chytrá inicializace,
(2)~\emph{generování podmnožin} --- typicky všechny dvojice z~$R$,
(3)~\emph{kombinace} --- aritmetické (či permutační) křížení,
(4)~\emph{vylepšení} --- lokální prohledávání, mutace, tabu search,
(5)~\emph{aktualizace $R$}. Referenční množinu lze aktualizovat
\emph{inkrementálně} (přidá se jeden či několik jedinců za generaci),
\emph{úplnou obnovou} každou generaci, nebo ji tvořit už ve vnitřním cyklu ---
noví jedinci se do $R$ zařazují okamžitě a~rekombinuje se rovnou s~nimi.
Vhodné pro úlohy s~velmi drahou fitness (např.\ black-box optimalizátor
OptQuest).

\subsection{Memetické algoritmy}

\begin{defbox}[Mem a~memetický algoritmus]
\emph{Mem} (Dawkins, 1976) je jednotka kulturní informace šířící se
napodobováním --- \uv{gen kultury}. \emph{Memetický algoritmus} (Moscato, 1989)
je hybridní metoda: \textbf{evoluční algoritmus s~vnořeným lokálním
prohledáváním}, které vylepšuje jednotlivé jedince (\uv{individuální učení}).
Synonyma: hybridní EA, genetické lokální prohledávání, lamarckovské či
baldwinovské EA, kulturní algoritmy.
\end{defbox}

\begin{algorithm}[!ht]
\caption{Memetický algoritmus}
\begin{algorithmic}[1]
\State inicializuj populaci
\While{není splněna ukončovací podmínka}
  \State ohodnoť jedince v~populaci
  \State vytvoř novou populaci stochastickými operátory (selekce, křížení, mutace)
  \State vyber podmnožinu $O$ jedinců, kteří projdou vylepšením
  \For{každého $x \in O$}
    \State proveď individuální učení $x$ pomocí memu (lokálního prohledávače),
      s~pravděpodobností $p$ po dobu $t$
    \State výsledek zpracuj \textbf{lamarckovsky} nebo \textbf{baldwinovsky}
  \EndFor
\EndWhile
\end{algorithmic}
\end{algorithm}

Podstata: lokální prohledávání vnořené do evolučního cyklu. Genetické operátory
přenášejí informaci \emph{mezi} jednotlivými běhy lokálního prohledávače, takže
celkové hledání je mnohem efektivnější než opakovaný restart lokální metody
(Krasnogor \& Smith). Jako mem lze použít hill climbing, simulované žíhání,
tabu search, kombinatorickou heuristiku (2-opt) nebo gradientní metodu
(zpětné šíření chyby při učení neuronové sítě).

\begin{defbox}[Lamarckismus vs.\ baldwinismus]
Co udělat s~vylepšeným jedincem $x' = \mathrm{LS}(x)$?
\begin{itemize}
\item \textbf{Lamarckovsky:} do populace se vloží $x'$ i~s~jeho fitness ---
  \emph{genotyp se změní podle fenotypu}. Z~darwinistického hlediska
  \uv{nekošer}, ale prakticky rychlejší.
\item \textbf{Baldwinovsky:} jedinci $x$ se přiřadí \emph{fitness}
  vylepšeného $x'$, ale genotyp $x$ zůstane \emph{nezměněn}. Darwinisticky
  korektní; fitness se interpretuje jako \uv{potenciál} jedince
  (\emph{Baldwinův efekt}: schopnost učit se zvýhodňuje jedince a~evoluce
  se posléze \uv{doučí} sama --- genetická asimilace).
\end{itemize}
Lamarckismus konverguje rychleji, ale více snižuje diverzitu a~může zkreslit
krajinu; baldwinismus zachovává diverzitu a~vyhlazuje fitness krajinu,
je však pomalejší.
\end{defbox}

\paragraph{Návrhové otázky.} Které jedince vylepšovat (všechny / jen elitu /
náhodný podíl $p$)? Jak dlouho ($t$ kroků)? Jak silně (hloubka lokálního
prohledávání)? Přílišné lokální prohledávání znamená zbytečné náklady
a~ztrátu diverzity; příliš malé nepřináší užitek.

\paragraph{Memetické počítání (\emph{memetic computing}).} Zobecnění:
memy nejsou známy předem, ale jsou samy předmětem hledání.
\begin{itemize}
\item \textbf{Multi-mem MA:} mem je zakódován \emph{jako součást genotypu}
  jedince; dědí se, kříží se (potomek zdědí mem lepšího rodiče, nebo náhodně)
  a~používá se pro vylepšení tohoto jedince.
\item \textbf{Meta-lamarckovský MA:} memy mají \emph{vlastní populaci} a~s~jedinci
  koevolvují; při vylepšení se jedinci přiřadí mem, úspěch jedince je odměnou
  pro mem, nad memy běží druhý evoluční algoritmus.
\item \emph{Memetický prostor} jako protějšek prostoru fenotypů: memy lze
  vybírat heuristikami, nechat je soutěžit podle úspěšnosti, dělit do
  subpopulací nebo řídit meta-memetickým algoritmem.
\end{itemize}

\subsection{Ladění a~řízení parametrů}

Úzce souvisí: EA má mnoho parametrů (velikost populace, $p_C$, $p_M$, velikost
turnaje, míra elitismu). Ty nejsou nezávislé a~vyčerpávající prohledání není
možné.
\begin{itemize}
\item \textbf{Ladění} (\emph{tuning}) --- parametry se určí \emph{před} během
  (pokus-omyl, grid search, meta-EA, iterované závodění --- irace).
\item \textbf{Řízení} (\emph{control}) --- parametry se mění \emph{během} běhu:
  \begin{itemize}
  \item \emph{Pevná strategie} (deterministická): závislost jen na čase
    (např.\ $\sigma(t) = 1 - 0{,}9\,t/T$, pokles $p_M$ o~0{,}1 \% za generaci);
    stochastická varianta = simulované žíhání.
  \item \emph{Adaptivní}: podle zpětné vazby z~běhu (kvalita řešení, rozptyl
    fitness, diverzita) --- např.\ pravidlo 1/5, hypermutace při poklesu
    fitness, adaptivní zjemňování reprezentace (zvyšování počtu bitů na
    reálné číslo, klesne-li diverzita měřená Hammingovou vzdáleností nejlepšího
    a~nejhoršího jedince).
  \item \emph{Samoadaptivní}: parametry jsou součástí genomu a~evolvují
    (ES, EP, multi-mem MA).
  \end{itemize}
\item Řídit lze i~\emph{velikost populace} nepřímo: každému novému jedinci se
  přidělí \uv{délka života} úměrná jeho fitness; po jejím vypršení zemře.
  Lepší jedinci žijí déle, velikost populace je odvozená dynamická veličina.
\end{itemize}

\subsection{Shrnutí ke~zkoušce}

\begin{itemize}
\item Hill climbing: steepest ascent / first improvement / stochastic /
  s~náhodnými restarty; problémy: lokální optima, plošiny, hřebeny.
\item SA: Metropolisovo kritérium $P = \min(1, e^{-\Delta/T})$; rozvrh chlazení
  (geometrický v~praxi, logaritmický pro důkaz konvergence); pro pevné $T$
  konverguje k~Boltzmannovu rozdělení $\propto e^{-f/T}$.
\item Tabu search: tabu list tahů, tenure, aspirační kritérium, střednědobá
  a~dlouhodobá (frekvenční) paměť; vždy se jde do nejlepšího nezakázaného
  souseda, i~když je horší.
\item Scatter search: referenční množina kombinující kvalitu a~diverzitu,
  aritmetické kombinace + vylepšení.
\item Memetický algoritmus = EA + lokální prohledávání nad jedinci;
  \textbf{lamarckismus} (mění genotyp, rychlejší) vs.\ \textbf{baldwinismus}
  (mění jen fitness, zachovává diverzitu). Memetic computing: memy v~genomu
  nebo v~koevolvující populaci.
\item Ladění vs.\ řízení parametrů; pevné / adaptivní / samoadaptivní strategie.
\end{itemize}


%=====================================================================
\section{Aplikace evolučních algoritmů}
\label{sec:aplikace}
%=====================================================================

Poslední věta okruhu vyjmenovává čtyři aplikační oblasti; každá tvoří jeden oddíl této
kapitoly. Společné je jim to, že v~každé rozhoduje \textbf{volba reprezentace a~operátorů},
nikoli volba \uv{lepšího} evolučního algoritmu.

\subsection{Evoluce pravidlových a~expertních systémů}
\label{sec:lcs}

\subsubsection{Michigan vs.\ Pittsburgh}

Cílem je naučit se \textbf{sadu pravidel} tvaru \texttt{IF podmínka THEN akce}
z~trénovacích dat nebo z~interakce s~prostředím (dobývání znalostí, expertní
systémy, řízení robotů, zpětnovazební učení). Existují dva základní přístupy:

\begin{center}
\begin{tabular}{@{}p{2.8cm}p{6.2cm}p{6.2cm}@{}}
\toprule
 & \textbf{Michigan} (Holland) & \textbf{Pittsburgh} (Smith) \\
\midrule
jedinec & \textbf{jedno pravidlo} & \textbf{celá sada pravidel} \\
řešení & celá populace dohromady & jeden jedinec \\
fitness & síla/přesnost pravidla (nutné rozdělení odměny) & kvalita celého
  systému (přímočará) \\
operátory & standardní, evoluce probíhá jen občas a~na části populace &
  složité, desítky operátorů nad sadami i~jednotlivými pravidly \\
výhody & on-line učení, malý prostor, přirozené pro RL & jednoznačná fitness,
  žádný credit assignment \\
nevýhody & přiřazení zásluh, pravidla musí spolupracovat, ale zároveň soutěží &
  velcí jedinci, drahé vyhodnocení, riziko bloatu \\
\bottomrule
\end{tabular}
\end{center}

\subsubsection{Learning classifier systems (Michigan)}

\begin{defbox}[LCS (\emph{learning classifier system})]
Systém, kde populace jednoduchých pravidel (\emph{klasifikátorů}) tvoří
dohromady expertní či řídicí systém. Klasifikátor má levou stranu
(podmínku) jako řetězec nad $\{0,1,*\}$ porovnávaný se vstupem, pravou stranu
(akci či klasifikační třídu) a~\textbf{váhu / sílu} odrážející jeho úspěšnost,
která slouží jako fitness. Evoluce neprobíhá generačně, ale průběžně a~jen na
části populace.
\end{defbox}

\paragraph{Architektura LCS.} Klíčové je uvědomit si, že \emph{expertním
systémem je celá populace}, nikoli jedinec. Systém tvoří smyčka:
\begin{enumerate}
\item \textbf{Detektory} převedou stav prostředí na vstupní zprávu.
\item Zpráva se uloží do \textbf{bufferu zpráv} a~porovná se s~levými stranami
  všech pravidel $\Rightarrow$ \emph{match set}.
\item Mezi vyhovujícími pravidly proběhne \textbf{aukce/soutěž} o~právo jednat
  (podle síly pravidla).
\item Vítězné pravidlo zapíše svou pravou stranu buď do bufferu (vnitřní
  zpráva), nebo přes \textbf{efektory} do prostředí jako akci.
\item Prostředí vrátí \textbf{odměnu}; ta se rozdělí mechanismem
  \emph{bucket brigade} mezi pravidla, která k~ní vedla.
\item Občas a~jen na části populace běží \textbf{GA}, který pravidla obměňuje.
\end{enumerate}

\paragraph{Problém reaktivity.} Čistě reaktivní systém nemá vnitřní paměť.
Holland proto zavedl \emph{zprávy} (\emph{messages}): pravá strana pravidla
může kromě akce zapsat vnitřní zprávu do bufferu, levá strana má
\emph{receptory} pro její zachycení. Vzniká tak řetězení pravidel --- systém
získá vnitřní stav (a~tím i~výpočetní sílu), ale zároveň s~ním přichází
problém, jak rozdělit odměnu mezi celý řetězec.

\begin{defbox}[Bucket brigade (algoritmus \uv{požárního řetězu})]
Mechanismus rozdělení odměny v~LCS. Pravidlo, které chce být aplikováno,
musí \uv{zaplatit} část své síly pravidlu, jež ho aktivovalo (jako by kupovalo
právo jednat). Odměna z~prostředí připadne poslednímu článku řetězu a~postupně
\uv{proteče} zpět k~pravidlům, která k~ní vedla. Prakticky je obtížné
vyvážit tuto ekonomiku pravidel, dnes se proto téměř nepoužívá.
\end{defbox}

\paragraph{ZCS (Wilson, 1994) --- \uv{zero-level} LCS.} Zjednodušení:
žádné vnitřní zprávy, žádná složitá redistribuce.
\begin{itemize}
\item $P$ --- populace pravidel; vstup $x$ z~detektorů.
\item $M$ --- \emph{match set}: pravidla, jejichž podmínka odpovídá $x$.
\item Akce $a$ se z~$M$ vybere ruletou podle síly pravidel;
  $A \subseteq M$ --- \emph{action set}: pravidla prosazující zvolenou akci.
\item \textbf{Posílení:} z~pravidel v~$A$ se odebere podíl $\beta$ jejich síly
  do \uv{kbelíku} $B$; po obdržení odměny $r$ se každému pravidlu v~$A$ přičte
  $\beta r/|A|$ a~pravidlům z~předchozího action setu $A_{-1}$ se přičte
  $\gamma B/|A_{-1}|$. Pravidla z~$M \setminus A$ (odpovídala, ale prosazovala
  jinou akci) jsou zdaněna.
\item \textbf{Cover operátor:} neexistuje-li pro aktuální vstup žádné pravidlo,
  vygeneruje se ad hoc (do podmínky se náhodně přidají hvězdičky, akce se zvolí
  náhodně).
\end{itemize}

\paragraph{XCS (Wilson, 1995) --- fitness založená na přesnosti.}
Slabiny ZCS: neevolvuje \emph{úplný} systém pravidel pokrývající všechny
situace; pravidla na začátku řetězců jsou zřídka odměňována a~vymírají;
pravidla vedoucí k~malým, ale správně předpovězeným odměnám také vymírají.
Příčinou je, že síla sloužila zároveň jako fitness pro GA i~jako kritérium
výběru akce. Wilsonova myšlenka: \emph{predikce} má odhadovat odměnu, ale
\emph{evoluce} má preferovat \textbf{spolehlivé (přesné)} klasifikátory.

\begin{defbox}[XCS --- klasifikátor jako pětice]
$(c, a, p, \epsilon, F)$: podmínka, akce, predikce odměny $p$, chyba predikce
$\epsilon$, fitness $F$. Přesnost se počítá jako klesající funkce chyby,
$\kappa = \alpha\,(\epsilon/\epsilon_0)^{-\nu}$ (pro $\epsilon > \epsilon_0$,
jinak $\kappa = 1$; $\alpha, \epsilon_0, \nu$ jsou parametry);
fitness je \emph{relativní} přesnost v~rámci action setu.
Predikce a~chyba se aktualizují pravidly Q-learningu.
\end{defbox}

Běh XCS: podle vstupu se sestaví match set $M$, rozdělí se na action sety
$A_i$ podle akcí; pro každou akci se predikuje výplata jako průměr predikcí
vážený fitness; akce se zvolí buď maximem (exploatace), nebo ruletou
(explorace); provede se; odměna se rozdělí do action setu předchozího kroku
(aktualizace $p$ pomocí $r + \gamma \max_a p_a$, pak $\epsilon$, pak $F$);
GA běží jako \textbf{nikový} --- pouze uvnitř action setu, nikoli nad celou
populací.

Výsledek: XCS propojuje LCS se zpětnovazebním učením (Q-learning jako
mechanismus přiřazení zásluh), evolvuje \textbf{úplnou a~minimální} mapu
vstup--výstup (nejen optimální politiku, ale i~kompaktní srozumitelný popis)
a~má řadu praktických aplikací (dobývání znalostí, řízení, bioinformatika).

\paragraph{Dnešní rodina LCS.} Podle úlohy se rozlišují:
\emph{XCS} (zpětnovazební učení, predikce odměny),
\textbf{UCS} (\emph{sUpervised Classifier System}) --- varianta pro učení
s~učitelem, kde je správná třída známa, takže se místo predikce odměny
rovnou měří přesnost klasifikace,
a~\emph{XCSF} pro aproximaci spojitých funkcí (pravá strana je lineární model,
nikoli konstanta). Podmínky se dnes běžně zapisují nejen nad
$\{0,1,*\}$, ale i~intervaly pro reálné atributy, což z~LCS dělá plnohodnotný
nástroj dobývání znalostí --- s~tou předností, že výsledkem je
\textbf{čitelná sada pravidel}.

\subsubsection{Pittsburghský přístup a~systém GIL}

Jedinec = celá sada pravidel, tedy kompletní klasifikační systém. Hodnocení je
složitější (priority pravidel, konflikty, falešně pozitivní a~negativní
odpovědi) a~operátorů je celá řada, na několika úrovních. Důraz se klade na
bohatou doménovou reprezentaci (množiny, výčty, intervaly).

\textbf{GIL} je klasický představitel: pro binární klasifikaci, jedinec
klasifikuje implicitně do jedné třídy (pravidla nemají pravou stranu).
Jedinec je \emph{disjunkce komplexů}, komplex je \emph{konjunkce selektorů}
(po jednom pro každou proměnnou) a~selektor je \emph{disjunkce hodnot}
z~domény dané proměnné, zakódovaná bitovou maskou. Například
\[
\big((X{=}A_1) \wedge (Z{=}C_3)\big) \vee \big((X{=}A_2) \wedge (Y{=}B_2)\big)
\;\equiv\;
[\,001|11|0011 \;\vee\; 010|10|1111\,].
\]
Operátory na třech úrovních:
\begin{itemize}
\item \emph{jedinec}: výměna pravidel, kopie pravidla, generalizace pravidla,
  smazání pravidla, specializace pravidla, zahrnutí pozitivního příkladu;
\item \emph{komplex}: rozdělení komplexu na jednom selektoru, generalizace
  selektoru (nahrazení $11\ldots1$), specializace, zahrnutí negativního
  příkladu;
\item \emph{selektor}: mutace $0 \leftrightarrow 1$, rozšíření $0 \to 1$,
  zúžení $1 \to 0$.
\end{itemize}
Tyto operátory přímo odpovídají operacím induktivního učení
(generalizace/specializace), takže GIL je vlastně evolučně řízený
induktivní algoritmus.

\subsubsection{Připomenutí: metriky klasifikace}

Z~matice záměn (TP, FP, FN, TN) se počítá
$\text{přesnost} = (TP+TN)/(TP+TN+FP+FN)$,
$\text{senzitivita} = TP/(TP+FN)$ a~$\text{specificita} = TN/(TN+FP)$.
Fitness pravidlového systému bývá kombinací přesnosti, pokrytí a~jednoduchosti
(počet pravidel) --- tedy přirozeně vícekriteriální úloha.

\subsection{Neuroevoluce}
\label{sec:neuro}

\begin{defbox}[Neuroevoluce (\emph{neuroevolution})]
\uv{Neuroevoluce je metoda pro úpravu vah, topologií nebo ansámblů neuronových
sítí za účelem naučení konkrétní úlohy. Evoluční výpočty se používají
k~hledání parametrů sítě maximalizujících fitness, která měří výkon v~úloze.
Ve srovnání s~jinými metodami učení je neuroevoluce velmi obecná --- umožňuje
učení bez explicitních cílových hodnot, s~nediferencovatelnými aktivačními
funkcemi a~s~rekurentními sítěmi.} (Miikkulainen, 2017)
\end{defbox}

Co lze evolvovat: váhy, další parametry (aktivační funkce, konstanty učení),
architekturu (topologii, propojení), architekturu i~váhy zároveň, případně
parametry samotného učicího algoritmu. Hlavní doménou je
\textbf{zpětnovazební učení} a~robotika, kde nejsou k~dispozici cílové výstupy
pro učení s~učitelem.

\subsubsection{Evoluce vah}

Přímočaré: váhy se zakódují do vektoru pevné délky a~použije se reálný GA,
ES nebo DE se standardními operátory. Vlastnosti:
\begin{itemize}
\item Pomalejší než specializované gradientní metody (zpětné šíření chyby)
  pro učení s~učitelem.
\item \textbf{+} snadná paralelizace, nepotřebuje gradient (funguje pro
  nediferencovatelné a~rekurentní sítě), použitelné pro RL.
\item \textbf{Problém konkurujících konvencí} (\emph{competing conventions},
  permutační problém): tatáž funkce sítě odpovídá mnoha různým genotypům
  (permutace skrytých neuronů), takže křížení dvou funkčně shodných rodičů
  často vytvoří nefunkčního potomka.
\item Novodobý návrat zájmu: evoluce vah je konkurenceschopná i~pro velké
  sítě, používá-li se hodnocení na minibatchích.
\end{itemize}

\subsubsection{Evoluce architektury}

Fitness architektury = odhad výkonu sítě: síť je nutné postavit, inicializovat
a~naučit (zpětným šířením či jiným EA), nejlépe vícekrát --- proto je
vyhodnocení fitness drahé a~zašuměné.

\begin{itemize}
\item \textbf{Přímé kódování} (\emph{direct encoding}): struktura propojení jako
  binární matice; jedinec je buď linearizovaná matice (standardní operátory),
  nebo matice sama (speciální 2D operátory). Jednoduché, ale neškáluje
  (velikost genomu roste kvadraticky).
\item \textbf{Gramatické kódování} (Kitano, 1990): matici propojení generuje
  jednoduchá 2D formální gramatika a~evolvují se \emph{pravidla gramatiky}.
  Kompaktní (logaritmické), umožňuje opakující se motivy, ale je nepřímé
  a~hůře laditelné.
\item \textbf{Buněčné kódování} (\emph{cellular encoding}, Gruau): jedincem je
  \textbf{program v~GP}, který popisuje, jak síť \emph{vyrůst} z~jediné buňky
  pomocí operací: přidej neuron, rozděl neuron sériově, rozděl paralelně,
  přepoj synapsi, změň váhu. Vývojový (developmental) přístup.
\item \textbf{Simulovaný růst} spojů ve 2D rovině --- rané pokusy v~evoluční
  robotice, neškálovalo.
\end{itemize}

\paragraph{GNARL (Angeline a~kol., 1993).} Evoluční programování nad grafy:
vstupní a~výstupní neurony jsou pevné, skryté neurony a~spoje se evolvují;
architektura i~váhy \emph{současně}. Dva druhy mutací: \emph{strukturální}
(přidej neuron bez spojů, přidej spoj s~nulovou váhou, smaž neuron, smaž spoj)
a~\emph{parametrická} (zaujatá mutace váhy). Síla mutace je řízena
\emph{teplotou} (přístup simulovaného žíhání); strukturální změny mají být
malé a~nepříliš destruktivní. Selekce: rodiči je lepší polovina populace.
Křížení se nepoužívá (EP).

\subsubsection{NEAT}

\begin{defbox}[NEAT (\emph{NeuroEvolution of Augmenting Topologies},
K.~Stanley, 2002)]
Síť je reprezentována jako \textbf{seznam hran}; každý gen-hrana obsahuje:
vstupní a~výstupní neuron, váhu, příznak \emph{enabled/disabled}
a~\textbf{globální inovační číslo} (\emph{innovation number}), které se
přiděluje při prvním vzniku dané strukturální novinky.
\end{defbox}

Tři klíčové myšlenky:
\begin{enumerate}
\item \textbf{Inovační čísla řeší problém konkurujících konvencí.}
  Kříží se pouze geny se \emph{shodným inovačním číslem} (mají tedy společný
  evoluční původ); geny přebývající u~jednoho rodiče
  (\emph{disjoint} a~\emph{excess}) se přebírají od zdatnějšího rodiče.
  Křížení tak nerozbíjí strukturu --- žádné \uv{bezhlavé kuře}.
\item \textbf{Speciace (niching) chrání inovace.} Na vektorech inovačních
  čísel se definuje míra podobnosti
  \[
  \delta \;=\; \frac{c_1 E}{N} + \frac{c_2 D}{N} + c_3 \overline{\Delta W},
  \]
  kde $E$ (\emph{excess}) je počet genů, jejichž inovační číslo leží
  \emph{za} nejvyšším inovačním číslem druhého genomu, $D$ (\emph{disjoint})
  počet genů chybějících u~druhého rodiče \emph{uvnitř} společného rozsahu
  inovačních čísel, $\overline{\Delta W}$ průměrný rozdíl vah genů se
  \emph{shodným} inovačním číslem a~$N$ velikost většího genomu (normalizace). Sítě s~$\delta$ pod prahem tvoří jeden \emph{druh};
  fitness se počítá jako sdílená (relativní v~rámci druhu). Nová topologie
  má tak čas \uv{doladit váhy\,} dřív, než ji vytlačí zavedená řešení ---
  jinak by strukturální novinka, která zpočátku fitness zhoršuje, okamžitě
  vymřela.
\item \textbf{Postupné narůstání složitosti} (\emph{complexification}):
  populace startuje z~minimálních sítí (jen vstupy a~výstupy) a~teprve
  mutacemi \emph{přidej spoj} a~\emph{přidej neuron} (rozdělí existující spoj:
  původní se zakáže, vzniknou dva nové) roste. Prohledávání tedy začíná
  v~prostoru nejnižší dimenze a~zvyšuje ji jen tehdy, když se to vyplatí.
\end{enumerate}

\subsubsection{HyperNEAT a~CPPN}

\noindent\mimoq{slidy zmiňují jednou větou, CPPN vůbec}

\textbf{HyperNEAT} rozšiřuje NEAT o~dvě myšlenky. \emph{Substrát}: váhy sítě se nekódují
jako vektor čísel, ale \emph{generuje je funkce} $S : \mathbb{R}^4 \to \mathbb{R}$, která
dostane souřadnice dvou neuronů v~rovině a~vrátí váhu jejich spojení (odtud \uv{hyper}).
\emph{CPPN} (\emph{compositional pattern-producing network}): funkce $S$ je reprezentována
sítí ve tvaru DAG, jejíž uzly počítají různé funkce ze slovníku (sigmoida, gaussovka, sinus,
absolutní hodnota); tuto CPPN evolvuje samotný NEAT. Přínos: reprezentace vah funkcí umožňuje
vyjádřit \textbf{symetrie a~opakující se motivy} a~generovat sítě s~milióny spojů z~malého
genomu.

\subsubsection{Neuroevoluce v~éře hlubokých sítí}

\noindent\mimo{}

\textbf{NAS} (\emph{neural architecture search}) hledá architekturu hlubokých sítí a~popisuje
se třemi osami: \emph{prohledávaný prostor} (posloupnost vrstev, DAG opakovaných buněk, nebo
podarchitektura jedné velké \uv{super sítě}), \emph{prohledávací strategie} (evoluční
algoritmy, náhodné prohledávání jako překvapivě silná baseline, bayesovská optimalizace,
zpětnovazební učení, gradientní DARTS) a~\emph{odhad výkonu} (zkrácené učení, sdílení vah).
Konkrétní systémy: \textbf{ENAS} (sdílení vah mezi kandidáty), \textbf{DeepNEAT} (uzlem
chromozomu je celá vrstva) a~\textbf{CoDeepNEAT} (koevoluce populace modulů a~populace
\uv{blueprintů}). Samostatnou linií je \textbf{ES pro RL} (OpenAI 2017): jednoduchá evoluční
strategie \emph{pouze s~mutacemi} nad milióny vah, triviálně paralelizovatelná --- uzly si
vyměňují jen semínka generátoru náhodných čísel.

\subsection{Kombinatorické úlohy a~práce s~omezeními}
\label{sec:kombi}

Evoluční algoritmy jsou oblíbeným nástrojem pro NP-těžké kombinatorické úlohy:
nezaručí optimum, ale poskytnou dobré řešení v~rozumném čase a~snadno se
přizpůsobí dodatečným (třeba i~ošklivým) omezením reálné úlohy.

\subsubsection{Problém batohu}

\textbf{Zadání.} Batoh o~kapacitě $C_{\max}$, $N$ předmětů s~cenami $v_i$
a~objemy $c_i$. Maximalizuj $\sum_i x_i v_i$ za podmínky
$\sum_i x_i c_i \le C_{\max}$, $x_i \in \{0,1\}$.

\textbf{Kódování} je triviální --- bitová mapa, např.\ $0110010$ znamená
\uv{vezmi předměty 2, 3 a~6}. Standardní operátory (křížení, bit-flip mutace)
fungují. \textbf{Problém je fitness}: jedinci mohou porušovat kapacitu.
Máme dvě protichůdná kritéria: maximalizovat $\sum v_i$ a~zároveň splnit
$\sum c_i \le C_{\max}$.

\subsubsection{Tři obecné strategie pro omezení}

\begin{defbox}[Práce s~omezeními v~EA]
\textbf{1. Penalizace} (\emph{penalty}): nepřípustná řešení zůstávají
v~populaci, ale jejich fitness se sníží, např.
$f'(x) = \sum_i x_i v_i - \alpha \max\big(0, \sum_i x_i c_i - C_{\max}\big)$.
Volba $\alpha$ je zásadní: příliš malá $\Rightarrow$ populace se plní
nepřípustnými řešeními, příliš velká $\Rightarrow$ ztrácí se možnost projít
nepřípustnou oblastí ke slibnému řešení. Používá se i~dynamická
($\alpha$ roste v~čase) či adaptivní penalizace.\\
\textbf{2. Oprava} (\emph{repair}): nepřípustné řešení se opraví
(u~batohu: vyhazuj předměty s~nejhorším poměrem $v_i/c_i$, dokud se vejde).
Opravu lze provést jen pro výpočet fitness (baldwinovsky) nebo trvale
zapsat do genotypu (lamarckovsky).\\
\textbf{3. Dekodér} (\emph{decoder}) / uzavřené operátory: kódování se navrhne
tak, že \emph{každý} genotyp reprezentuje přípustné řešení. U~batohu: bit
1 znamená \uv{přidej předmět, pokud se ještě vejde}, jinak jej přeskoč.
Elegantní a~účinné; nevýhodou je ne\-jed\-no\-znač\-nost mapování (více
genotypů dává tentýž fenotyp) a~jeho horší lokalita.
\end{defbox}

Čtvrtá možnost: chápat porušení omezení jako \textbf{druhé kritérium}
a~použít vícekriteriální EA (odst.~\ref{sec:moea}) --- výsledkem je pak celá
Pareto fronta kompromisů mezi kvalitou řešení a~mírou porušení omezení.

\subsubsection{Problém obchodního cestujícího}

TSP: $N$ měst, najdi nejkratší okružní jízdu. Fitness je triviální (délka
cesty), \textbf{problémem je kódování a~zejména křížení}.

\paragraph{Reprezentace.}
\begin{itemize}
\item \textbf{Sousednostní} (\emph{adjacency}): na pozici $i$ je město $j$
  právě tehdy, vede-li hrana $i \to j$. Např.\ $(248397156)$ odpovídá cestě
  $1{-}2{-}4{-}3{-}8{-}5{-}9{-}6{-}7$. Každá cesta má jedinou reprezentaci,
  ale některé seznamy nejsou platné cesty. Klasické křížení nefunguje;
  zajímavé je, že schémata zde mají smysl (např.\ $(*3*\ldots)$ = všechny cesty
  s~hranou 2--3). \emph{Nepoužívat.}
\item \textbf{Ordinální} (\emph{buffer}): udržuje se seznam měst, gen $i$ je
  \emph{pozice} města v~aktuálním seznamu a~po použití se město ze seznamu
  vyjme. Pro seznam $(123456789)$ je cesta $1{-}2{-}4{-}3{-}8{-}5{-}9{-}6{-}7$
  zakódována jako $(112141311)$. Motivací bylo umožnit standardní jednobodové
  křížení --- to sice produkuje platné cesty, ale výsledek nemá s~rodiči nic
  společného. \emph{Také nepoužívat.}
\item \textbf{Cestová (permutační)}: cesta $5{-}1{-}7{-}8{-}9{-}4{-}6{-}2{-}3$
  je $(517894623)$. Nejpřirozenější; vyžaduje speciální operátory ---
  PMX, OX, CX, ER (odst.~\ref{sec:reprez}).
\item \textbf{Maticová}: binární matice, kde $1$ na pozici $(i,j)$ znamená buď
  hranu $i \to j$, nebo (častěji) že $i$ předchází $j$. Speciální operátory:
  \emph{konjunkce} (bitové AND obou rodičů a~náhodné doplnění chybějících hran),
  \emph{disjunkce} (rozdělení na kvadranty, dva se zahodí, odstraní se spory
  a~zbytek se doplní náhodně).
\end{itemize}

\paragraph{Inicializace.} Náhodné permutace, nebo konstrukční heuristiky:
\emph{nejbližší soused} (začni náhodným městem, vždy jdi do nejbližšího
nenavštíveného) a~\emph{vkládání hran} (k~částečné cestě $T$ vyber nejbližší
město $c$ mimo $T$, najdi hranu $k$--$j$ v~$T$ minimalizující
$d(k,c) + d(c,j) - d(k,j)$, hranu nahraď dvojicí $k$--$c$, $c$--$j$).

\paragraph{Mutace.} Inverze úseku (nejdůležitější --- odpovídá 2-opt kroku),
vložení města jinam, posun podcesty, prohození dvou měst, prohození podcest,
heuristické mutace 2-opt a~3-opt (vezmi dvě hrany a~přepoj čtyři města jinak).
Kombinace ES či GA s~\uv{chytrými mutacemi} typu 2-opt je v~praxi velmi
účinná --- de facto memetický algoritmus.

\subsubsection{SAT}

\noindent\mimoq{slidy řeší batoh, TSP a~rozvrhování; SAT je z~Michalewicze}

\textbf{Zadání.} Formule $f: \mathbb{B}^n \to \mathbb{B}$ v~CNF; najdi ohodnocení $x$
s~$f(x)=1$. 2-SAT je v~P, 3-SAT a~výše je NP-úplný.

\textbf{Fitness.} Samotné $f$ je nepoužitelné (nabývá jen 0/1, krajina je \uv{jehla v~kupce
sena}). Standardem je \textbf{počet splněných klauzulí}; zlepšením je vážení problematických
klauzulí a~\emph{zjemňující funkce} --- regularizační člen, který rozliší jedince se stejným
počtem splněných klauzulí a~pomůže tak z~plošin.

\textbf{Reprezentace} bývá přímo bitový řetězec. Zajímavá je \emph{reálná} varianta:
proměnná $x_j$ se nahradí reálným $y_j$ (pravda $=1$, nepravda $=-1$), literály se přepíší
jako $x \mapsto (1-y)^2$ a~$\neg x \mapsto (1+y)^2$ a~výraz se \emph{minimalizuje}. Hodnota
literálu je tedy \emph{penalta}, a~proto se \textbf{disjunkce} (klauzule) přepíše jako
\textbf{součin} a~\textbf{konjunkce} (formule) jako \textbf{součet}. Klasickou lokální
heuristikou je \textbf{WalkSAT}; kombinace EA + WSAT je typické memetické řešení.

\subsubsection{Rozvrhování a~plánování výroby}

\textbf{Školní rozvrh.} Jedinec je rozvrh jako matice (řádky = učitelé, sloupce = třídy,
hodnoty = kódy předmětů). Mutace míchá předměty, křížení vyměňuje mezi jedinci lepší řádky.
Fitness kombinuje kvalitu jednotlivých řádků a~měkká kritéria (okna, rovnoměrnost).
\textbf{Tvrdá omezení} (učitel nemůže učit dvě třídy naráz) je nutné respektovat
\emph{přímo v~operátorech} --- jinak vzniká příliš mnoho nepřípustných jedinců.

\textbf{Job shop scheduling.} Výrobky se skládají z~dílů, pro každý díl existuje více plánů
výroby na strojích a~přestavení stroje stojí čas; fitness je celkový čas výroby
(\emph{makespan}). Kritické je kódování: \emph{permutační} jedinec (jen pořadí výrobků,
zbytek dourčí dekodér) umožňuje použít křížení z~TSP, ale ukazuje se málo efektivní ---
složitou část řeší dekodér. \emph{Přímé} kódování celého plánu je výkonnější, ale vyžaduje
specializované operátory.

\subsection{Vícekriteriální optimalizace}
\label{sec:moea}

\subsubsection{Dominance a~Pareto fronta}

Místo jediné fitness máme vektor kritérií $f_1,\dots,f_k$ (pro jednoduchost
vše minimalizujeme). Kritéria jsou obvykle v~konfliktu (cena vs.\ kvalita,
přesnost vs.\ velikost modelu), takže neexistuje jediné nejlepší řešení.

\begin{defbox}[Dominance]
Pro jedince $x, y$ definujeme:
\begin{itemize}
\item $x$ \textbf{slabě dominuje} $y$ ($x \preceq y$), právě když
  $f_i(x) \le f_i(y)$ pro všechna $i = 1,\dots,k$.
\item $x$ \textbf{dominuje} $y$ ($x \prec y$), právě když $x \preceq y$
  a~existuje $j$ s~$f_j(x) < f_j(y)$.
\item $x$ a~$y$ jsou \textbf{neporovnatelné}, pokud ani $x \prec y$,
  ani $y \prec x$.
\end{itemize}
Relace $\prec$ je \emph{ostré} (striktní) částečné uspořádání --- je
ireflexivní a~tranzitivní, ale nikoli úplné: proto obecně existuje celá
množina navzájem neporovnatelných optim.
\end{defbox}

\begin{defbox}[Pareto optimalita a~Pareto fronta]
Řešení $x^*$ je \emph{Pareto-optimální}, není-li dominováno žádným
přípustným řešením. Množina všech Pareto-optimálních řešení v~prostoru
proměnných je \emph{Pareto množina}, její obraz v~prostoru kritérií je
\textbf{Pareto fronta}. \emph{Aproximace} Pareto fronty populací je cílem
vícekriteriálních EA.
\end{defbox}

Dvojí cíl algoritmu: (a) \textbf{konvergence} --- dostat se co nejblíže
skutečné frontě; (b) \textbf{diverzita} --- pokrýt frontu rovnoměrně
v~celém rozsahu. Právě proto jsou EA pro tuto úlohu ideální: pracují
s~populací, a~mohou tedy vrátit celou aproximaci fronty v~jediném běhu.

\textbf{Příklad.} Body $A(1,10)$, $B(2,7)$, $C(4,5)$, $D(8,2)$, $E(5,8)$
(minimalizace obou kritérií). $E$ je dominován $C$ ($4 \le 5$, $5 \le 8$,
alespoň jedna ostrá) --- $E$ tedy do fronty nepatří. Body $A, B, C, D$ jsou
navzájem neporovnatelné a~tvoří Pareto frontu.

\subsubsection{Skalarizace --- jednoduchá řešení}

\begin{itemize}
\item \textbf{Lineární skalarizace}: $f = \sum_i w_i f_i$, poté standardní
  jednokriteriální EA (v~kontextu MOEA označované jako SOEA). Nevýhody:
  neumíme volit váhy; jeden běh dá jeden bod fronty; a~zásadně ---
  \textbf{nekonvexní části fronty nelze lineární skalarizací nikdy nalézt}.
\item \textbf{$\varepsilon$-constraint}: minimalizuj $f_1$ s~omezeními
  $f_i \le \varepsilon_i$ pro $i \ge 2$. Umí najít i~nekonvexní části, ale
  je třeba volit konstanty $\varepsilon_i$ a~řešit úlohu s~omezeními.
\item \textbf{Čebyševova (Chebyshev) skalarizace}: minimalizuj
  $\max_i \lambda_i |f_i(x) - z_i^*|$, kde $z^*$ je referenční (ideální) bod.
  Na rozdíl od lineární umí dosáhnout na libovolný bod fronty včetně
  nekonvexních částí.
\end{itemize}

\subsubsection{Historické algoritmy}

\paragraph{VEGA (Schaffer, 1985)} --- jeden z~prvních MOEA. Populace $N$
jedinců se pro každé kritérium $f_i$ seřadí a~vybere se $N/k$ nejlepších
podle $f_i$; tyto podmnožiny se sloučí, zkříží a~zmutují.
Nevýhody: obtížné udržení diverzity a~konvergence ke krajním bodům
optimálním pro jednotlivá kritéria (\uv{středy} fronty se ztrácejí).

\paragraph{NSGA (1994)} --- první použití dominance přímo jako fitness.
Populace se rozdělí do front: $F_1$ = nedominovaní jedinci z~$P$,
$F_2$ = nedominovaní z~$P \setminus F_1$, $F_3$ = nedominovaní
z~$P \setminus (F_1 \cup F_2)$ atd. Jedinci z~první fronty dostanou
\uv{fiktivní} fitness, která se dělí \emph{nikovým faktorem}
$\sum_j sh(i,j)$, kde
$sh(i,j) = 1 - (d(i,j)/d_{\mathrm{share}})^2$ pro $d(i,j) < d_{\mathrm{share}}$
a~0 jinak. Druhá fronta dostane fitness menší než nejhorší jedinec první
fronty, opět dělenou nikovým faktorem, atd.
Nevýhody: nutnost nastavit $d_{\mathrm{share}}$, chybějící elitismus,
výpočetní složitost $O(kN^3)$.

\subsubsection{NSGA-II}

\begin{defbox}[NSGA-II (Deb a~kol., 2000)]
Opravuje nedostatky NSGA. Tři pilíře:
\textbf{(1) rychlé nedominované třídění} do front (složitost $O(kN^2)$),
\textbf{(2) crowding distance} místo parametru $d_{\mathrm{share}}$,
\textbf{(3) elitismus} --- rodiče i~potomci se spojí, setřídí a~lepší polovina
tvoří novou populaci.
\end{defbox}

\begin{defbox}[Crowding distance]
Pro každou frontu zvlášť: pro každé kritérium $m$ seřaď jedince podle $f_m$,
krajním jedincům přiřaď $\infty$ a~ostatním přičti normalizovanou vzdálenost
sousedů:
\[
d_i \;=\; \sum_{m=1}^{k}
\frac{f_m(i+1) - f_m(i-1)}{f_m^{\max} - f_m^{\min}} .
\]
Vyjadřuje \uv{obvod kvádru} tvořeného nejbližšími sousedy --- velká hodnota
znamená řídce obsazenou oblast, kterou chceme zachovat.
\end{defbox}

\begin{defbox}[Crowded comparison operator]
Jedinec $i$ je lepší než $j$, právě když
(a) $i$ leží v~\emph{nižší} frontě ($rank_i < rank_j$), nebo
(b) fronty se shodují a~$i$ má \emph{větší} crowding distance.
Žádná skalární fitness se tedy vůbec nepočítá --- porovnávají se jen tato
dvě čísla, typicky v~binárním turnaji.
\end{defbox}

\begin{algorithm}[!ht]
\caption{NSGA-II (jedna generace)}
\begin{algorithmic}[1]
\State $R_t \gets P_t \cup Q_t$ \Comment{rodiče a~potomci, celkem $2N$ jedinců}
\State $(F_1, F_2, \dots) \gets$ nedominované třídění $R_t$
\State $P_{t+1} \gets \emptyset$; $i \gets 1$
\While{$|P_{t+1}| + |F_i| \le N$}
  \State spočítej crowding distance ve $F_i$; $P_{t+1} \gets P_{t+1} \cup F_i$; $i \gets i+1$
\EndWhile
\State seřaď $F_i$ sestupně podle crowding distance a~doplň $P_{t+1}$ na $N$ jedinců
\State $Q_{t+1} \gets$ turnajová selekce (crowded comparison) $+$ křížení (SBX) $+$ mutace
\end{algorithmic}
\end{algorithm}

\paragraph{Číselný příklad crowding distance.} Fronta
$A(1,10)$, $B(2,7)$, $C(4,5)$, $D(8,2)$; rozsahy: $f_1 \in [1,8]$
(rozpětí 7), $f_2 \in [2,10]$ (rozpětí 8). Krajní body $A$ a~$D$ dostanou
$\infty$ (v~obou kritériích jsou krajní).
\[
d_B = \frac{f_1(C)-f_1(A)}{7} + \frac{f_2(A)-f_2(C)}{8}
    = \frac{4-1}{7} + \frac{10-5}{8} = 0{,}43 + 0{,}63 = 1{,}05,
\]
\[
d_C = \frac{f_1(D)-f_1(B)}{7} + \frac{f_2(B)-f_2(D)}{8}
    = \frac{8-2}{7} + \frac{7-2}{8} = 0{,}86 + 0{,}63 = 1{,}48 .
\]
Při nutnosti vybrat jednoho z~$B$, $C$ zvítězí $C$ --- leží v~řidčeji
obsazené části fronty.

\paragraph{NSGA-III.} Pro úlohy s~mnoha kritérii ($k \ge 4$,
\emph{many-objective}), kde crowding distance selhává (téměř vše je
nedominované). Crowding distance je nahrazena množinou rovnoměrně
rozmístěných \textbf{referenčních bodů} (jejich počet zhruba odpovídá velikosti
populace, $N = \binom{p+k-1}{p}$ pro $p$ dělení). Po nedominovaném třídění se
přednostně doplňují ty referenční směry, které mají nejméně přiřazených
řešení; nemá-li směr žádné řešení, přežije jedinec s~nejmenší kolmou
vzdáleností od příslušného referenčního vektoru.

\subsubsection{Další významné algoritmy}

\noindent\mimoq{slidy končí u~NSGA-II}

\textbf{SPEA2} používá externí \textbf{archiv} nedominovaných řešení; fitness jedince je
součet \emph{sil} všech jedinců, kteří ho dominují (síla = počet jedinců, které daný jedinec
dominuje), plus hustota odvozená ze vzdálenosti k~$\sigma$-tému nejbližšímu sousedovi.
\textbf{MOEA/D} rozloží úlohu na $\mu$ \emph{skalárních} podúloh s~rovnoměrně rozloženými
váhovými vektory (Čebyševova skalarizace); každý jedinec odpovídá jedné podúloze
a~spolupracuje jen se svým sousedstvím --- odtud nízká složitost a~dobrá škálovatelnost.
\textbf{SMS-EMOA} vybírá přímo podle příspěvku k~\emph{hypervolume} (teoreticky nejlépe
podložené, ale exponenciálně drahé v~počtu kritérií). \textbf{NSGA-III} nahrazuje crowding
distance rovnoměrně rozmístěnými \textbf{referenčními body} a~cílí na úlohy s~mnoha kritérii
($k \ge 4$), kde crowding distance selhává.

\subsubsection{Metriky kvality aproximace}

\noindent\mimo{}

\begin{defbox}[Hypervolume ($S$-metrika)]
Objem oblasti dominované nalezenou frontou a~ohraničené \emph{referenčním bodem} $r$
(zvoleným tak, aby byl dominován všemi řešeními). Jediná běžná metrika, která je
\emph{striktně monotónní vůči dominanci} a~nevyžaduje znalost skutečné Pareto fronty.

\emph{Příklad:} fronta $\{(2,7), (4,5), (8,2)\}$, $r = (10,10)$. Po svislých pruzích:
$(8,2)$ dá $2\cdot8 = 16$, $(4,5)$ dá $4\cdot5 = 20$, $(2,7)$ dá $2\cdot3 = 6$; celkem
$HV = 42$. Přidání dominovaného bodu hodnotu nezmění.
\end{defbox}

Dále \emph{GD} (průměrná vzdálenost nalezených bodů od skutečné fronty --- měří konvergenci),
\emph{IGD} (opačně; měří konvergenci i~pokrytí) a~\emph{spread} (rovnoměrnost rozložení);
GD i~IGD vyžadují znalost skutečné fronty.

\subsubsection{Souhrnné srovnání MOEA}

\begin{center}
\small
\begin{tabular}{@{}llll@{}}
\toprule
\textbf{algoritmus} & \textbf{princip selekce} & \textbf{diverzita} & \textbf{poznámka} \\
\midrule
VEGA & podpopulace podle $f_i$ & žádná & konverguje ke krajům \\
NSGA & fronty + fiktivní fitness & fitness sharing ($d_{\mathrm{share}}$) & bez elitismu, $O(kN^3)$ \\
NSGA-II & fronty + crowding & crowding distance & elitismus, $O(kN^2)$, standard \\
NSGA-III & fronty + referenční body & referenční směry & many-objective \\
SPEA2 & síla dominujících & $\sigma$-tý soused, truncation & externí archiv \\
MOEA/D & skalarizované podúlohy & rovnoměrné váhové vektory & sousedství, škáluje \\
SMS-EMOA & příspěvek k~hypervolume & implicitní & teoreticky nejlépe podložené, drahé \\
\bottomrule
\end{tabular}
\end{center}

\subsection{Shrnutí ke~zkoušce}

\begin{itemize}
\item \textbf{Michigan}: jedinec = jedno pravidlo, řešením je celá populace,
  nutné rozdělení odměny (credit assignment).
  \textbf{Pittsburgh}: jedinec = celá sada pravidel, fitness přímočará,
  operátory složité (GIL).
\item Hollandův LCS: podmínky nad $\{0,1,*\}$, vnitřní zprávy a~receptory,
  bucket brigade (pravidlo platí předchůdci za právo jednat).
\item ZCS: bez zpráv; match set $M$ $\to$ ruletový výběr akce $\to$ action set
  $A$; posílení $A$ a~$A_{-1}$, daň pro $M \setminus A$; cover operátor.
\item XCS: klasifikátor $(c,a,p,\epsilon,F)$, \textbf{fitness podle přesnosti
  predikce} $\kappa = \alpha(\epsilon/\epsilon_0)^{-\nu}$, aktualizace
  Q-learningem, nikový GA v~action setu; evolvuje úplnou a~minimální mapu
  vstup--výstup. Varianty: UCS (učení s~učitelem), XCSF (aproximace funkcí).
\item Umět vysvětlit, proč přesnost jako fitness řeší nedostatky
  \uv{strength-based} systémů.
\item Co se evolvuje: váhy / topologie / obojí / hyperparametry učení.
  Výhody EA: bez gradientu, rekurentní sítě, RL, paralelizace.
\item Problém konkurujících konvencí (permutační problém) --- křížení sítí
  je bez opatření destruktivní.
\item Kódování architektury: přímé (matice), gramatické (Kitano), buněčné
  (Gruau, GP program pro růst sítě); GNARL = EP nad grafy, strukturální
  i~parametrické mutace řízené teplotou.
\item NEAT: seznam hran s~\textbf{inovačními čísly} (křížení jen shodných
  genů), \textbf{speciace} podle podobnosti genomů se sdílenou fitness
  (ochrana nových topologií), \textbf{complexification} od minimální sítě.
\item HyperNEAT: substrát $S:\mathbb{R}^4 \to \mathbb{R}$ generující váhy
  z~souřadnic neuronů, reprezentovaný CPPN (DAG s~různými aktivačními
  funkcemi) evolvovaným pomocí NEAT; zachycuje symetrie a~regularity.
\item NAS: prostor (lineární / DAG buněk / super síť), strategie (EA, RL,
  bayesovská, gradientní), odhad výkonu (sdílení vah --- ENAS);
  DeepNEAT a~CoDeepNEAT (moduly + blueprinty); ES pro deep RL.
\item Batoh: bitová mapa, triviální operátory, problém s~omezením kapacity.
  Tři strategie pro omezení: \textbf{penalizace} (volba váhy $\alpha$),
  \textbf{oprava} (lamarckovsky/baldwinovsky), \textbf{dekodér}
  (\uv{přidej, pokud se vejde} --- každý genotyp je přípustný);
  čtvrtá možnost = vícekriteriální formulace.
\item TSP: fitness triviální, problém je kódování. Sousednostní a~ordinální
  reprezentace nepoužívat; permutační + PMX/OX/CX/ER, případně maticová.
  Inicializace: nejbližší soused, vkládání hran. Mutace: inverze (2-opt).
\item SAT: fitness = počet splněných klauzulí (samotné $f$ nestačí),
  vážení klauzulí, zjemňující funkce; reprezentace bitová, reálná,
  klauzulová, cestová; WalkSAT jako lokální heuristika.
  U~reálné reprezentace ($x \mapsto (1-y)^2$, $\neg x \mapsto (1+y)^2$,
  minimalizace) pozor: hodnoty jsou \emph{penalty}, takže disjunkce
  (klauzule) je \textbf{součin} a~konjunkce (formule) \textbf{součet}.
\item Rozvrhování: přímé maticové kódování, tvrdá omezení řešit v~operátorech,
  měkká ve fitness; job shop --- permutace s~dekodérem vs.\ přímé kódování.
\item Obecné pravidlo: u~kombinatorických úloh rozhoduje kódování a~operátory
  respektující strukturu úlohy; hybridizace s~lokální heuristikou (2-opt,
  WSAT) je téměř vždy výhodná.
\item Dominance ($x \prec y$: ve všech kritériích ne horší a~alespoň v~jednom
  lepší), Pareto množina a~Pareto fronta; dvojí cíl = konvergence + diverzita.
\item Skalarizace: lineární (jednoduchá, ale nenajde nekonvexní části fronty),
  $\varepsilon$-constraint, Čebyševova.
\item NSGA-II: nedominované třídění do front, crowding distance
  $d_i = \sum_m (f_m(i+1)-f_m(i-1))/(f_m^{\max}-f_m^{\min})$, krajní body
  $\infty$; crowded comparison (nižší fronta, při shodě větší crowding);
  elitismus spojením $P_t \cup Q_t$. Umět spočítat crowding distance na
  příkladu.
\item SPEA2 (archiv, síla dominujících, hustota přes $\sigma$-tého souseda),
  MOEA/D (Čebyševova skalarizace, váhové vektory, sousedství),
  SMS-EMOA (příspěvek k~hypervolume), NSGA-III (referenční body pro
  many-objective).
\item Hypervolume: objem dominovaný frontou vůči referenčnímu bodu; jediná
  běžná metrika monotónní vůči dominanci a~bez znalosti skutečné fronty.
  Dále GD, IGD, spread, $\varepsilon$-indikátor.
\end{itemize}

%=====================================================================
\section*{Souhrnné srovnání a~typické zkouškové otázky}
\addcontentsline{toc}{section}{Souhrnné srovnání a~typické zkouškové otázky}
\label{sec:zaver}
%=====================================================================
\subsection*{Mapa oboru}

\begin{itemize}
\item \textbf{Evoluční algoritmy} (populace, fitness, křížení, mutace, selekce):
  genetické algoritmy (binární, celočíselné, permutační, reálné), evoluční
  strategie, diferenciální evoluce, evoluční programování, genetické
  programování (stromové, lineární, kartézské, gramatická evoluce),
  learning classifier systems.
\item \textbf{Přírodou inspirované metody bez evoluce}: rojové algoritmy
  (PSO, ACO, ABC, firefly).
\item \textbf{Nebiologicky inspirované metaheuristiky}: tabu search,
  scatter search, novelty search, simulované žíhání, memetické algoritmy.
\item \textbf{Aplikační oblasti} s~vlastními reprezentacemi a~operátory:
  kombinatorická optimalizace, neuroevoluce, zpětnovazební učení,
  vícekriteriální optimalizace, AutoML/NAS.
\end{itemize}

\subsection*{Souhrnná srovnávací tabulka algoritmů}

\begin{center}
\small
\begin{tabular}{@{}p{2.5cm}p{2.9cm}p{3.1cm}p{3.1cm}p{3.1cm}@{}}
\toprule
\textbf{algoritmus} & \textbf{reprezentace} & \textbf{variace} &
\textbf{selekce} & \textbf{co si zapamatovat} \\
\midrule
SGA & binární řetězec & 1-bod.\ křížení, bit-flip & ruleta, generační &
  teorie schémat, BBH \\
ES & $\mathbb{R}^n$ + $\sigma$ & gaussovská mutace, rekombinace &
  náhodní rodiče, $(\mu,\lambda)$/$(\mu{+}\lambda)$ & pravidlo 1/5,
  samoadaptace \\
CMA-ES & $\mathbb{R}^n$, model $\mathcal N(m,\sigma^2 C)$ & vzorkování
  z~rozdělení & pořadová, $\mu$ nejlepších & evoluční cesty, rank-1/rank-$\mu$ \\
DE & $\mathbb{R}^n$ & $\vec x_a + F(\vec x_b - \vec x_c)$, binom.\ křížení &
  souboj rodič--potomek & měřítko z~populace \\
EP & doménová (automat) & jen mutace & turnaj rodiče+potomci &
  genotyp = fenotyp \\
GP & syntaktický strom & výměna podstromů, mutace & turnaj &
  bloat, ADF \\
PSO & $\mathbb{R}^n$ + rychlost & pohybové rovnice & žádná (jen paměť) &
  $w$, $c_1$, $c_2$, gbest/lbest \\
ACO & konstrukce cesty & feromon + heuristika & implicitní &
  $\tau^\alpha\eta^\beta$, odpařování \\
SA & jedno řešení & soused & Metropolis & rozvrh chlazení \\
MA & libovolná & EA operátory + lokální hledání & libovolná &
  Lamarck vs.\ Baldwin \\
NSGA-II & libovolná & SBX + polynom.\ mutace & fronty + crowding &
  elitismus, crowding distance \\
NEAT & seznam hran & mutace struktury, křížení dle ID & sdílená fitness
  v~druhu & inovační čísla, speciace \\
\bottomrule
\end{tabular}
\end{center}

\subsection*{Průřezová témata --- co se ptá napříč otázkami}

\begin{itemize}
\item \textbf{Explorace vs.\ exploatace} --- jak ji řeší každý algoritmus:
  GA (křížení vs.\ selekční tlak), ES ($\sigma$ a~pravidlo 1/5), SA (teplota),
  PSO ($w$ a~topologie), ACO ($q_0$, odpařování), tabu (tabu list),
  novelty search (jen explorace).
\item \textbf{Udržení diverzity} --- fitness sharing, crowding, speciace,
  ostrovní model, nenáhodné páření, restart, hypermutace.
\item \textbf{Adaptace parametrů} --- ladění vs.\ řízení; pevná, adaptivní
  a~samoadaptivní strategie; kde se v~kterém algoritmu vyskytuje.
\item \textbf{Práce s~omezeními} --- penalizace, oprava, dekodér, uzavřené
  operátory, vícekriteriální formulace.
\item \textbf{Lamarckismus vs.\ baldwinismus} --- memetické algoritmy,
  oprava jedinců, neuroevoluce s~douč\-ováním vah.
\item \textbf{Teoretické zázemí} --- věta o~schématech a~její kritika, Voseho
  model, Markovovy řetězce a~konvergence elitistického GA, No Free Lunch,
  konvergence SA.
\end{itemize}

\subsection*{Typické zkouškové otázky}

\begin{enumerate}
\item Popište obecné schéma evolučního algoritmu a~vyjmenujte typy selekce
  včetně jejich vlastností. Spočtěte ruletovou selekci pro danou populaci.
\item Formulujte a~dokažte větu o~schématech. Co je řád a~definující délka?
  Jaké jsou slabiny věty?
\item Vysvětlete hypotézu stavebních bloků a~implicitní paralelismus.
  Co je deceptivní úloha?
\item Popište Voseho model SGA: co jsou vektory $p$ a~$s$, co dělá operátor
  $\mathcal{F}$, co $\mathcal{M}$, jaké mají pevné body? Jak vypadá model
  konečné populace jako Markovův řetězec?
\item Jaký je rozdíl mezi $(\mu,\lambda)$ a~$(\mu+\lambda)$-ES?
  Vysvětlete pravidlo 1/5 a~samoadaptaci $\sigma$. Jak funguje CMA-ES?
\item Popište jeden krok diferenciální evoluce včetně vzorců
  a~číselného příkladu. Jaké jsou varianty mutace?
\item Jak funguje stromové genetické programování? Co je bloat a~jak se proti
  němu bojuje? K~čemu slouží ADF?
\item Vysvětlete gramatickou evoluci a~kartézské GP --- jaký je genotyp
  a~fenotyp?
\item Co je koevoluce, jaké má typy a~jaké patologie? Jak se hodnotí fitness?
\item Co je otevřená evoluce, jaké má znaky a~proč se umělé systémy
  \uv{zasekávají}? Jak funguje novelty search a~proč se u~deceptivních úloh
  vyplatí zahodit cíl?
\item Vysvětlete iterované vězňovo dilema, strategii TFT a~vlastnosti
  úspěšných strategií.
\item Napište rovnice PSO a~vysvětlete význam jednotlivých členů.
  Jaký je rozdíl mezi gbest a~lbest topologií?
\item Popište ACO: pravidlo přechodu, aktualizaci feromonu, rozdíl mezi
  AS a~ACS.
\item Vysvětlete simulované žíhání, Metropolisovo kritérium a~rozvrh chlazení.
  Kdy konverguje?
\item Co je memetický algoritmus? Vysvětlete rozdíl mezi lamarckovským
  a~baldwinovským učením.
\item Porovnejte michiganský a~pittsburghský přístup. Jak funguje XCS a~proč
  je fitness založena na přesnosti?
\item Popište NEAT: k~čemu slouží inovační čísla a~speciace?
  Co přidává HyperNEAT?
\item Jak se řeší TSP evolučním algoritmem? Popište PMX, OX, CX a~ER.
\item Jak se v~EA pracuje s~omezeními (na příkladu problému batohu)?
\item Definujte dominanci a~Pareto frontu. Popište NSGA-II a~spočtěte
  crowding distance na příkladu. Co je hypervolume?
\end{enumerate}

\subsection*{Závěrečné shrnutí}

\begin{itemize}
\item Evoluční a~rojové algoritmy jsou \textbf{robustní metaheuristiky}
  pro úlohy, kde selhávají exaktní a~gradientní metody: black-box,
  nedi\-fe\-ren\-co\-va\-tel\-né, diskrétní, struk\-tu\-ro\-va\-né,
  ví\-ce\-kri\-te\-ri\-ál\-ní a~dynamické úlohy.
\item Jejich síla není v~jednom \uv{nejlepším} algoritmu (No Free Lunch),
  ale ve \textbf{schopnosti přizpůsobit reprezentaci a~operátory dané doméně}
  a~v~možnosti \textbf{hybridizace} s~lokálními heuristikami a~doménovými
  znalostmi.
\item Teorie (schémata, Voseho model, konvergenční věty, analýza doby běhu)
  poskytuje intuici a~a\-sym\-pto\-tic\-ké záruky, ale praktický návrh
  zůstává z~velké části ex\-pe\-ri\-men\-tál\-ní dis\-cip\-lí\-nou.
\end{itemize}

\end{document}
